\chapter{Spatial Reasoning}
\label{sec:reasoning}

\FramedInitial{S}{patial} reasoning is the step that turns perceived structure into decisions: given detections, depths, poses, and a task, it works out the relations, the causes, and the actions that follow, the why and how of spatial interaction rather than merely the what and where. It is where perception becomes goal-directed, supporting inference, planning, and problem-solving in physical and virtual environments.

The chapter is organized along the spectrum by which spatial reasoning has actually developed, from explicit symbolic models to learned and foundation-model-driven ones. At one end, reasoning runs over hand-built representations of geometry, physics, and logic, whose structure is inspectable and whose conclusions can be certified. In the middle, it emerges from data, as neural models acquire spatial competence directly rather than being told the rules. More recently, large pre-trained models supply reasoning as a general capability to be elicited and steered. Hybrid approaches recombine these, pairing the semantic reach of foundation models with the precision and verifiability of symbolic and learned components; a closing survey of benchmarks then records how the field measures spatial reasoning as it shifts from passive observation toward active, situated interaction.

\section{Rule-based Symbolic Spatial Reasoning}
\label{sec:p3_s7_1}

\FramedInitial{R}{ule-based} symbolic spatial reasoning covers the methods that operate over \emph{explicit, model-based representations} of space (geometric primitives, physical laws, graph connectivity, formal constraints, action rules, or causal mechanisms) rather than over learned latent representations alone. What sets this family apart is that reasoning happens inside a model that can be inspected: its state and transitions are legible, so claims about feasibility, optimality, or safety can be backed by certificates rather than empirical averages, and when a plan fails it fails legibly, as an infeasible inverse kinematics, a violated constraint, a missing precondition, or an inconsistent causal attribution, not as an opaque correlation.

The section surveys six kinds of explicit structure, ordered from the most concrete to the most abstract: geometric and algebraic models of position, orientation, motion, and feasibility; physical models of contact, friction, and stability; topological and graph abstractions of space that turn planning into search; formal constraint systems whose solvers return certified plans or proofs of infeasibility; rule-based and programmatic logic that sequences and supervises actions; and causal models that explain, predict counterfactually, and prescribe interventions.
% Table~\ref{tab:p3_s7_1_paradigm_comparison} at the end of the section compares the six.

These paradigms are not isolated alternatives. Modern robotic systems routinely compose them: geometric solvers serve as feasibility oracles inside constraint-based search; behavior trees wrap symbolic plans for fault tolerance; causal action languages compile to SAT; topological maps anchor both classical and learning-based navigation policies. The unifying thread is that each paradigm carries explicit structure (a state, an operator, a constraint, a rule, or a mechanism) that can be inspected, composed, and reasoned about.

\subsection{Geometric and Algebraic Reasoning}
\label{subsec:p3_s7_1_1}
% - Coordinate transformation, workspace computation, FoV estimation
% - Euclidean, projective geometry, Jacobian-based reasoning
% - Applications: reachability analysis, sensor coverage, grasp range

Geometric and algebraic reasoning uses explicit mathematical models of position, orientation, shape, and motion to turn local measurements (image points, depths, joint angles, odometry) into quantities that can be analyzed and verified: where things are in a common frame, what a sensor can see, what a mechanism can reach, and whether a grasp will hold. The same models extend to time-varying settings, where relationships evolve under kinematic and differential constraints and the question becomes not only where entities can be but how they can move. Precision, interpretability, and domain independence have made this the common substrate of robotics~\cite{Murray1994RoboticManipulation,Lynch2017ModernRobotics,siciliano2010robotics}, AR/VR~\cite{hartley2003multiple,azuma1997survey}, and autonomous systems~\cite{Thrun2005Probabilistic,sola2018micro,LaValle2006Planning}, and the source of the hard constraints that learned components are checked against. Figure~\ref{fig:p3_s7_1_1_geo_algebra_abstract} gives the abstract formulation.

\begin{figure}[!ht]
    \centering
    \includegraphics[width=0.99\linewidth]{figures/part3/section7/711_geo.pdf}
    \caption{Abstract task formulation of geometric and algebraic reasoning: (i) local spatial measurements and system models are mapped by (ii) a core of reusable symbolic formalisms to (iii) analyzable, verifiable spatial outputs.}
    \label{fig:p3_s7_1_1_geo_algebra_abstract}
\end{figure}

\subsubsection{Task Definition}
Given measurements $\{\mathbf{z}_k\}$ expressed in local frames (pixels, ranges, joint readings, pose increments) and a system model $\Theta$ (sensor calibration, kinematic parameters, environment geometry), geometric reasoning computes a world-frame quantity and a feasibility verdict,
\begin{equation}
    \mathbf{x} = g(\{\mathbf{z}_k\};\, \Theta), \qquad
    \psi(\mathbf{x};\, \Theta) \in \{\mathrm{true}, \mathrm{false}\},
\end{equation}
where $\mathbf{x}$ may be a pose, a reconstructed surface, a workspace or visibility region, or a trajectory, and $\psi$ is a chosen predicate, such as reachability, visibility, or force closure, checked against the model. In the time-varying case the inputs are indexed over time and $\mathbf{x}$ becomes a trajectory with reachable sets and feasibility certificates along it.

\subsubsection{Methodology}

\paragraph{Transformation and Pose Representation.}
Rigid-body relationships between sensors, objects, and mechanisms are elements of the special Euclidean group
\begin{equation}
\mathbf{T} =
\begin{bmatrix}
\mathbf{R} & \mathbf{t} \\
\mathbf{0}^\top & 1
\end{bmatrix} \in SE(3),
\end{equation}
where $\mathbf{R} \in SO(3)$ is a rotation and $\mathbf{t} \in \mathbb{R}^3$ a translation; composing them moves measurements across frames. Early robotics parameterized serial chains through the Denavit--Hartenberg convention~\cite{Denavit1955Kinematic,Paul1981Robot}, which gives forward kinematics from a minimal set of link parameters. Lie-group formulations instead write motion as exponentials of screw axes, the product-of-exponentials representation~\cite{Murray1994RoboticManipulation,Lynch2017ModernRobotics}, which is coordinate-free and behaves well where the DH parameterization has discontinuities. Quaternions~\cite{Shoemake1985QuaternionCurves,Horn1987Orientation} and dual quaternions~\cite{Daniilidis1999DualQuat,horn2021online} give compact, numerically stable encodings of rotation and screw motion for calibration and sensor fusion.

\paragraph{Projective and Multi-View Geometry.}
Reasoning from images must recover 3D structure and camera motion from 2D measurements under perspective distortion and scale ambiguity. The constraint between corresponding points in two views is the epipolar condition
\begin{equation}
\mathbf{x}'^\top \mathbf{F} \mathbf{x} = 0,
\end{equation}
with $\mathbf{F}$ the fundamental matrix. The eight-point algorithm~\cite{LonguetHiggins1981Scene} solves for $\mathbf{F}$ linearly, Hartley's normalization~\cite{Hartley1997EightPoint} made it numerically stable, and Nistér's five-point solver~\cite{Nister2004FivePoint} reaches the minimal case for calibrated cameras. Zhang's homography-based calibration~\cite{zhang2000flexible} removed the need for special rigs; the field is consolidated in~\cite{hartley2003multiple}, and structure-from-motion systems such as COLMAP~\cite{schonberger2016structure} scale it to thousands of images, with RANSAC~\cite{Fischler1981RANSAC} enforcing the algebraic constraints against outliers throughout.

\paragraph{Jacobian and Differential Formulations.}
For motion rather than static configuration, the map from configuration to task space,
\begin{equation}
\mathbf{x} = f(\mathbf{q}),
\end{equation}
is differentiated to relate velocities,
\begin{equation}
\dot{\mathbf{x}} = \mathbf{J}(\mathbf{q}) \dot{\mathbf{q}},
\end{equation}
and the Jacobian $\mathbf{J}(\mathbf{q})$ carries the information about feasibility, singularities, and controllability along a motion. Whitney's resolved motion rate control~\cite{Whitney1969ResolvedMotion} used it to command Cartesian velocities without closed-form inverse kinematics; Nakamura's damped least squares~\cite{Nakamura1986IKSingularity} trades exactness for stability near singularities; redundancy resolution~\cite{Liegeois1977Redundancy,Siciliano1991Redundancy} uses null-space motion for secondary objectives; and Khatib's operational-space formulation~\cite{Khatib1987OperationalSpace} separates task and null-space dynamics so that objectives are written directly in task space.

\paragraph{Visibility and Coverage Geometry.}
What a sensor at pose $p$ can observe is the set
\begin{equation}
\mathcal{V}(p) = \{ x \in \mathcal{E} \mid \text{line}(p,x) \cap \mathcal{O}
= \emptyset \},
\end{equation}
with $\mathcal{E}$ the environment and $\mathcal{O}$ the obstacles. The Art Gallery theorem~\cite{Chvatal1975ArtGallery,ORourke1987ArtGallery} bounds the number of viewpoints needed for full coverage; polygon visibility~\cite{Lee1983Visibility} and shortest-path queries~\cite{Guibas1987ShortestPath} made the computations algorithmic; and Tarabanis \etal~\cite{Tarabanis1995SensorPlanning} cast sensor planning as geometric optimization over fields of view modeled as cones, frusta, or polygons.

\paragraph{Workspace and Reachability Analysis.}
The reachable workspace of a mechanism is
\begin{equation}
\mathcal{R} = \{ \mathbf{x} \mid \exists\, \mathbf{q} \in \mathcal{C}_{\mathrm{free}}
\text{ such that } \mathbf{x} = f(\mathbf{q}) \},
\end{equation}
with $\mathcal{C}_{\mathrm{free}}$ the feasible configuration space. Roth~\cite{Roth1975Performance} and Gosselin~\cite{Gosselin1990Singularity} added performance and dexterity measures so that manipulability is characterized along with reachability; Merlet~\cite{Merlet1998Workspace} classified workspace types; Zlatanov \etal~\cite{Zlatanov1994Constraints} analyzed constraint-induced singularities; and sampling-based reachability maps~\cite{lu2026richmap} trade closed-form guarantees for empirical coverage in high-dimensional systems.

\paragraph{Contact and Grasp Geometry.}
Whether a set of contacts can hold an object is decided in wrench space. Ferrari and Canny~\cite{Ferrari1992Grasp} model the forces each contact can apply as a convex set of wrenches $\mathcal{W}$; the grasp is force-closed if
\begin{equation}
\mathbf{0} \in \text{int}\,\text{conv}(\mathcal{W}),
\label{eq:p3_s7_1_1_force_closure}
\end{equation}
so that any external wrench can be resisted. GraspIt!~\cite{Miller2004GraspIt} made these metrics computable for arbitrary hands and objects; task-oriented measures~\cite{Lin2015TaskGrasp} weight the wrench space by what the task requires; and Dex-Net~\cite{Mahler2017DexNet} samples analytic grasps at scale to train data-driven grasp predictors that keep the geometric criterion as their label.

\subsubsection{Benchmarks}
Geometric reasoning is evaluated against geometrically grounded ground truth: camera poses and intrinsics, depth maps and meshes, object models, grasp annotations. The pose, odometry, SLAM, and multi-view reconstruction benchmarks (TUM RGB-D~\cite{Sturm2012TUMRGBD}, EuRoC~\cite{Burri2016EuRoC}, TUM-VI~\cite{Schubert2018TUMVI}, KITTI~\cite{KITTI-GeigerLU12}, ETH3D~\cite{Schoeps2017ETH3D}, Tanks and Temples~\cite{knapitsch2017tanks}, DTU~\cite{Aanaes2016DTU}, Middlebury~\cite{Seitz2006Middlebury,Scharstein2002Stereo}, HPatches~\cite{Balntas2017HPatches}) and the long-term localization suites under seasonal and day-night change (Aachen, RobotCar Seasons, CMU Seasons~\cite{Sattler2018Benchmarking}, InLoc~\cite{Taira2018InLoc}, 4Seasons~\cite{wenzel2021seasons}) are those of Section~\ref{sec:visual_odometry} and Section~\ref{sec:mapping_slam}, where their protocols and headline results are tabulated. Specific to the reasoning problems above are BAL~\cite{agarwal2010bundle}, which standardizes large bundle-adjustment problems for solver evaluation; the Active Vision Dataset~\cite{Johns2016AVD}, ScanNet~\cite{dai2017scannet}, and HM3D~\cite{Ramakrishnan2021HM3D} for next-best-view, coverage, and exploration protocols; Cornell~\cite{Jiang2011Cornell}, Jacquard~\cite{Depierre2018Jacquard}, and GraspNet-1Billion~\cite{Fang2020GraspNet} for grasp quality (rectangle labels, simulated grasp trials, and dense 6-DoF annotations with seen / unseen / novel splits respectively), with the YCB object set~\cite{Calli2015YCB} as the shared physical test set; and MoveIt~\cite{Cohen2012MoveIt}, whose reachability tooling is the de facto harness for inverse kinematics and motion-planning evaluation.

\subsubsection{Challenges and Trends}
The formulations above assume exact geometry, and real sensing is noisy, incomplete, and subject to calibration drift~\cite{Sturm2012TUMRGBD,Burri2016EuRoC}; bundle adjustment~\cite{agarwal2010bundle} and reachability analysis scale poorly with problem size; many methods fail near singularities or degenerate feature configurations~\cite{Nakamura1986IKSingularity,Hartley1997EightPoint}; the frameworks are strongest for static rigid bodies and approximate deformable and human-centric geometry~\cite{SicilianoKhatib2016}; and evaluation is concentrated on a few canonical datasets and manipulators, which limits cross-domain comparison.

The response has been to keep the geometry and wrap it in probabilistic or learned machinery: factor-graph SLAM and robust bundle adjustment attach uncertainty to the symbolic model (Section~\ref{sec:mapping_slam}); learned correspondence priors are refined by RANSAC~\cite{Fischler1981RANSAC}; and the long-term benchmarks~\cite{wenzel2021seasons,Sattler2018Benchmarking} push the same pipelines toward continual adaptation under changing conditions.




\subsection{Physical Reasoning}
\label{subsec:p3_s7_1_2}
% - Concerned with physical laws governing spatial environments (mechanics, dynamics, collisions, friction, stability, etc.).
% - Newtonian mechanics, quasi-static models, collision simulation
% - Applications: grasp stability, push/pull prediction, object stacking

Physically-aware spatial reasoning predicts how a spatial configuration evolves under physical law: rigid-body dynamics, contact and friction, and the conditions for static stability. Will a grasp hold? How will an object move when pushed? Will a stack of blocks stand or fall? Answering these from a model of mechanics, rather than from observed correlations, is what lets a robot anticipate the outcome of an interaction before committing to it; Figure~\ref{fig:physical-reasoning} shows the three canonical cases.

\subsubsection{Task Definition}
Let the environment at time $t$ be the state
\begin{equation}
    \mathbf{s}_t = (\mathbf{q}_t, \mathbf{v}_t, \theta),
\end{equation}
with generalized coordinates $\mathbf{q}_t \in \mathbb{R}^n$, velocities $\mathbf{v}_t=\dot{\mathbf{q}}_t$, and physical parameters $\theta$ (masses and inertias in the mass matrix $M(\mathbf{q})$, friction coefficients $\{\mu_k\}$, restitution, geometry). An action $\mathbf{u}_t$ (a push, a grasp) enters as a generalized force $\boldsymbol{\tau}_{\mathrm{ext},t+1}$, and the next state follows from a \emph{contact-implicit time-stepping} update:
\begin{equation}
\begin{aligned}
M(\mathbf{q}_{t})\,\frac{\mathbf{v}_{t+1}-\mathbf{v}_{t}}{\Delta t}
+ \mathbf{h}(\mathbf{q}_{t},\mathbf{v}_{t+1})
&= \boldsymbol{\tau}_{\mathrm{ext},t+1}
+ \sum_{k=1}^{K}\Big(J_{n,k}^\top \lambda_{n,k} + J_{t,k}^\top \boldsymbol{\lambda}_{t,k}\Big), \\
\mathbf{q}_{t+1}&=\mathbf{q}_{t}+\Delta t\,\mathbf{v}_{t+1},
\end{aligned}
\end{equation}
where $\mathbf{h}$ collects Coriolis, centrifugal, and gravitational terms, and for each potential contact $k$ with gap function $\phi_k(\mathbf{q})$ and Jacobians $J_{n,k},J_{t,k}$ the normal and tangential contact forces satisfy non-penetration, complementarity, and Coulomb friction:
\begin{equation}
\begin{aligned}
\phi_k(\mathbf{q}_{t+1}) \ge 0,\quad \lambda_{n,k}&\ge 0,\quad \lambda_{n,k}\,\phi_k(\mathbf{q}_{t+1})=0,\\
\|\boldsymbol{\lambda}_{t,k}\| &\le \mu_k\,\lambda_{n,k}.
\end{aligned}
\label{eq:p3_s7_1_2_contact}
\end{equation}
Physical reasoning is the computation, or approximation, of this transition $\mathbf{s}_{t+1}=\mathcal{T}(\mathbf{s}_t,\mathbf{u}_t)$ and of the queries derived from it (stability, collision, reachability). When motion is slow enough that inertial terms vanish, the \emph{quasi-static} form
\begin{equation}
\sum_{k=1}^{K}\Big(J_{n,k}^\top \lambda_{n,k} + J_{t,k}^\top \boldsymbol{\lambda}_{t,k}\Big) + \boldsymbol{\tau}_{\mathrm{ext}} - \mathbf{h}(\mathbf{q},\mathbf{0}) \approx \mathbf{0},
\qquad \|\boldsymbol{\lambda}_{t,k}\| \le \mu_k\,\lambda_{n,k},
\end{equation}
reduces reasoning to an equilibrium feasibility test, which is the setting of the classical analytic models of grasping, pushing, and stacking.

\begin{figure}[t]
  \centering
  \includegraphics[width=0.95\linewidth]{figures/part3/section6/physical_reasoning.pdf}
  \caption{Physics-aware spatial reasoning in robotics: (a) grasp stability from contact forces and friction cones, (b) planar pushing under quasi-static assumptions, and (c) stacking stability under equilibrium and support constraints.}
  \label{fig:physical-reasoning}
\end{figure}

\subsubsection{Methodology}
Model-based physical reasoning is carried out either numerically, in a physics engine that integrates the dynamics above, or analytically, with closed-form criteria for the quasi-static cases. Both rest on the same Newtonian foundation and differ in whether fidelity or interpretability is the priority.

\paragraph{Physics Engines.}
MuJoCo~\cite{todorov2012mujoco}, Bullet~\cite{coumans2021}, ODE, and DART~\cite{DART2023} integrate rigid-body dynamics with contact by time stepping, and they resolve contact in one of two ways. Penalty methods treat a contact as a compliant spring-damper that resists interpenetration,
\begin{equation}
f_{n,k} = \max\!\left(0, k_p\,\delta_k + k_d\,\dot{\delta}_k\right),
\end{equation}
with penetration depth $\delta_k$, stiffness $k_p$, and damping $k_d$, and regularize friction similarly; they are cheap and simple, but stiff parameters destabilize the integration and soft ones introduce artificial compliance, which is acceptable in graphics and approximate planning but not in control. Constraint-based methods enforce Eq.~\eqref{eq:p3_s7_1_2_contact} directly as complementarity constraints~\cite{trinkle1997dynamic}, solved per step as a linear or nonlinear complementarity problem~\cite{cottle2009linear}, a convex quadratic program~\cite{lin1987iterative}, or a second-order cone program for exact friction cones~\cite{alizadeh2003second} (projected Gauss--Seidel and sequential-impulse schemes are the usual implementations). They give physically consistent, nearly penetration-free contact forces at higher cost and with harder scaling in the number of contacts, so engines in practice hybridize: hard constraints on normal forces, regularized tangential friction, warm starting, contact aggregation, and stabilization terms.

\paragraph{Analytic Quasi-Static Models.}
For the three settings of Figure~\ref{fig:physical-reasoning}, closed-form criteria replace simulation. At a point contact the admissible forces form the friction cone
\begin{equation}
\mathcal{F}_k = \{(f_{n,k},\mathbf{f}_{t,k}) \mid f_{n,k} \ge 0, \ \|\mathbf{f}_{t,k}\|\le \mu_k f_{n,k}\},
\end{equation}
which a polyhedral approximation turns into linear constraints. Mapping each contact's cone through its moment arm into wrench space and taking the convex hull gives the grasp wrench space, and the grasp is force-closed when the origin lies strictly inside it, Eq.~\eqref{eq:p3_s7_1_1_force_closure}~\cite{bicchi2000robotic}; form closure is the stronger, purely geometric condition in which contact geometry alone prevents motion, and the Ferrari--Canny $\epsilon$-metric (the distance from the origin to the hull boundary) grades grasp quality. For planar pushing the limit surface relates the applied wrench to the resulting slip, and its ellipsoidal approximation yields a motion cone of object velocities reachable from a given push, the basis of stable pushing and closed-loop push planning. For stacking, stability under gravity requires that the gravitational wrench be balanced by the supporting contacts~\cite{groth2018shapestacks}: without friction, the center of mass must project into the support polygon; with friction, the contacts' friction cones must generate a wrench cone containing the negative of the gravitational wrench.

\subsubsection{Benchmarks}
PHYRE~\cite{bakhtin2019phyre} poses 2D physics puzzles in which an object must be placed so that simulated physics reaches a goal, scored by the log-weighted area under the success-rate curve as the attempt budget increases (AUCCESS) within and across puzzle templates; ShapeStacks~\cite{groth2018shapestacks} labels 3D stacks as stable or unstable; PlasticineLab~\cite{huang2021plasticinelab} is a differentiable-physics suite for deformable manipulation; and the video suites IntPhys~\cite{riochet2018intphys}, CLEVRER~\cite{Yi2020CLEVRER}, and Physion~\cite{Bear2021Physion} test prediction and counterfactual judgment of physical outcomes, with ComPhy~\cite{chen2022comphy} asking for hidden properties such as mass to be inferred from motion. Grasp quality under quasi-static analysis is evaluated on the grasping benchmarks of §\ref{subsec:p3_s7_1_1} (Cornell, YCB, GraspNet-1Billion). Metrics are predictive: position or trajectory error against the simulated outcome, success rate (stable stack, successful push or grasp), binary stability accuracy, and physical-plausibility terms such as momentum conservation and constraint violation.

\subsubsection{Challenges and Trends}
The parameters that matter most, mass, friction, elasticity, internal stability, are rarely observable. They must be inferred from interaction, from how far an object moves when pushed or how it responds to a collision; ComPhy~\cite{chen2022comphy} tests exactly this, and most methods still either assume the parameters known, as classical simulation does, or approximate them in ways that fail outside the lab. Latent-dynamics learning~\cite{battaglia2016interaction} and system identification~\cite{an1985estimation,muralidhar2020phynet} are the two established routes.

The simulation-to-reality gap is the second limit. Small errors in restitution, micro-contact behavior, or unmodeled environmental factors compound once a policy leaves the engine, and fidelity trades against speed: high-fidelity simulation is too slow for large-scale training or real-time control, lightweight approximation too inaccurate. Hybrid physics-plus-learning models, domain adaptation between simulated and real dynamics, and uncertainty-aware transfer each address one side of this trade-off; none closes it.

Perception is the third. Geometry, material, and object state come from noisy, incomplete sensing, and a small error in shape, pose, or contact location propagates into a large error in predicted force or stability. Systems that couple vision to a physics engine, Galileo~\cite{wu2015galileo} and PhyNet~\cite{muralidhar2020phynet} among them, show that physical priors improve perception and that physical consistency can constrain it in return, but end-to-end pipelines in which the two regularize each other remain the exception.



\subsection{Topological and Graph-based Reasoning}
\label{subsec:p3_s7_1_3}
% - Roadmap construction, A*, Dijkstra, Voronoi diagrams
% - Applications: navigation planning, environment mapping, multi-agent pathfinding

Topological and graph-based reasoning replaces continuous geometry with discrete connectivity, so that path planning, exploration, and coordination become combinatorial search over nodes and edges: places, corridors, and routes rather than metric trajectories. The abstraction is what makes large-scale, dynamic, or partially known environments tractable, where maintaining globally consistent geometry is hard but connectivity and reachability are stable. Figure~\ref{fig:p3_s7_1_3_topo_graph_abstract} gives the abstract formulation.

\subsubsection{Task Definition}
Given a configuration space $\mathcal{C}$ with free subset $\mathcal{C}_{\mathrm{free}} \subseteq \mathcal{C}$ (known a priori or inferred online from sensor streams), robot constraints, and optional priors such as floor plans or road networks, most methods build a planning graph $G=(V,E,\mathbf{w})$ in which vertices $V$ are places, sampled configurations, or abstract states, edges $E$ are feasible transitions, and $\mathbf{w}: E \to \mathbb{R}_{\geq 0}$ assigns metric or semantic costs (distance, time, energy, risk); in some methods the graph is implicit and expanded locally during search. The output is a collision-free path $\pi=\langle v_0=s,\dots,v_k=g\rangle$ of minimal accumulated cost, a topological map of place connectivity (often linked to local metric submaps), or, with several agents, conflict-free trajectories $\{\pi_i(t)\}_{i=1}^N$ that minimize makespan or total cost under vertex and edge conflict constraints.


\subsubsection{Methodology}

\paragraph{Visibility and Line-of-Sight Graphs.}
Connecting configurations that can see each other reduces shortest-path planning to graph search with interpretable waypoints. The visibility graph of Lozano-P{\'e}rez and Wesley~\cite{LozanoPerez1979Visibility} links obstacle, start, and goal vertices and recovers exact Euclidean shortest paths in polygonal environments by Dijkstra or A*; the configuration-space construction via Minkowski sums~\cite{LozanoPerez1983Cspace} extends it to robots of finite size by growing the obstacles instead of shrinking the robot; continuous-Dijkstra algorithms settle the complexity at $O(n\log n)$~\cite{Hershberger1999EuclideanShortest,Mitchell2000Handbook}. Grid planners bring line-of-sight into the search itself: Field D* interpolates across cell faces for smooth, replannable paths~\cite{Ferguson2006FieldD}, Theta* updates parents by visibility to produce near-Euclidean paths on the fly~\cite{Nash2007Theta}, and Jump Point Search prunes symmetric grid paths by jumping to points where successor choices change~\cite{Harabor2011JPS}.


\begin{figure}[!ht]
    \centering
    \includegraphics[width=0.99\linewidth]{figures/part3/section7/713_topo.pdf}
    \caption{Abstract task formulation of topological and graph-based reasoning: (i) free space, sensing, start and goal states, and constraints are turned by (ii) construction and search over a graph $G=(V,E,w)$ into (iii) decision-ready graph artifacts such as collision-free paths, topological maps, coordinated multi-agent trajectories, and coverage plans. These are alternative task outputs. In the multi-agent example, rows denote vertices and columns denote successive time steps; horizontal segments represent waiting.}
    \label{fig:p3_s7_1_3_topo_graph_abstract}
\end{figure}


\paragraph{Clearance-Oriented Skeletons.}
Where safety matters more than length, free space is retracted to a skeleton that keeps connectivity and maximizes distance to obstacles. Voronoi diagrams and medial axes give maximal-clearance paths, with Fortune's sweepline~\cite{Fortune1987Voronoi} for efficient planar construction; the generalized and hierarchical Voronoi graphs~\cite{Choset2000HGVG,Choset2000GVGExploration} let a robot follow ridges with its sensors and build the clearance roadmap incrementally while detecting junctions and loops; and distance-transform updates~\cite{Felzenszwalb2012EDT,Lau2013DynamicEDT} maintain the skeleton in dynamic environments for real-time replanning.

\paragraph{Decomposition Graphs and Coverage.}
Discretizing space into cells trades geometric fidelity for simplicity and completeness. Occupancy grids~\cite{Moravec1985SonarMaps,Elfes1989Occupancy} are probabilistic cell maps updated from range data on which planning is grid search; exact decompositions serve coverage, with boustrophedon decomposition~\cite{Choset1998Boustro} partitioning free space into monotone cells and ordering them through an adjacency graph for complete, low-redundancy sweeps; and the coverage literature is systematized in~\cite{Galceran2013CoverageSurvey}.

\paragraph{Sampling-Based Roadmaps and Trees.}
In high-dimensional configuration spaces, the graph is discovered by sampling. Probabilistic roadmaps~\cite{Kavraki1996PRM} sample and connect configurations in a preprocessing phase amortized over many queries and are probabilistically complete; RRT and RRT-Connect~\cite{LaValle1998RRT,Kuffner2000RRTConnect} grow trees toward unexplored regions for single-query planning under kinematic and kinodynamic constraints; PRM* and RRT*~\cite{Karaman2011RRTStar} add connection-radius and rewiring rules that make them asymptotically optimal; and FMT*~\cite{Janson2015FMT} performs lazy dynamic programming over a batch of samples, while BIT*~\cite{Gammell2015BIT} combines heuristic search with progressively denser sample batches for anytime improvement.

\paragraph{Optimal, Any-Angle, and Incremental Graph Search.}
On a given graph, Dijkstra's algorithm and A*~\cite{Dijkstra1959,Hart1968Astar} are the optimal single-source and heuristic searches; LPA*~\cite{Koenig2004LPA} and D* Lite~\cite{Koenig2002DstarLite} repair a shortest path incrementally as edge costs change; and the any-angle planners above (Field D*, Theta*) improve path quality, while JPS accelerates optimal search on its underlying grid.

\paragraph{Multi-Agent and Hierarchical Graph Reasoning.}
With several agents on one graph, coordination is search over a time-expanded network. Conflict-Based Search~\cite{Sharon2015CBS} splits the problem into a high-level search over conflict constraints and low-level single-agent planning and is complete and optimal; ICTS~\cite{Sharon2013ICTS} searches joint cost allocations backed by individual searches; ECBS and EECBS~\cite{Barer2014ECBS,Li2021EECBS} are bounded-suboptimal with focal heuristics and better conflict handling; and SAT or ILP compilations~\cite{Surynek2016SATMAPF,Yu2016OptimalMAPF} reach global optimality on moderate fleets at the cost of solver scalability. Hierarchical and multi-resolution planning~\cite{Botea2004HPA} and contraction hierarchies~\cite{Geisberger2008CH} precompute abstraction layers that answer shortest-path queries on continental road networks in well under a millisecond, with standardized maps and instances~\cite{Sturtevant2012Benchmarks,MovingAI2025} for comparison.

\subsubsection{Benchmarks}
The MovingAI suite~\cite{Sturtevant2012Benchmarks,MovingAI2025} supplies standardized 2D maps from games, cities, and synthetic generators for single-agent search (path cost, runtime, memory, node expansions) and the MAPF instances on which CBS, ICTS, ECBS, and compilation-based solvers are compared by makespan, sum of costs, and success under time and memory limits, with the ICAPS MAPF tracker~\cite{Stern2019MAPFsurvey} as the community leaderboard. OMPL's Planner Arena~\cite{OMPL2025} compares PRM, RRT, FMT, and BIT variants on success rate, planning time, path quality, and anytime behavior, and the DIMACS shortest-path challenge~\cite{DIMACS2025} evaluates query latency, preprocessing cost, and optimality on continental road graphs. Topological mapping and pose-graph SLAM are stressed by place recognition under appearance change: RADISH~\cite{radish2003} provides classic multi-sensor logs, and Oxford RobotCar~\cite{Maddern2017RobotCar}, NCLT~\cite{Carlevaris2016NCLT}, and Nordland~\cite{olid2018singleviewplace} span seasons and weather, with KITTI and EuRoC (Section~\ref{sec:visual_odometry}) for urban and indoor pose graphs.

\subsubsection{Challenges and Trends}
Graph density sets the trade-off: sparse graphs lose connectivity or optimality, dense ones lose real-time response, and the problem sharpens in high-dimensional and lifelong settings~\cite{Karaman2011RRTStar,Koenig2002DstarLite}. Perceptual aliasing, drifting pose graphs, and dynamic obstacles complicate planning under partial information~\cite{Thrun2005Probabilistic,Stachniss2006Mapping}, and on the topological SLAM benchmarks the binding difficulty is not search but appearance robustness, with cross-season and cross-weather queries far harder than same-condition retrieval. Multi-agent coordination strains conflict resolution and time-expanded representations as fleets grow, while fairness and deadlock avoidance must be kept~\cite{Sharon2015CBS,Yu2016OptimalMAPF,Stern2019MAPFsurvey}.

The directions are structured rather than end-to-end: hierarchical and multi-resolution planning that abstracts and refines~\cite{Botea2004HPA,Geisberger2008CH}, risk-aware and formally verified planning under uncertainty~\cite{Luders2010ChanceConstraints,zhang2023safety}, bounded-suboptimal and distributed coordination for responsiveness~\cite{Barer2014ECBS,Li2021EECBS}, and neuro-symbolic methods that enrich the graph with learned perception, heuristics, and semantics~\cite{kimura-etal-2021-neuro}.





\subsection{Constraint-based Spatial Solving}
\label{subsec:p3_s7_1_4}
% - Continuous geometric/motion feasibility: identifying configurations or trajectories that satisfy geometric, dynamic, collision, and temporal constraints under a given goal.
% - CSPs, SAT/SMT, ILP/MILP, QP/NLP, IK solvers, feasible region modeling
% - Applications: manipulation planning, multi-robot coordination, assembly validation

Many robotics problems are not about finding \emph{a} plan but one that \emph{must} satisfy strict requirements: no collisions, no joint-limit violations, stable contacts, correct timing, sometimes guaranteed dynamics or safety margins. Constraint-based spatial solving states the problem as a collection of constraints, what must hold and what must not, and hands it to a solver that either returns a satisfying solution or proves that none exists. Compared with sampling- and heuristic-driven planning, this buys certificates of feasibility, an explanation of why a candidate fails, and, under an explicit model, optimality or bounded-suboptimality guarantees, which is what manipulation in clutter, multi-robot coordination in shared space, and assembly demand. Figure~\ref{fig:p3_s7_1_4_constraint_solving_abstract} gives the abstract formulation.

\begin{figure}[!ht]
    \centering
    \includegraphics[width=0.99\linewidth]{figures/part3/section7/714_constraint.pdf}
    \caption{Abstract task formulation of constraint-based spatial solving: (i) decision variables and constraints, (ii) a solver core that satisfies or optimizes them, and (iii) feasible trajectories, schedules, and assembly sequences, with feasibility or infeasibility certificates and optimality-gap bounds where supported by the solver and problem formulation. The output types are alternatives; failure of a local solver alone is not an infeasibility proof.}
    \label{fig:p3_s7_1_4_constraint_solving_abstract}
\end{figure}

\subsubsection{Task Definition}
Let $x=(\mathbf{q}_{1:T},\mathbf{u}_{1:T},\mathbf{z})$ collect continuous states and controls together with discrete mode and logic variables $\mathbf{z}$. Constraint-based solving seeks
\begin{equation}
\text{find } x\in\mathcal{X}\quad\text{s.t.}\quad g_i(x)\le 0,\ \ h_j(x)=0,\ \ \phi(\mathbf{z},x)=\top,
\end{equation}
where the inequalities $g_i$ encode geometric and dynamic feasibility (collision, joint, and dynamics limits), the equalities $h_j$ encode kinematics and contact conditions, and $\phi$ encodes task logic such as preconditions and effects. The optimization variant minimizes an objective $f(x)$ as well. The output is a feasible $x$ (a collision-free and dynamically admissible trajectory, a task-robot schedule, an assembly sequence), with an objective value or optimality gap when optimizing, or a certificate that no such $x$ exists.

\subsubsection{Methodology}
The solver families differ in the variables and constraints they handle and the guarantees they return; practical systems are rarely pure, combining discrete reasoning over actions and modes with continuous optimization over motions and contacts, and calling fast geometric oracles (collision checking, inverse kinematics) inside the solver loop.

\paragraph{Constraint Programming (CSP/CP).}
When the decisions are combinatorial (which object, which placement, which grasp) but each must pass a geometric test, constraint programming encodes the choices as variables under constraints and interleaves symbolic branching with propagation, pruning, and geometric checks. Lozano-P{\'e}rez and Kaelbling~\cite{LozanoPerez2014ConstraintTAMP} cast task-and-motion planning as constraint satisfaction over action skeletons, pruning symbolic branches as soon as a geometric feasibility test fails; Lagriffoul \etal~\cite{Lagriffoul2012ConstraintPropagation} propagate and filter continuous domains before any expensive collision or IK call, so that many failures become cheap inconsistencies; and Dantam \etal~\cite{Dantam2018IncrementalTAMP} interleave discrete search with incremental feasibility queries that reuse information across related constraints.

\paragraph{SAT/SMT Compilations.}
Compiling planning and feasibility into Boolean formulas, or into formulas over richer theories such as real arithmetic, puts clause learning, incremental solving, and conflict-driven pruning to work. dReal~\cite{Gao2013dReal} decides nonlinear real arithmetic up to a numerical tolerance ($\delta$-completeness), which gives certified answers where exact decidability is impossible; SMT-based planning for hybrid systems~\cite{Shoukry2016SMTPlanning} reasons first over an abstraction and adds constraints as counterexamples appear; SAT encodings of multi-agent pathfinding~\cite{Surynek2016SATMAPF} obtain optimal schedules on time-expanded graphs (the MAPF problem itself is treated in §\ref{subsec:p3_s7_1_3}); and SMT over linear real arithmetic extends the encodings to continuous time and to robustness against execution delays~\cite{Kolarik2023MAPFSMT,Tan2024RobustMAPF}.

\paragraph{Mixed-Integer Formulations (MICP/MILP).}
When discrete mode choices (which region, which contact mode, which precondition) interact with continuous geometry and dynamics, binary variables encode the nonconvex choice and convex constraints capture feasibility within a mode. IRIS~\cite{Deits2014IRIS} inflates large convex obstacle-free regions that serve as certified safe sets, turning collision avoidance into membership constraints with explicit margins; Graphs of Convex Sets~\cite{Marcucci2024GCSShortest} compile motion planning through clutter into a shortest-path problem over convex nodes and edges that convex optimization, with mixed-integer selection where needed, solves with certificates; and later work constructs convex sets directly in configuration space and in space-time~\cite{Marcucci2023SciRobotics,Dai2024CIRIS}. In each case a globally nonconvex feasibility problem is replaced by a sequence of convex ones without giving up the guarantee.

\paragraph{Trajectory Optimization.}
Even with the task fixed, an executable motion must satisfy kinematic limits, collision avoidance, smoothness, and often dynamics, which is a continuous, usually nonconvex, constrained optimization whose local optimum comes with optimality conditions and sensitivities. CHOMP~\cite{Zucker2013CHOMP} deforms an initial path by functional gradient descent on a signed-distance field; STOMP~\cite{Kalakrishnan2011STOMP} replaces the gradient with stochastic perturbations to handle non-smooth costs and local minima; TrajOpt~\cite{Schulman2014TrajOpt} solves a sequence of convex programs with hinge-loss collision penalties; and GPMP2~\cite{Mukadam2018GPMP2} writes planning as inference on a sparse factor graph with Gaussian-process priors, which gives scalability and incremental replanning. Increasingly these optimizers run inside the safe sets of IRIS and GCS, which supply initialization, constrain the search, or verify the result.

\paragraph{Inverse Kinematics as a Constraint Subproblem.}
``Is there a joint configuration that reaches this pose under limits and collisions?'' is the most frequent feasibility query in any hybrid planner, so IK solvers are built as fast oracles. IKFast~\cite{Diankov2010IKFast} derives closed-form solutions and generates optimized code for repeated real-time queries; TRAC-IK~\cite{Beeson2015TRACIK} combines complementary numerical strategies to raise the solve rate and avoid stalling the planner above it; and global IK by mixed-integer programming~\cite{Dai2019GlobalIK} returns certifiable solutions for full chains when a feasibility proof is needed.

\paragraph{Hybrid Logic--Geometric Planning.}
In long-horizon problems the discrete decision (what to do next) and the continuous one (how to do it) cannot be separated, and task-and-motion planners close the loop between them. Logic-Geometric Programming~\cite{toussaint2015lgp,Toussaint2017LGPcoop} couples symbolic action sequences to bound-constrained nonlinear programs through a multi-bound tree search, so that continuous feasibility and cost decide which symbolic branches are worth expanding; PDDLStream~\cite{garrett2020pddlstream} integrates black-box geometric generators (grasps, placements, motions) into symbolic planning by optimistic refinement, requesting continuous witnesses only when needed while remaining complete under suitable stream assumptions; R-LGP~\cite{Ly2023RLGP} biases the search with reachability guidance; and COAST~\cite{Vu2024COAST} uses constraints to communicate between task planner and stream generators, cutting wasted geometric proposals by an order of magnitude in planning time.

\subsubsection{Benchmarks}
Evaluation asks whether a method returns solutions that satisfy strict feasibility conditions and, increasingly, whether it can show that it did (margins, certificates, optimality gaps) under a comparable oracle budget of collision checks and IK calls. OMPL's benchmark suite and Planner Arena~\cite{Sucan2012OMPL} score planners on standardized problems with feasibility diagnostics and solution cost; the platform-independent TAMP benchmarks of Lagriffoul \etal~\cite{Lagriffoul2018TAMPBench} specify symbolic operators together with IK and collision feasibility so that a jointly feasible task-motion plan is required; and the assembly benchmarks RAMP~\cite{Collins2023RAMP} (real-world assemblies requiring geometric fit, stability under gravity, and grasp or fixture feasibility), ASAP~\cite{Tian2024ASAP} (large-scale 3D assembly sequences validated by physics-based feasibility checks), and FurnitureBench~\cite{Heo2023FurnitureBench} (long-horizon real-world furniture assembly scored by success rate and completion time) stress feasibility across geometry, stability, and sequencing at once. Multi-robot coordination under space-time collision constraints uses the MAPF benchmarks of §\ref{subsec:p3_s7_1_3}.

\subsubsection{Challenges and Trends}
The central difficulty is integrating a discrete task structure with continuous, highly nonconvex feasibility at scale and with certificates. Task-level choices (modes, grasps, placements, contact sets, schedules) interact combinatorially with geometric and dynamic constraints and fragment the feasible set; mixed-integer formulations grow expensive with the number of binary decisions~\cite{Deits2014IRIS,Marcucci2023SciRobotics}, trajectory optimizers return locally feasible motions that may violate a symbolic assumption such as the intended contact mode~\cite{Zucker2013CHOMP,Schulman2014TrajOpt}, and symbolic planners propose sequences that are geometrically unrealizable. Information flows poorly in both directions: continuous solvers rarely return explanations (active sets, separating hyperplanes, unsatisfiable cores) that the symbolic layer could use to prune, and symbolic abstractions are often too coarse to keep the continuous subproblems tractable~\cite{garrett2020pddlstream,toussaint2015lgp}. On the hybrid and assembly benchmarks this coupling, not raw search speed, is what separates planners: those that interleave symbolic branching with continuous oracles waste fewer feasibility checks and degrade more gracefully across domains.

Assurance under uncertainty is the second difficulty. Real systems need machine-checkable feasibility and fast infeasibility proofs, but nonlinear dynamics, contact, and hybrid interactions make strong certificates hard; convex surrogates and bounded-horizon relaxations buy speed at the price of completeness or tightness; execution requires real-time replanning with expensive oracles in the loop; and robustness to noise and model error calls for explicit risk and temporal-logic constraints~\cite{Blackmore2011ChanceConstraints,Raman2014STL,Ames2019CBF}. Evaluations still under-report certificate quality, oracle counts, and failure explanations.

The trend is toward structure with certificates: graphs of convex or geodesically convex safe sets that reduce global nonconvex feasibility to sequences of convex subproblems with explicit proofs~\cite{Marcucci2024GCSShortest,Dai2024CIRIS}, coupled to trajectory-optimization and MPC back-ends that return sensitivities and constraint activations for the search above them; chance constraints, temporal logic, and control barrier functions for persistent safety~\cite{Blackmore2011ChanceConstraints,Ames2019CBF}; GPU acceleration for real-time replanning; and learning that guides branching, cutting, or warm starts without discarding the guarantee~\cite{Ly2023RLGP,Vu2024COAST}.





\subsection{Rule-based and Programmatic Reasoning}
\label{subsec:p3_s7_1_5}
% - Discrete task logic and sequential decision-making: using rules/operators to specify when and in what order actions should be performed.
% - FSM, PDDL, STRIPS, Behavior Trees, If-then rules
% - Applications: multi-step task planning, logic-guided assembly, agent behavior switching

Rule-based and programmatic reasoning is the explicit, human-auditable layer that decides \emph{when} an action is allowed, \emph{in what order} actions happen, and \emph{what to do next} when conditions change, using symbolic operators, transition rules, or program-like control structures. In a robot it is the supervisory layer that sequences reusable skills (navigate, grasp, insert) and invokes geometric solvers or learned policies only when their preconditions hold, which is what makes a decision explainable (why the detour when the corridor was blocked), verifiable, and able to carry commonsense and safety constraints through a long task. Figure~\ref{fig:p3_s7_1_5_rule_programmatic_abstract} gives the abstract formulation.

\begin{figure}[t]
    \centering
    \includegraphics[width=0.99\linewidth]{figures/part3/section7/715_program.pdf}
    \caption{Abstract task formulation of rule-based and programmatic reasoning: (i) symbolic task structure, (ii) four representations for executing decision logic, and (iii) executable plans, reactive policies, and monitored execution traces with replanning after a disturbance.}
    \label{fig:p3_s7_1_5_rule_programmatic_abstract}
\end{figure}

\subsubsection{Task Definition}
A symbolic planning problem is a tuple $\Pi=\langle \mathcal{S},\mathcal{A},s_0,G\rangle$ with symbolic states $\mathcal{S}$, an operator library $\mathcal{A}$, an initial state $s_0$, and a goal $G$; each operator $a\in\mathcal{A}$ carries a precondition and effects,
\begin{equation}
a=\big(\mathrm{pre}(a),\ \mathrm{eff}^{+}(a),\ \mathrm{eff}^{-}(a)\big),\qquad s'=\big(s\setminus\mathrm{eff}^{-}(a)\big)\cup\mathrm{eff}^{+}(a),
\end{equation}
and is applicable in $s$ iff $\mathrm{pre}(a)\subseteq s$. The planning output is a sequence $\pi=\langle a_1,\dots,a_k\rangle$ whose successive application to $s_0$ reaches a state satisfying $G$; the reactive variant instead returns a policy $\mu$ that maps the current observation and symbolic state to a behavior. Depending on the formalism, the inputs also carry resources, temporal constraints, and exogenous events.

\subsubsection{Methodology}
The formalisms split by whether they \emph{plan} (search for a correct sequence before acting) or \emph{react} (choose a behavior continually at run time); deployed systems usually do both, with a planner proposing the sequence and an execution structure (FSM, behavior tree, rule engine) monitoring conditions and triggering recovery.

\paragraph{Classical Symbolic Planning (STRIPS/PDDL and Heuristic Search).}
STRIPS~\cite{Fikes1971STRIPS} fixed the representation: planning as search over symbolic states with action schemas, so that a plan can be checked against the model and the continuous motion layer sits behind a clean abstraction boundary. PDDL~\cite{McDermott1998PDDL} standardized domain and problem specification so that planners could be compared across navigation, logistics, and assembly, with extensions (ADL, numeric fluents, temporal constructs) that capture robotic requirements without abandoning explicit semantics. Graphplan~\cite{Blum1997PlanningGraph} builds a planning graph whose mutual-exclusion constraints prune incompatible action combinations early, and FF~\cite{Hoffmann2001FF} made forward search fast with relaxed-plan heuristics that ignore delete effects, which capture ordering constraints such as unstacking before stacking in object-rich domains. Task-and-motion planning (§\ref{subsec:p3_s7_1_4}) extends these planners with geometric feasibility checks~\cite{Kaelbling2011HierarchicalNow,Garrett2021TAMPReview,Zhao2024TAMPsurvey}; the common pattern is to keep the symbolic reasoning explicit, delegate feasibility to motion and IK solvers, and wrap the plan in an execution structure with monitors and recovery.

\paragraph{Finite-State Machines and Reactive Mode Switching.}
An FSM switches among a few well-defined behaviors (navigate, align, grasp) with explicit transition guards, which makes it transparent and fast but prone to state explosion as contingencies multiply. Statecharts~\cite{Harel1987Statecharts} add hierarchy and concurrency, so that complex controllers are nested reusable submachines with precise semantics; Brooks' subsumption architecture~\cite{Brooks1986Subsumption} layers behaviors so that higher ones override lower ones, a blueprint for composing navigation and avoidance without global planning, systematized as behavior-based robotics~\cite{Arkin1998BehaviorBased}; teleo-reactive programs~\cite{Nilsson1994TeleoReactive} evaluate conditions continuously and fire the highest-priority satisfied rule; and ROS tooling such as FlexBE~\cite{conner2023flexbe} provides hierarchical, event-driven state machines at deployment scale.

\paragraph{Behavior Trees.}
A behavior tree re-evaluates its conditions on every tick and chooses among structured alternatives (sequence, fallback, parallel), so that long-horizon execution survives disturbances without committing to one fixed sequence; subtrees encapsulate reusable routines, guards and fallbacks trigger immediate recovery, and the composition semantics fit perception and motion skills. BTs originated in game AI~\cite{Isla2005Halo2}; Marzinotto \etal~\cite{Marzinotto2014BT} formalized their semantics and their relation to other switching structures, and Colledanchise and {\"O}gren~\cite{Colledanchise2017BT,Colledanchise2018BTbook} showed which compositions preserve safety and robustness and consolidated the design principles. Compiling or wrapping a symbolic plan into a BT keeps the nominal task structure while exposing runtime guards and recovery branches~\cite{mart2021behaviourtree}; ARIAC systems pair PDDL planning with BT execution monitors to handle failures, occlusions, and defective parts~\cite{liu2024autonomous}; and event-driven BTs~\cite{Agis2020EDBT} react to non-periodic stimuli rather than waiting for the next tick.

\paragraph{Logic-based Execution and Production Rule Systems.}
A production system keeps a working memory of facts and fires forward-chaining rules when their conditions hold, which suits event-condition-action policies, mode switching, and safety enforcement (no-go zones, mutual exclusions, priorities) and depends on fast incremental matching. OPS5~\cite{Forgy1981OPS5} set the model; the Rete algorithm~\cite{Forgy1982Rete} caches partial matches so that large rule sets update in milliseconds as facts change; CLIPS~\cite{Riley1991CLIPS} made the toolchain portable; and Soar~\cite{Laird1987Soar} integrated large rule bases with planning and execution in one architecture. In robotics, KnowRob~\cite{tenorth2013knowrob} and CRAM~\cite{beetz2010cram,beetz2023cram} join symbolic knowledge to rule-driven execution for reasoning about actions, effects, and everyday manipulation, and LTLMoP~\cite{finucane2010ltlmop} compiles temporal-logic specifications into correct-by-construction reactive controllers whose safety and sequencing constraints are explicit at the program level.

\subsubsection{Benchmarks}
The International Planning Competition supplies the canonical PDDL domains and problems across classical, temporal and numeric, probabilistic, and HTN tracks, compared by plan quality, makespan, and runtime~\cite{McDermott1998PDDL}, with \texttt{planning.domains}~\cite{PlanningDomains} mirroring and extending the sets for reproducible evaluation. Long-horizon embodied suites test execution under partial observability and variation: BEHAVIOR-100 in iGibson and BEHAVIOR-1K in OmniGibson~\cite{Srivastava2022BEHAVIOR,li2024behavior1k} specify household activities at the predicate level, ALFRED~\cite{shridhar2020alfred} and TEACh~\cite{Padmakumar2022TEACh} score instruction-conditioned action sequences, and RoboTHOR~\cite{Deitke2020RoboTHOR} and the Habitat challenge~\cite{szot2021habitat2} pair high-level decision logic with navigation policies under SPL. ARIAC~\cite{liu2024autonomous} injects faulty parts, occlusions, and failures into kitting and assembly, and FurnitureBench~\cite{Heo2023FurnitureBench} and RoboCup@Home~\cite{Wisspeintner2009RoboCupHome} stress recovery and explainability on real tasks; qualitative solvers over RCC-8~\cite{Randell1992RCC} and Allen relations~\cite{janhunen2019allen} are tested on consistency problems, and composite pipelines reuse the TAMP and OMPL harnesses of §\ref{subsec:p3_s7_1_4}.

\subsubsection{Challenges and Trends}
The results on these benchmarks divide along the closed-world boundary. On IPC-style domains with fully specified preconditions and effects, classical planners remain strong; on long-horizon embodied tasks the limit is perception noise and grounding ambiguity rather than search, and a logically applicable action may be geometrically unrealizable because of a collision, an IK failure, or a missing grasp. Hand-authored operators, rules, and recovery logic also grow combinatorially with object count and contingency. Integrated task-and-motion planning (§\ref{subsec:p3_s7_1_4}) answers the coupling problem by interleaving symbolic search with feasibility checks or compiling symbolic commitments into optimization back-ends~\cite{LozanoPerez2014ConstraintTAMP,garrett2020pddlstream,toussaint2015lgp}, and hierarchical abstraction and compositional skill models contain the search while keeping it auditable~\cite{Garrett2021TAMPReview,Zhao2024TAMPsurvey}.

Robustness to noise and nonstationarity is why the industrial and instruction-conditioned benchmarks reward compositions rather than single formalisms: a behavior tree around a symbolic plan makes disturbances and recoveries explicit through runtime guards, fallbacks, and safety-preserving composition~\cite{Marzinotto2014BT,Colledanchise2017BT,Colledanchise2018BTbook}, and Rete-style engines~\cite{Forgy1982Rete} supply fast event-condition-action policies over streaming facts, the pattern ARIAC entries converge on~\cite{liu2024autonomous}.

Learned perception and language models increasingly propose the high-level procedure or heuristic while the symbolic structure stays as the safety shell that enforces preconditions, effects, and constraints: SayCan~\cite{ahn2022icanisay} grounds language in robot skills by value-based scoring, and plan-to-BT compilation~\cite{mart2021behaviourtree} exposes parallelism, guards, and recovery at run time. Domain modeling, heuristic guidance, and recovery synthesis are being automated piece by piece while the programmatic structure is retained for auditing, debugging, and verification.



\subsection{Causal Reasoning}
\label{subsec:p3_s7_1_6}
% - Cause--effect relationships between spatial events (``why'' and ``what-if'').
% - Structural Causal Models (SCM), Counterfactual SCM, PC Algorithm, GES, LiNGAM
% - Task dependency planning, event explanation, intervention & prevention, multi-robot coordination

Causal reasoning moves from describing \emph{what} happened to explaining \emph{why} it happened and predicting \emph{what would have happened} under a different action. For a robot the relevant outcomes are chains of spatial events (a collision, a delay, a deadlock, an unstable grasp, a near-miss in shared space), and making the causal assumptions explicit, rather than leaving them implicit in ``actions cause state changes'', is what lets the system attribute a failure to its causes rather than its symptoms, evaluate counterfactual alternatives such as a different approach direction or timing, and, with several agents, separate primary causes from downstream effects when choosing an intervention. Figure~\ref{fig:p3_s7_1_6_causal_reasoning_abstract} gives the abstract formulation.

\begin{figure}[!th]
    \centering
    \includegraphics[width=0.99\linewidth]{figures/part3/section7/716_causal.pdf}
    \caption{Abstract task formulation of causal reasoning for spatial intelligence: (i) event streams, candidate variables, prior knowledge, and interventional data feed (ii) methods that discover and verify causal structure, yielding (iii) causal models, explanations, counterfactual answers, and intervention or prevention strategies. The method families serve different roles rather than forming a serial pipeline. Verification returns a satisfying witness or an unsatisfiability proof/core where supported; counterfactual scenes keep the objects fixed while changing the blue agent's action.}
    \label{fig:p3_s7_1_6_causal_reasoning_abstract}
\end{figure}

\subsubsection{Task Definition}
The inputs are spatiotemporal event streams (agent actions, object states, contact and visibility relations, environment variables) mapped to causal variables by measurement functions (\texttt{in\_contact}, \texttt{blocked}, \texttt{priority\_violated}), optional prior knowledge from physics or commonsense, and, when available, interventional data in which a controllable variable was set deliberately. A structural causal model over these variables,
\begin{equation}
    \mathcal{M} = \langle \mathcal{G}, \{f_i\}, P(U) \rangle, \qquad X_i := f_i(\mathrm{Pa}(X_i), U_i),
\end{equation}
consists of a directed acyclic graph $\mathcal{G}$, a mechanism $f_i$ that generates each variable from its parents $\mathrm{Pa}(X_i)$ and exogenous noise $U_i$, and a noise distribution $P(U)$~\cite{Pearl2009Causality}. The outputs are the model itself when it must be learned, the explanation of an observed outcome (which variables caused it), the answer to an interventional query $P(Y \mid do(X=x))$ or a counterfactual one ($Y_{x}$ for a unit whose actual trace was observed), and an intervention or prevention strategy that changes controllable causes to reach a desired spatial outcome.

\subsubsection{Methodology}
The methods differ in where the causal structure comes from (assumed, learned from data, or written as domain knowledge) and in what they are asked (intervention effects, counterfactuals, responsibility, verification); in spatial domains they also differ in how they use locality, temporal precedence, and network connectivity as inductive bias.

\paragraph{Structural Causal Models and Counterfactuals.}
An SCM makes mechanisms explicit (``clearance determines collision risk'', ``traffic density sets delay'') and separates intervention from observation: the $do(\cdot)$ operator replaces a structural equation instead of conditioning on a value, which is what is needed when, say, congestion and delay are both driven by time of day. Spatial pipelines typically start from a graph fixed by temporal precedence and locality priors (edges among nearby regions, interacting agents, adjacent cells), estimate the mechanisms from data, and validate them against implied conditional independences and refutation tests~\cite{Peters2017Elements}. Counterfactuals follow the abduction-action-prediction recipe: infer the exogenous noise consistent with the observed trace, apply the hypothetical intervention, predict the outcome~\cite{Pearl2009Causality}. DoWhy and DoWhy-GCM~\cite{DoWhy2020,Bloebaum2024DoWhyGCM} implement identification, estimation, and refutation for this workflow on real spatiotemporal data.

\paragraph{Causal Discovery: Constraint- and Score-Based.}
When the graph is unknown but observational data are plentiful, constraint-based methods read structure off conditional independences: the PC algorithm builds the skeleton by independence tests and orients what it can, returning a Markov equivalence class, and FCI~\cite{Spirtes2000CPS} handles latent confounding and selection bias by returning a partial ancestral graph. Score-based methods instead search graph space under an objective such as BIC: GES~\cite{Chickering2002GES} greedily adds and then removes edges over equivalence classes, and its forward-backward structure often recovers mesoscale organization such as bottlenecks that drive downstream congestion. Spatial data strain both. Large conditioning sets are expensive and statistically weak, so practice restricts conditioning to a radius or $k$-hop neighborhood, forbids edges that violate adjacency or reachability priors, and uses order-independent variants (PC-stable~\cite{Colombo2014PCStable}) whose high-dimensional consistency is understood~\cite{Kalisch2007PC}; ambiguous orientations are then resolved by targeted interventions. causal-learn~\cite{zheng2024causal} and the TETRAD / Py-Tetrad implementations of FGES~\cite{Ramsey2018Tetrad,Ramsey2023PyTetrad} provide these pipelines at facility and city scale.

\paragraph{Independent-Noise Discovery (LiNGAM).}
In linear-Gaussian settings many edge directions are unidentifiable from independences alone. LiNGAM~\cite{Shimizu2006LiNGAM} assumes a linear structural model with independent non-Gaussian noise, under which the causal ordering is identifiable, and DirectLiNGAM~\cite{Shimizu2011DirectLiNGAM} estimates it by regression rather than ICA. Spatiotemporal signals need care because autocorrelation violates i.i.d. sampling assumptions and aggregation may change the noise distribution; temporal blocking, prewhitening, and locality constraints are the usual remedies, and the recovered ordering then scaffolds an SCM for intervention analysis.

\paragraph{Causal Action Languages and Plan-Space Causal Links.}
Much causal structure in robotics is designed knowledge rather than discovered: actions have direct and indirect effects, and the system must reason about change under incomplete information. The action language C\textsuperscript{+}~\cite{Giunchiglia2004CPlus} writes causal theories of action with inertia and nonmonotonic semantics, compactly describing how fluents such as \texttt{occupied(c)} or \texttt{visible(a,o)} evolve; CCALC~\cite{Lifschitz2000CCalc} compiles them to SAT so that prediction, diagnosis, and plan synthesis use the same solvers; and multi-valued fluents and external geometric predicates (collision checks, IK) are being integrated without abandoning the formal semantics~\cite{giunchiglia2023causal}. In plan space, partial-order causal-link planning~\cite{McAllester1991SNLP} represents a plan as partially ordered actions plus links $A \xrightarrow{p} B$ (action $A$ establishes precondition $p$ for $B$), which makes threats explicit and repairable by promotion, demotion, or binding constraints; UCPOP~\cite{Penberthy1992UCPOP} extended the representation to ADL and conditional effects and VHPOP~\cite{Younes2004VHPOP} added heuristics and reachability analysis. In spatial domains a causal link is often a physical dependency, ``unblock door'' enabling ``move cart'', which bridges symbolic causality and spatial feasibility.

\paragraph{Execution and Verification.}
A causal plan must be realized in time and checked. Simple temporal networks~\cite{Dechter1991STN} enforce difference constraints between events with polynomial consistency and dispatchable forms~\cite{Tsamardinos1998Dispatch}, and STNUs with dynamic-controllability analysis~\cite{Morris2001STNU} handle contingent durations, so that ordering and clearance constraints stay satisfiable under bounded disturbance. The SAT/SMT machinery of §\ref{subsec:p3_s7_1_4} enters here as an explanation device: when a schedule is unsatisfiable, the unsat core is a set of conflicting commitments, a concrete statement of which decisions caused the failure~\cite{Surynek2016SATMAPF,Yu2016OptimalMAPF}; and for hybrid dynamics $\delta$-complete SMT (dReal)~\cite{Gao2013dReal} and dReach~\cite{Kong2015dReach} verify safety and synthesize interventions with numerically sound proofs or counterexamples.

\subsubsection{Benchmarks}
No single benchmark covers the causal pipeline, and the suites stress different links of it. CLEVRER~\cite{Yi2020CLEVRER} pairs synthetic collision videos with descriptive, explanatory, predictive, and counterfactual questions, and CLEVRER-Humans~\cite{CLEVRERHumans2024} adds human-annotated causal event graphs against which model attributions can be compared; the video suites of §\ref{subsec:p3_s7_1_2} (PHYRE~\cite{bakhtin2019phyre}, Physion and Physion++~\cite{Bear2021Physion,tung2023physion}) score intervention selection and latent-property inference, and CRaFT~\cite{Saha2022CRaFT}, ContPhy~\cite{Zheng2024ContPhy}, NExT-QA~\cite{Xiao2021NExTQA}, AGQA~\cite{GrundeMcLaughlin2021AGQA,Gupta2023AGQA2}, and SUTD-TrafficQA~\cite{Xu2021TrafficQA} ask causal questions in natural language about synthetic and real-world videos, including traffic scenes. Embodied platforms with predicate-level state (BEHAVIOR-1K, ALFRED, iGibson 2.0~\cite{Li2021iGibson2}, ManiSkill2~\cite{Gu2023ManiSkill2}, RLBench~\cite{james2020rlbench}; §\ref{subsec:p3_s7_1_5}) test whether inferred dependencies improve intervention selection and recovery; the MAPF suites and Flatland~\cite{Mohanty2020Flatland} (§\ref{subsec:p3_s7_1_3}) test reasoning over space-time constraints in which one agent's decision causes another's delay; and ARCH-COMP~\cite{ARCHCOMP2024} with the dReal / dReach problem sets~\cite{Gao2013dReal,Kong2015dReach} evaluate causal safety reasoning on hybrid dynamics by proof or counterexample and solve time. Discovery itself is evaluated on recovery of the Markov equivalence class under faithfulness, with scalability bounded by conditioning-set size.

\subsubsection{Challenges and Trends}
Identification is brittle. Latent confounding, selection bias, and temporal aggregation can make a query non-identifiable from logs, and estimates become unstable when locality and precedence assumptions fail~\cite{Pearl2009Causality,Peters2017Elements}; constraint- and score-based discovery is expensive on dense telemetry with strong short-range dependence and sensitive to test strength, ordering, and regularization~\cite{Spirtes2000CPS,Colombo2014PCStable,Kalisch2007PC,Chickering2002GES}; and LiNGAM's non-Gaussianity is exactly what smoothed trajectories and aggregated traffic signals lack~\cite{Shimizu2006LiNGAM,Shimizu2011DirectLiNGAM}. The response is structure-aware discovery that constrains search by topology and locality and then resolves the remaining ambiguity by targeted robot interventions, connecting learning to actuation~\cite{Spirtes2000CPS,Chickering2002GES}.

Causal reasoning is only as good as its variables, and in embodied settings contact, occlusion, blockage, and object state must be grounded from noisy perception. The errors of that mapping propagate into SCM estimates and counterfactual conclusions~\cite{Peters2017Elements,DoWhy2020}. Knowledge-enabled frameworks such as KnowRob and CRAM (§\ref{subsec:p3_s7_1_5}) show the shape of the interface, perception populating an explicit, interpretable model whose parameters are then estimated statistically, and the trend is toward such perception-to-symbol interfaces that tolerate sensor noise~\cite{tenorth2013knowrob,beetz2010cram,beetz2023cram}.

Between a causal recommendation and a safe execution there is a verification-synthesis gap: the recommendation must become a temporally controllable schedule and a certified controller, and execution uncertainty invalidates a symbolic causal plan unless temporal networks and hybrid verification are integrated with it~\cite{Dechter1991STN,Morris2001STNU}. SCM toolchains are therefore being embedded in decision stacks so that counterfactual policies compile into action-language or plan-space structures and then into dispatchable networks~\cite{Bloebaum2024DoWhyGCM,Giunchiglia2004CPlus,Morris2001STNU}, with SAT/SMT certificates and unsat-core explanations supplying the accountability~\cite{Surynek2016SATMAPF,Gao2013dReal,Kong2015dReach}; evaluation that scores explanation quality, interventional effectiveness, and certification cost on the same task is what the field still lacks.



% \subsection{Summary}
% \label{subsec:p3_s7_1_summary}

% The six paradigms in §\ref{sec:p3_s7_1} build spatial intelligence on explicit, inspectable structure (geometric, mechanical, topological, formal, programmatic, or causal) rather than on learned correlations alone, and that choice buys four things. Plans, schedules, and explanations are artifacts that can be debugged and verified rather than distributions to be sampled. Each paradigm increasingly produces machine-checkable certificates (force-closure conditions, optimality witnesses, $\delta$-complete SMT proofs, unsat cores) on which safety-critical deployment can rely. Explicit interfaces (operators, constraints, rules, causal variables) let the paradigms compose with one another and with learned components without losing their guarantees. And where the model captures the relevant structure, classical methods generalize from a compact specification rather than from a large corpus.

% Table~\ref{tab:p3_s7_1_paradigm_comparison} sets the six side by side; they are complementary rather than redundant, each targeting a structure that the others presuppose (geometry beneath physics), wrap (behavior trees around symbolic plans), or extend (causal mechanisms over constraint solutions).

% \begin{table}[t]
%     \centering
%     \footnotesize
%     \caption{Comparison of the six foundational paradigms of rule-based symbolic spatial reasoning covered in §\ref{sec:p3_s7_1}.}
%     \label{tab:p3_s7_1_paradigm_comparison}
%     \begin{tabularx}{\linewidth}{@{}l X X X X@{}}
%         \toprule
%         \textbf{Paradigm} & \textbf{Core representation} & \textbf{Output / artifact} & \textbf{Notable strength} & \textbf{Main limitation} \\
%         \midrule
%         §\ref{subsec:p3_s7_1_1} Geometric \& algebraic         & $SE(3)$ poses, projective geometry, Jacobians, contact wrenches                  & Frame-aligned coordinates, workspace / reachability sets, feasibility certificates      & Closed-form precision; coordinate-free formulations                          & Brittle near singularities; assumes exact geometry \\
%         §\ref{subsec:p3_s7_1_2} Physical                       & Rigid-body dynamics, contact / friction laws, quasi-static equilibrium           & Force-closure verdict, motion cone, support-polygon test                                & Grounded in mechanics; analytically interpretable                            & Sim-to-real gap; latent parameters (mass, friction) often unobserved \\
%         §\ref{subsec:p3_s7_1_3} Topological / graph-based      & Roadmap / decomposition / sampling graph $G=(V,E,\mathbf{w})$                    & Collision-free path $\pi$, topological map, multi-agent trajectories                    & Combinatorial tractability; mature solvers (A*, CBS, CH)                     & Density vs.\ optimality trade-off; appearance robustness across seasons \\
%         §\ref{subsec:p3_s7_1_4} Constraint-based               & Decision variables + collision / IK / dynamics / logic constraints               & Certified-feasible plan, optimality gap, infeasibility witness                          & Strong correctness guarantees; certificate-producing                         & Mixed-integer scaling; symbol-grounding from perception \\
%         §\ref{subsec:p3_s7_1_5} Rule-based / programmatic      & STRIPS / PDDL operators, FSMs, Behavior Trees, production rules                  & Action plan, reactive policy, monitored execution trace                                 & Auditable; reactive recovery; modular composition                            & Modeling burden; brittle under perception noise without grounding \\
%         §\ref{subsec:p3_s7_1_6} Causal                         & Structural causal models, temporal networks (STN / STNU)                         & Causal model, $do(\cdot)$-effect estimate, counterfactual answer, intervention strategy & Counterfactual / responsibility queries; identifiability theorems            & Grounding causal variables; latent confounding; LiNGAM needs non-Gaussian noise \\
%         \bottomrule
%     \end{tabularx}
% \end{table}

% The same design philosophy exposes the limitations. Explicit models must be authored and maintained, they presume that the relevant structure is observable, and they degrade gracelessly when reality violates their idealizations (perception noise, latent parameters, unmodeled contact, drifting causal mechanisms). These are the conditions under which data-driven methods excel: learned pipelines absorb unobservable variables and perception noise into end-to-end objectives and generalize across regimes that would be intractable to specify by hand.

% The rest of Chapter~\ref{sec:reasoning} takes up spatial reasoning from that complementary angle: §\ref{sec:p3_s7_2} on reasoning learned from data, §\ref{sec:p3_s7_3} on reasoning elicited from foundation models, §\ref{sec:p3_s7_4} on hybrids of the two with the classical paradigms, and §\ref{sec:p3_s7_5} on how the field measures all of them. Throughout, the classical paradigms of §\ref{sec:p3_s7_1} reappear as feasibility oracles, safety shells, structural priors, and explainability layers, and the most effective systems are those that compose the two regimes: explicit structure where it is authoritative, learned structure where it is unavoidable.


% \newpage

\section{Learning-based Spatial Reasoning}
\label{sec:p3_s7_2}

\FramedInitial{L}{\textbf{earning-based spatial reasoning}} refers to the class of methods that acquire and perform spatial reasoning from data, rather than from the hand-authored geometric, mechanical, or symbolic models surveyed in §\ref{sec:p3_s7_1}.
The six subsections that follow are ordered by the \emph{locus of the structural prior}: it begins inside the learned representation (latent feature interaction, explicit programs, hard-coded symmetries, optimization layers in the forward pass), moves to the agent's action-conditioned behavior (reinforcement learning with spatial inductive biases), and ends with joint decision-making across coordinated agents.



\subsection{Implicit Neural Spatial Reasoning}
\label{subsec:p3_s7_2_1}



We define \emph{implicit neural spatial reasoning} as the family of neural approaches in which spatial relations are encoded in learned feature representations, attention patterns, latent object interactions, or architectural inductive biases, rather than specified through hand-written rules, symbolic programs, or fully explicit relational graphs.
\emph{Implicit} does not mean that these models never use boxes, masks, object proposals, or task-level labels; it means that the reasoning procedure itself is never given as a symbolic program, so the model must learn spatial dependencies from end-to-end or weakly structured supervision.
Several families below (latent scene graphs, capsule networks) expose partly inspectable structure, yet remain implicit in that their reasoning representations are learned rather than fixed by hand-written symbolic rules.
Figure~\ref{fig:p3_s7_2_1_implicit_formulation} shows the representative tasks.


\begin{figure}[!th]
    \centering
    \includegraphics[width=0.99\linewidth]{figures/part3/section7/implicit_neural_spatial_reasoning_task.pdf}
    \caption{Representative tasks in which spatial relations must be inferred without explicit symbolic supervision: (1) visual question answering over compositional scenes, (2) image captioning with spatial attention, and (3) visual relational reasoning.}
    \label{fig:p3_s7_2_1_implicit_formulation}
\end{figure}




\subsubsection{Task Definition}
Given one or more images $I$, optionally paired with a language query or task goal $q$ (and possibly boxes, masks, or proposals as features), an implicit reasoner computes a latent spatial representation and a task output,
\begin{equation}
  z=f_\theta(I,q),\qquad y=h_\psi(z),
\end{equation}
where $z$ encodes the relations relevant to the task and $y$ is an answer label, a caption, or a relation label.
In a purely implicit setting, supervision reaches the model through $y$; semi-explicit variants may additionally use object or relation labels without prescribing a symbolic reasoning program.
The methods below differ in the mechanism by which $z$ is formed (Figure~\ref{fig:p3_s7_2_1_implicit_method}).




\begin{figure}[!th]
    \centering
    \includegraphics[width=0.99\linewidth]{figures/part3/section7/implicit_neural_spatial_reasoning_method.pdf}
    \caption{Four mechanisms for learning spatial representations: (a) relational aggregation (Relation Network, GNN message passing), (b) spatial attention with multi-scale fusion (FPN), (c) conditional modulation (FiLM), and (d) geometry-aware inductive biases (capsule routing, equivariant features). Blue marks selected features, messages, and operations. The mechanisms can be combined in one model, and attention and FPN support spatial representation rather than forming complete reasoning procedures.}
    \label{fig:p3_s7_2_1_implicit_method}
\end{figure}

\subsubsection{Methodology}

\paragraph{Relational Aggregation over Latent Objects.}
The most direct approach treats a set of feature vectors as latent ``objects'' and learns to aggregate their interactions.
\textbf{Relation Networks (RN)}~\cite{santoro2017simple} score every pair: given objects $O=\{o_1,\dots,o_n\}$ (\eg, feature vectors from spatial regions of a CNN map) and an optional query $q$,
\begin{equation}
  \text{RN}(O, q) = f_{\phi}\!\left(\sum_{i,j} g_{\theta}(o_i, o_j, q)\right),
\end{equation}
where $g_{\theta}$ produces a query-conditioned relation embedding for each pair and $f_{\phi}$ maps the sum to the task output.
Considering all pairs forces the network to learn pairwise spatial and semantic relations, at $O(n^2)$ cost in the number of objects.
\textbf{Graph neural networks and latent scene graphs} generalize this from a complete graph to a learned, sparser topology, with proposals or feature regions as nodes and learned edges as spatial dependencies; iterative message passing for scene-graph generation~\cite{xu2017scene} alternates between refining node features and updating relation embeddings, learning the topology end-to-end.
Such models are the clearest \emph{semi-explicit} case in this subsection: the graph is inspectable and often trained with relation annotations (\eg, from Visual Genome), yet the intermediate relational representation is learned rather than specified by a program, and its quality is bounded by the upstream proposal stage.

\paragraph{Spatial Attention and Multi-Scale Feature Fusion.}
A second family controls \emph{where} and \emph{how strongly} the model attends across the spatial layout, and fuses evidence across scales, rather than scoring all pairs.
CNN spatial attention weights locations in a feature grid, from the hard and soft attention of \citet{xu2015show} for captioning and grounded VQA~\cite{zhu2016visual7w} to SCA-CNN's joint spatial and channel attention~\cite{chen2017sca}; vision transformers make relating image parts intrinsic through self-attention over patches~\cite{vit}, with locality-aware variants (Swin~\cite{liu2021swin}, PVT~\cite{wang2021pyramid}, Twins~\cite{chu2021twins}) trading global receptive field for cost; and Feature Pyramid Networks~\cite{lin2017feature} add a top-down pathway so that a relation coupling a fine object with a coarse region (a car ``on'' a street) can be read from one map.
These designs are representation backbones that preserve and aggregate spatial information for the reasoning built on top of them, not reasoning modules in themselves.

\paragraph{Conditional Feature Modulation.}
A third family conditions one representation on another rather than aggregating or attending.
In \textbf{FiLM (Feature-wise Linear Modulation)}~\cite{perez2018film}, a conditioning vector (\eg, an encoded query or global context) predicts affine parameters $(\gamma, \beta)$ that are applied feature-wise to a target map, so the same backbone can be steered per task.
For a query like ``What is to the left of the green cube?'', a FiLM layer can learn to amplify features to the left of regions identified as the green cube, making reasoning adaptive and context-dependent.

\paragraph{Geometry-Aware Inductive Biases.}
A fourth family bakes geometric regularities into the architecture, trading data-driven flexibility for sample efficiency where the relevant geometry is known.
\textbf{Capsule Networks}~\cite{sabour2017dynamic} replace scalar activations with vector-valued capsules whose orientation encodes instantiation parameters (pose, deformation) and whose magnitude encodes presence; \emph{dynamic routing} decides how lower-level capsules compose into higher-level ones by agreement in predicted pose, a stronger part-whole bias than translation-only convolution, though capsules have been hard to scale beyond small digit and shape datasets.
Equivariant and steerable convolutions~\cite{cohen2016group}, which commute with a chosen transformation group, are the subject of §\ref{subsec:p3_s7_2_3}; both share the cost that the symmetry must be chosen in advance.


\subsubsection{Benchmarks}

CLEVR~\cite{johnson2017clevr} is the canonical benchmark, pairing rendered scenes with multi-step relational questions whose ground-truth program structure is known; Sort-of-CLEVR~\cite{santoro2017simple} isolates a controlled relational vocabulary.
The decisive stress tests are the \emph{systematic-generalization} splits: CLEVR-CoGenT~\cite{johnson2017clevr} holds out attribute-composition combinations, and GQA~\cite{hudson2019gqa} scales compositional questions to real images.
Open-domain VQA v2~\cite{goyal2017making} and Visual7W~\cite{zhu2016visual7w} probe whether implicit reasoning transfers beyond synthetic templates to natural images and free-form language.
On CLEVR, RN reaches 95.5\% and FiLM 97.7\%; Table~\ref{tab:p3_s7_2_2_clevr_gqa} in §\ref{subsec:p3_s7_2_2} places these against the program-structured reasoners, which close the remaining gap only on in-distribution questions.


\subsubsection{Challenges and Trends}
The \emph{interpretability challenge} is intrinsic: because relations live in distributed activations rather than named symbols, it is hard to audit which relations a model actually exploits, and attention-map visualizations are only an indirect proxy.
\emph{Compositional generalization} is the most striking failure: even at 97--98\% CLEVR accuracy, the same models drop sharply on CLEVR-CoGenT~\cite{johnson2017clevr} and GQA~\cite{hudson2019gqa}, revealing dataset-specific shortcuts rather than transferable spatial concepts; the related \emph{shortcut-learning} concern in open-domain VQA, where language priors and object co-occurrence dominate the signal, is the imbalance VQA~v2's balanced design targets~\cite{goyal2017making}.
\emph{Scalability} divides the families: RN-style pairwise modeling is $O(n^2)$ in the number of objects, attention backbones trade receptive field against compute~\cite{liu2021swin,wang2021pyramid}, and GNNs inherit the brittleness of the proposal stage~\cite{xu2017scene}.
The more fundamental limit is the \emph{2D-to-3D} gap: these methods reason over 2D feature grids and image-plane relations, whereas much spatial reasoning depends on depth ordering, occlusion, viewpoint, and metric layout, and strong 2D backbones and vision--language models degrade sharply once relations hinge on 3D geometry~\cite{daxberger2025mmspatial,xu2025spatialbench}.

Two responses run through the rest of this section.
The explicit-structured methods of §\ref{subsec:p3_s7_2_2} reintroduce programs and graphs over learned features, recovering inspectability while keeping the representational power of deep networks; the geometry-aware methods of §\ref{subsec:p3_s7_2_3} answer the 2D-to-3D gap directly by building the symmetry group into the architecture, where it is known.
Whether web-scale pretraining alone closes the systematic-generalization gap, or explicit structure remains necessary, is the open question that §\ref{sec:p3_s7_3} takes up.








\subsection{Explicit Structured Neural Reasoning}
\label{subsec:p3_s7_2_2}


% IEP/EE 本质上是 program-structured / neuro-symbolic 方法，更适合放在下一节 explicit structured reasoning

Explicit structured neural reasoning decomposes a spatial problem into an inspectable intermediate structure, a program, a scene graph, a sequence of reasoning states, or executable code that calls vision tools, and obtains the answer by executing or traversing that structure over visual representations.
Unlike the rule-based reasoning of §\ref{sec:p3_s7_1}, the structure is learned or inferred from data rather than written by hand; unlike the implicit reasoners of §\ref{subsec:p3_s7_2_1}, it exposes spatial predicates directly, so that an inspectable edge or a named operation can represent containment, adjacency, support, or relative position, making them auditable and composable instead of entangled in a feature space.
Figure~\ref{fig:p3_s7_2_2_explicit_formulation} previews the representative tasks.

\begin{figure}[!th]
    \centering
    \includegraphics[width=0.99\linewidth]{figures/part3/section7/explicit_structured_neural_reasoning_task.pdf}
    \caption{Representative tasks for explicit structured neural reasoning, in which the reasoning trace is inspectable: (1) program-style VQA on compositional scenes, (2) state-machine traversal of a probabilistic scene graph, and (3) execution of an LLM-generated visual program.}
    \label{fig:p3_s7_2_2_explicit_formulation}
\end{figure}




\subsubsection{Task Definition}
Given a visual observation $\mathbf{I}$ (a single image or a set of views) and a query or goal $q$, the system constructs an intermediate structure and executes it,
\begin{equation}
  \mathcal{S}=\mathrm{Parse}_\phi(\mathbf{I},q),\qquad y=\mathcal{E}(\mathcal{S};\mathbf{I}),
\end{equation}
where $\mathcal{S}$ may be a program layout $P=(m_1,\dots,m_k)$, a recurrent sequence of reasoning states, a scene graph $G=(V,E)$ with a finite-state traversal over its nodes, or a code trace that calls external tools, and the executor $\mathcal{E}$ realizes $\mathcal{S}$ through neural modules, symbolic operators, graph transitions, or tool calls.
The output $y$ may be an answer, a localized region, a predicted relation, or an action.
The defining property is therefore not a specific program syntax but the presence of an inspectable $\mathcal{S}$ that mediates between perception and prediction, which is what makes the reasoning path explicit and controllable.




\subsubsection{Methodology}

The four families in Figure~\ref{fig:p3_s7_2_2_explicit_formulation_method} differ in how much structure they expose, enough to make reasoning inspectable and compositional, but not so rigid that a perception error, a parser error, or a missing primitive causes deterministic failure.
Each exposes a different facet of space: recurrent architectures a temporal trace of attentional focus, scene-graph methods object-centric relations, program-based methods named operations such as filtering, relating, counting, and comparing, and LLM-generated programs tool-call traces that combine detection, segmentation, localization, and external knowledge.


\paragraph{Structured Recurrent Reasoning Architectures.}
This family imposes a step-wise reasoning structure in the architecture itself, so the model is fully differentiable and trained end-to-end from input--output pairs while still exposing a step-by-step trace.
The Memory, Attention, and Composition (MAC) network~\cite{hudson2018compositional} chains $p$ recurrent cells, each of which attends to the question to choose the operation for the current step (control unit), extracts the relevant visual information by spatial attention guided by that control signal (read unit), and integrates it into a reasoning state passed to the next cell (write unit); the attention maps of successive read units provide a transparent trace of where the model is looking, without a rigid symbolic program.
MAC is thus a boundary case: explicit at the level of reasoning dynamics, implicit at the level of spatial representation, since the operations run over dense feature maps rather than named objects or relations.
Because the spatial logic stays in the learned weights, it is prone to statistical shortcuts (if ``red spheres'' co-occur with ``cubes'' in training, the model may learn that correlation rather than ``left of''), and soft attention maps are often diffuse, so ``to the left'' bleeds into irrelevant regions in cluttered scenes.
\begin{figure}[!th]
    \centering
    \includegraphics[width=0.99\linewidth]{figures/part3/section7/explicit_structured_neural_reasoning_method.pdf}
    \caption{Four method families for explicit structured neural reasoning, ordered by the intermediate structure $\mathcal{S}$ each exposes: (A) the MAC recurrent reasoning cell~\cite{hudson2018compositional}; (B) the Neural State Machine~\cite{hudson2019learning}; (C) program-based reasoning, a parse-to-execute pipeline whose executor ranges from neural modules (NMN~\cite{andreas2016neural}, IEP~\cite{johnson2017inferring}) through the symbolic executor of NS-VQA~\cite{yi2018neural} to the quasi-symbolic executor of NS-CL~\cite{mao2019neuro}; and (D) LLM-as-program-generator pipelines that call vision tools (\eg,~ViperGPT~\cite{suris2023vipergpt}, VisProg~\cite{gupta2023visual}, Code-as-Policies~\cite{liang2023code}).}
    \label{fig:p3_s7_2_2_explicit_formulation_method}
\end{figure}

\paragraph{Structured Scene Representations.}
A related line strengthens the structural prior on the reasoning substrate.
Where MAC reasons over a soft grid of features, the Neural State Machine (NSM)~\cite{hudson2019learning} reasons over a probabilistic semantic scene graph and performs differentiable state transitions over it: a concept-detection module (Mask R-CNN) identifies entities, each node $v_i$ is initialized with semantic properties over a learned concept vocabulary, each edge $e_{ij}$ carries a distribution over spatial and semantic relation types (left of, behind, on top of, holding), and the query is answered by propagating a distribution over nodes through query-conditioned transitions.
More broadly, scene graphs act as a generic structured substrate: a downstream reasoner can propagate probabilistic states over the graph (as in NSM), attend over it as a soft relational prior, or search it for graph-based question answering, in each case obtaining an inspectable cut-point between perception and reasoning.
The price is that reasoning quality becomes tightly coupled to object detection, attribute grounding, and relation prediction, a dependency we return to in the challenges.


\paragraph{Modular and Program-Based Reasoning.}
A third family translates the query itself into an executable program, and its members differ mainly in the nature of the \emph{executor}.
The Neural Module Network (NMN)~\cite{andreas2016neural} is the canonical \emph{neural-executor} instance: an external semantic parser converts a question such as ``What color is the llama?'' into a layout, \eg,~\texttt{classify[color](attend[llama])}, whose slots are instantiated by distinct learnable modules (attention, re-attention, classification, combination, measurement) trained jointly to maximize the likelihood of the gold answer.
Because NMN relies on an off-the-shelf parser, IEP~\cite{johnson2017inferring} replaces it with a sequence-to-sequence program generator that emits the program as a sequence of functions, trained jointly with the execution modules via REINFORCE~\cite{williams1992simple} or by maximizing answer likelihood over a distribution of programs.
Moving further toward symbols, Neuro-Symbolic VQA (NS-VQA)~\cite{yi2018neural} adopts a \emph{symbolic executor}: a scene parser turns the image into a symbolic table, after which the program reduces to deterministic functional operators, yielding fully interpretable execution at the cost of a near-perfect-perception assumption.
The Neuro-Symbolic Concept Learner (NS-CL)~\cite{mao2019neuro} occupies the middle ground with a \emph{quasi-symbolic} executor: a neural perception module produces object-centric representations and a semantic parser emits symbolic programs, but the executor relaxes discrete operations such as filtering and relating into differentiable ones, so concept learning and semantic parsing can be trained from question--answer pairs, while the object detector is pretrained using bounding-box annotations.

\paragraph{LLM-Generated Visual Programs.}
The most active recent extension lets a large language model emit the program directly, replacing the seq2seq parser of IEP or the hand-crafted parser of NMN with a \emph{tool-based executor}.
ViperGPT~\cite{suris2023vipergpt} and VisProg~\cite{gupta2023visual} pair a frozen LLM with a Python-style API of vision tools (detectors, captioners, attribute classifiers) that the generated code invokes, making the approach compositional and zero-shot without per-task training.
Code-as-Policies~\cite{liang2023code} pushes the same program-synthesis paradigm into embodied control, treating the LLM-emitted program as a closed-loop policy over perception and actuation primitives; it is included here as an embodied extension of the paradigm rather than as a VQA method.
The interpretability gain is real, since the emitted code is human-readable, but these systems re-inherit the brittleness of the underlying LLM (hallucinated API calls, ungrounded tool routing) as a failure mode the earlier hand-coded DSLs were designed to avoid.

\paragraph{Unifying View: Parser, Representation, Executor, and Supervision.}
The four families are points in a shared design space with four axes.
The \emph{parser} that maps a query to structure may be hand-built (NMN), learned end-to-end (IEP, NS-CL), or a frozen LLM (ViperGPT, VisProg); the \emph{representation} over which reasoning runs may be a dense feature map (MAC), an object-centric scene graph (NSM, and the scene-graph-generation literature~\cite{xu2017scene} supervised at scale by Visual Genome~\cite{krishna2017visual}), or a symbolic scene table (NS-VQA); the \emph{executor} may be neural, symbolic, quasi-symbolic, or a code interpreter; and the \emph{supervision} may be answer-only, program-supervised, or, in tool-based systems, augmented with execution traces and tool feedback.
The inspectable structure between perception and answer is the common thread: it provides a natural cut-point at which symbolic and learned components can be swapped without retraining the rest of the stack, and the same four axes reappear in the foundation-model reasoners of §\ref{sec:p3_s7_3}.

\paragraph{Beyond VQA: 3D, Temporal, and Embodied Extensions.}
The same design pattern extends beyond static images: referring-expression grounding has been lifted from 2D images to 3D scenes~\cite{achlioptas2020referit3d}, where object-centric and spatial-relation graphs supply the substrate for reasoning about support, containment, proximity, and navigability, and in video and embodied settings the intermediate structure becomes a temporal scene graph or a task plan that links perception to action (§\ref{sec:p3_s7_3}, Part~IV).
These extensions also expose the central limitation of current explicit methods: most available program primitives and graph predicates are still designed for static 2D relations, whereas viewpoint change, object permanence, physical interaction, and multi-step navigation need 3D, temporal, and causal operators.


\subsubsection{Benchmarks and Performance}

The CLEVR family~\cite{johnson2017clevr} of §\ref{subsec:p3_s7_2_1} is the canonical testbed here as well, with the added property that each question is generated alongside a ground-truth functional program, which gives both program-supervised and answer-only training signals; MAC, IEP, NS-CL, and NS-VQA all converge to the high-90s on it, making it attractive but saturated.
GQA~\cite{hudson2019gqa} is the dominant open-domain benchmark: it pairs natural-image scenes with compositional questions whose ground-truth scene graphs (from Visual Genome~\cite{krishna2017visual}) make scene-graph-based reasoners directly applicable, and its out-of-distribution splits expose the systematic-generalization gap hidden by saturated CLEVR scores.
VQA v2~\cite{goyal2017making} is a general VQA benchmark, while the cited visual-program systems ViperGPT and VisProg report task-specific comparisons on GQA.
Referring-expression grounding benchmarks (RefCOCO and RefCOCO+~\cite{yu2016modeling}, RefCOCOg~\cite{mao2016generation}, and the earlier ReferItGame~\cite{kazemzadeh2014referitgame}) are equally relevant, since they test whether a system can bind spatial language (``the object left of / behind / next to \dots'') to a localized region even when the output is a box or mask rather than an answer.
Embodied program-as-policy systems such as Code-as-Policies are evaluated by task success in robotic or simulated manipulation environments and should not be compared against VQA accuracy.
Table~\ref{tab:p3_s7_2_2_clevr_gqa} collects the headline numbers.

\begin{table}[!ht]
    \centering
    \caption{Implicit and explicit structured reasoners on CLEVR and GQA (overall accuracy, \%, higher is better). CLEVR numbers use the evaluation splits reported in the original papers; NS-CL reports validation accuracy after training on 5K images; GQA numbers are test-dev or test-standard as reported. ViperGPT and VisProg run zero-shot with a frozen LLM, and VisProg is evaluated on a curated test-dev subset, so its number is not directly comparable. RN and FiLM are the implicit reference points of §\ref{subsec:p3_s7_2_1}.}
    \label{tab:p3_s7_2_2_clevr_gqa}
    \fontsize{9}{11}\selectfont
    \setlength{\tabcolsep}{1.5mm}
    \begin{tabular}{L{0.22\textwidth} L{0.26\textwidth} L{0.20\textwidth} C{0.10\textwidth} C{0.10\textwidth}}
        \toprule
        \textbf{Method} & \textbf{Structure exposed} & \textbf{Supervision} & \textbf{CLEVR} & \textbf{GQA} \\
        \midrule
        RN~\cite{santoro2017simple}            & none (pairwise latent)              & answers                       & 95.5 & --   \\
        FiLM~\cite{perez2018film}              & none (feature modulation)           & answers                       & 97.7 & --   \\
        MAC~\cite{hudson2018compositional}     & recurrent control / read / write    & answers                       & 98.9 & 54.1 \\
        IEP~\cite{johnson2017inferring}        & generated program, neural executor  & programs (700k) $+$ answers   & 96.9 & --   \\
        NS-CL~\cite{mao2019neuro}              & program, quasi-symbolic executor    & boxes (detector) $+$ answers (curriculum) & 98.9 & --   \\
        NS-VQA~\cite{yi2018neural}             & program, symbolic executor          & scene tables $+$ 270 programs & 99.8 & --   \\
        NSM~\cite{hudson2019learning}          & probabilistic scene graph, automaton & answers                      & --   & 63.2 \\
        ViperGPT~\cite{suris2023vipergpt}      & LLM-generated code, tool executor   & none (zero-shot)              & --   & 48.1 \\
        VisProg~\cite{gupta2023visual}         & LLM-generated code, tool executor   & none (zero-shot)              & --   & 50.5 \\
        \bottomrule
    \end{tabular}
\end{table}


\subsubsection{Challenges and Trends}
Exposing programs, graphs, state transitions, or tool calls buys interpretability but creates new failure modes, and the numbers above locate them.
The first bottleneck is \textbf{parsing}: if a question is mapped to the wrong program or traversal path, even a correct executor produces the wrong answer.
The second is \textbf{symbol grounding and underspecification}: spatial predicates such as \texttt{Left}, \texttt{Behind}, \texttt{Near}, or \texttt{Between} depend on viewpoint, scale, object extent, and task context, so many spatial queries do not determine a unique program, a single deterministic DSL operator is often too rigid, and new predicates demand manual re-engineering of the executor.
The third is \textbf{error propagation}: mistakes by the detector, scene parser, relation classifier, or program generator are inherited by the executor and can cause deterministic failure, which is why NSM reaches 63\% on GQA while the same family saturates CLEVR; on natural images the bottleneck has moved from reasoning execution to perception, scene-graph construction, language grounding, and long-tail relation coverage.
A fourth, more structural tension is the \textbf{differentiability--interpretability trade-off}: the more symbolic the executor, the more auditable but the harder to train end-to-end; the more differentiable, the more robust but the less transparent.
Because the value of explicit structure lies in the intermediate trace, \textbf{evaluation} should move beyond final-answer accuracy to whether the parsed program is faithful to the query, whether graph edges correspond to real relations, whether tool calls are necessary and grounded, and whether the execution trace actually supports the answer.

Most explicit systems also reason over static 2D layouts or object graphs rather than 3D geometry, temporal dynamics, or physical causality: they can answer whether $A$ is left of $B$, but rarely simulate how a scene changes under motion, contact, containment, support, or occlusion, which is exactly the mental-simulation capacity human-level spatial reasoning requires.
LLM-generated programs widen the scope by replacing fixed parsers with flexible code generation, trading a small accuracy gap on saturated benchmarks for zero-shot generality across VQA, grounding, and knowledge-grounded VQA, but they add hallucinated tool calls, invalid APIs, ungrounded intermediate variables, and the absence of a verifiable execution loop.
The emerging direction is therefore not a return to purely symbolic reasoning but \emph{verifiable hybrid} systems, in which learned perception produces object-centric representations, structured programs or graphs provide an inspectable scaffold, and neural or tool-based executors absorb the variability of real visual inputs, the direction of §\ref{sec:p3_s7_3} and §\ref{sec:p3_s7_4}.









\subsection{Geometry-aware Neural Spatial Reasoning}
\label{subsec:p3_s7_2_3}
Where the implicit and explicit methods above learn spatial relationships from data, geometry-aware approaches build the relevant symmetries into the network itself: an object rotated, translated, or re-indexed keeps its identity, and the architecture is constructed so that this holds by design rather than being learned from many transformed examples.
The gain is sample efficiency and robustness under the chosen transformations; the cost is that the symmetry group has to be chosen in advance.
Figure~\ref{fig:p3_s7_2_3_geometry_formulation} pairs three symmetry-driven task settings with the method families that supply the matching inductive bias.

\begin{figure}[t]
    \centering
    \includegraphics[width=\linewidth]{figures/part3/section7/geometry-aware_neural_spatial_reasoning.pdf}
    \caption{Symmetry-driven tasks (top) and geometric neural mechanisms (bottom). Top: (i) rotation-invariant classification, (ii) SE(3)-equivariant grasping, shown with a planar rotation, and (iii) a molecular scalar invariant to rigid motion and atom reindexing. Bottom: (a) STN~\cite{jaderberg2015spatial}, (b) G-CNN~\cite{cohen2016group}, (c) steerable CNN~\cite{cohen2017steerable}, (d) PointNet~\cite{qi2017pointnet}, and (e) TFN / SE(3)-Transformer~\cite{thomas2018tensorfield,fuchs2020se}. Blue marks transformations, geometric features, and preserved outputs; all examples are schematic.}
    \label{fig:p3_s7_2_3_geometry_formulation}
\end{figure}

\subsubsection{Task Definition}
The input $x$ lives in a space with a known geometric structure (a 2D image, a 3D point cloud, a molecular graph), optionally with a query $q$, and the network $M_{\theta,G}$ is designed to respect a chosen group $G$ of transformations: translations for standard CNNs, rotations SO(3), rigid motions SE(3), or permutations of set elements.
Its output $\mathbf{z}$ is either invariant or equivariant to the action of $G$.

\begin{quotebox}
\noindent\textbf{Invariance vs.\ equivariance.} An \emph{invariant} representation stays fixed under the group action, $M_{\theta,G}(g\cdot x)=M_{\theta,G}(x)$, and deliberately discards pose; an \emph{equivariant} representation co-transforms with the input, $M_{\theta,G}(g\cdot x)=\rho(g)\cdot M_{\theta,G}(x)$, where $\rho(g)$ is the representation of the group action on the feature space, and thereby preserves it. In either case geometric consistency is a native constraint that the architecture guarantees by construction.
\end{quotebox}

\subsubsection{Methodology}

Invariance is the right target when the output should not depend on orientation or position (classification), equivariance when pose and relative arrangement are the point (manipulation, keypoints, instance segmentation, molecular modeling); the layers below realize each.

\paragraph{Invariance by Pooling and Augmentation.}
The baseline strategies work without changing the architecture.
Augmenting the training set with rotated, flipped, or scaled copies pushes a standard CNN toward assigning the same label to all variants, but buys only an approximate invariance and requires examples of every transformation to be tolerated; global average pooling~\cite{lin2013network} at the end of a feature extractor discards the ``where'' of each feature and keeps only its presence, which is the main reason classification CNNs are robust to translation.
Neither guarantees invariance to anything beyond translation.

\paragraph{Active Alignment with Spatial Transformer Networks.}
Rather than ignoring transformations, a Spatial Transformer Network (STN)~\cite{jaderberg2015spatial} learns to undo them: a small localization network regresses the parameters of an affine transformation that it predicts will bring the object into a canonical pose, a grid generator turns these parameters into a sampling grid, and a differentiable sampler produces the rectified feature map, so the localization network is trained by backpropagation from the task loss alone.
The module can be inserted at any point in a network and acts as learned attentional alignment, normalizing a region of interest for the layers downstream; the invariance it provides is learned, not guaranteed.

\paragraph{Permutation Invariance over Point Sets.}
A separate symmetry arises when the input is an unordered set rather than a grid: the model must be invariant to the ordering of its elements.
DeepSets~\cite{zaheer2017deep} shows that any permutation-invariant function over a set can be written as $f(X)=\phi\!\left(\sum_{x\in X}\psi(x)\right)$ for suitable $\phi$ and $\psi$, encode each element independently, sum, then decode, which turns set-level invariance from an ad-hoc choice into a constructive recipe and makes explicit which aggregators preserve invariance and which silently break it; the same framework yields permutation-\emph{equivariant} layers for per-element outputs (per-point labels, per-atom properties).
PointNet~\cite{qi2017pointnet} is the canonical instance on raw point clouds, a shared multilayer perceptron per point followed by symmetric max pooling, and PointNet++~\cite{qi2017pointnetplus} adds locality by repeatedly grouping points within a metric ball and applying the same shared-MLP-plus-pool pattern at multiple scales.
Permutation invariance comes for free in this family; rigid-motion invariance does not, which is the gap the equivariant methods below address.

\paragraph{From Translation to Group Equivariance.}
A standard convolution is translation-equivariant (shifting the image shifts the feature map) but not rotation-equivariant: rotating the input changes the feature map unpredictably.
Group equivariant convolutions (G-CNNs)~\cite{cohen2016group} generalize the operation to a larger discrete group such as the four planar rotations: a single filter is rotated into a set of transformed copies, all are correlated with the input, and the feature map gains an orientation dimension, so a feature that reads ``edge at $45^\circ$'' becomes ``edge at $135^\circ$'' at the corresponding location when the input rotates by $90^\circ$.
Equivariance is therefore not learned but imposed, layer by layer.

\paragraph{Steerable Filters and Equivariance in Every Layer.}
Steerable CNNs~\cite{cohen2017steerable} replace the group-indexed filter copies by filters expressed in a steerable basis, so that the response at a transformed orientation follows by linear combination of basis responses rather than by explicit rotation; the general E(2) framework~\cite{weiler2019general} extends this to continuous rotations and reflections with circular-harmonic bases under one configurable interface, and library efforts such as e3nn~\cite{geiger2022e3nn} do the same for E(3) / SE(3).
Full equivariance also requires that pooling, nonlinearities, and normalization respect the group: a standard ReLU applied to the components of a geometric vector breaks it (activations are therefore applied to invariant quantities such as the vector norm), and channel-wise batch normalization breaks the coupling between channels, which motivates equivariant normalization~\cite{meng2024geometric}.
The guarantee is only as strong as its weakest layer; one non-equivariant operation anywhere in the stack silently undoes the symmetry the rest preserve.

\paragraph{Equivariant Attention and Continuous SE(3) Symmetries.}
A more recent line lifts the treatment from convolutions to attention, with queries, keys, and values that transform predictably under a chosen Lie group.
The SE(3)-Transformer~\cite{fuchs2020se} is the canonical instance, building each attention layer from tensor-field representations so that rotating the input scene rotates the whole stack of outputs consistently; Tensor Field Networks~\cite{thomas2018tensorfield} use the same machinery without an attention head and have been especially influential in molecular modeling, where SE(3) equivariance is a problem constraint rather than an optional bias.
Both parameterize the kernel as a combination of spherical-harmonic basis functions, which is computationally heavy, so a parallel line pursues lighter equivariance: E($n$)-equivariant graph networks (EGNN)~\cite{satorras2021egnn} achieve E($n$)-equivariant message passing from relative distances and coordinate updates alone, and Vector Neurons~\cite{deng2021vector} lift scalar point-cloud features to 3D vectors, obtaining SO(3) equivariance as a drop-in primitive, each trading representational richness for scalability.
With the libraries above, trying equivariance on a new task now costs about as much engineering as a convolutional baseline; the choice of group still has to be made by the designer.


\subsubsection{Benchmarks}

Benchmarks in this area are constructed to expose behavior under symmetry transformations.
MNIST-rot is the standard 2D testbed, separating methods that learn rotation robustness through augmentation from those that build it in; ModelNet40~\cite{wu20153d} evaluates point-cloud classification under three protocols, aligned training and test (z/z), training with azimuthal rotations and testing under arbitrary SO(3) rotations (z/SO(3)), and arbitrary rotations in both (SO(3)/SO(3)), which is where augmentation, STN-style alignment, and PointNet's permutation-invariant pooling are separated from genuine rotation handling; ScanObjectNN~\cite{uy2019scanobjectnn} adds real scan noise and occlusion.
For equivariance as a hard constraint, QM9~\cite{ramakrishnan2014quantum} (small-molecule property prediction) and interatomic-potential tasks are the reference, where the SE(3)-equivariant models NequIP~\cite{batzner2022e3} and MACE~\cite{batatia2022mace} set the current points for data efficiency and accuracy.
Table~\ref{tab:p3_s7_2_3_symmetry} reports the two ladders on which the method families can be compared directly.

\begin{table}[!ht]
    \centering
    \caption{Symmetry benchmarks. MNIST-rot: test error (\%, lower is better); plain CNN and P4CNN figures from \citet{cohen2016group}, E(2)-steerable figure from \citet{weiler2019general}. ModelNet40: classification accuracy (\%, higher is better) under the train/test rotation protocols z/z, z/SO(3), and SO(3)/SO(3), all as re-evaluated by \citet{deng2021vector}.}
    \label{tab:p3_s7_2_3_symmetry}
    \fontsize{9}{11}\selectfont
    \setlength{\tabcolsep}{3.0mm}
    \begin{tabular}{L{0.30\textwidth} C{0.14\textwidth} C{0.11\textwidth} C{0.11\textwidth} C{0.13\textwidth}}
        \toprule
        \textbf{Method} & \textbf{MNIST-rot err.} & \textbf{z/z} & \textbf{z/SO(3)} & \textbf{SO(3)/SO(3)} \\
        \midrule
        Plain CNN (Z2CNN)~\cite{cohen2016group}                & 5.03 & --   & --   & --   \\
        P4CNN (G-CNN)~\cite{cohen2016group}                    & 2.28 & --   & --   & --   \\
        E(2)-steerable CNN~\cite{weiler2019general}            & 0.68 & --   & --   & --   \\
        PointNet~\cite{qi2017pointnet}                         & --   & 85.9 & 19.6 & 74.7 \\
        PointNet++~\cite{qi2017pointnetplus}                   & --   & 91.8 & 28.4 & 85.0 \\
        VN-PointNet~\cite{deng2021vector}                      & --   & 77.5 & 77.5 & 77.2 \\
        VN-DGCNN~\cite{deng2021vector}                         & --   & 89.5 & 89.5 & 90.2 \\
        \bottomrule
    \end{tabular}
\end{table}


\subsubsection{Challenges and Trends}
The most fundamental difficulty is choosing the symmetry group: invariance to a group that is too large (SO(3) when only SO(2) is task-relevant) destroys discriminative information, invariance to one too small leaves the model brittle to the symmetries the task actually exhibits, and the trade-off has no purely data-driven solution~\cite{cohen2016group,cohen2017steerable}.
Cost scales with the group: classical equivariant architectures index filters by group elements, so parameters and memory grow with the size of the chosen group, limiting how freely reflections, multi-scale rotations, or higher-order tensor representations can be included~\cite{weiler2019general}.
The assumption of \emph{exact} symmetry is also brittle: point clouds with self-occlusion, partial-view RGB, and noisy LiDAR scans do not realize an SO(3)-orbit, so an architecture that is mathematically equivariant can still fail when the input violates the assumed structure.
Table~\ref{tab:p3_s7_2_3_symmetry} shows where the two targets diverge: on 2D classification both hand-coded invariance and learned equivariance saturate, so data efficiency rather than accuracy becomes the discriminating axis, while on ModelNet40 the invariance-only PointNet family collapses under unseen rotations and the equivariant family holds its accuracy, because the task requires geometric information to be \emph{preserved}, not discarded.
Finally, hard equivariance sits uneasily with the attention-based representations that dominate the rest of the stack (§\ref{subsec:p3_s7_2_1}, §\ref{subsec:p3_s7_2_2}, vision--language pretraining): self-attention is permutation-equivariant but not rotation-equivariant, and bolting an equivariant layer onto an attention pipeline often degrades the inductive bias that motivated the choice.

Equivariant attention is the main compromise on offer, and frame averaging~\cite{puny2022frameaveraging} moves in the opposite direction, turning any off-the-shelf backbone into an invariant or equivariant one by symmetrizing its outputs over a small input-dependent set of group frames, which decouples the choice of symmetry from the choice of architecture and sidesteps the parameter cost of group-indexed filters.
These threads are pulling geometry-aware reasoning back toward the rest of the chapter: vision--language and vision--language--action models attach equivariant heads where 3D pose and geometric consistency matter, as in grasp policies and pose-conditioned generation, and primitives first devised for object classification are finding second lives in world models, embodied control, and 3D generation.




\subsection{Differentiable Optimization Layers for Spatial Reasoning}
\label{subsec:p3_s7_2_4}

Where the three preceding subsections put structure into the learned representation, differentiable optimization layers put it into the computation: selected feed-forward operations are replaced by the solution of a constrained optimization problem, so that hard logical, geometric, or physical constraints are enforced inside the forward pass while the whole system stays end-to-end differentiable~\cite{gould2022exploiting}.
In spatial reasoning this is how a network enforces distances, collision avoidance, or alignment, or solves a sub-problem such as pose estimation, within a learnable pipeline.
Figure~\ref{fig:p3_s7_2_4_differentiable_optimization_formulation} previews the representative tasks (top row) and the method families (bottom row).

\begin{figure}[!ht]
    \centering
    \includegraphics[width=\linewidth]{figures/part3/section7/differentiable_optimization_layers_for_spatial_reasoning.pdf}
    \caption{Differentiable optimization layers for spatial reasoning. Top, tasks: (1) QP relaxation of binary assignment for Sudoku~\cite{amos2017optnet,wang2019satnet}, (2) differentiable convex MPC~\cite{amos2018differentiable,agrawal2019differentiable}, and (3) bundle adjustment as differentiable nonlinear least squares for end-to-end SLAM~\cite{teed2021droid}. Bottom, methods: (1) OptNet~\cite{amos2017optnet}, (2) CVXPY Layers~\cite{agrawal2019differentiable,diamond2016cvxpy}, (3) Deep Declarative Networks and differentiable geometry solvers~\cite{gould2021deep,campbell2020solving}, and (4) SATNet-style differentiable SAT layers~\cite{wang2019satnet,shindo2024learning}.}
    \label{fig:p3_s7_2_4_differentiable_optimization_formulation}
\end{figure}


\subsubsection{Task Definition}
Given an input parameter vector $\theta$ produced by the preceding layers, a differentiable optimization layer computes
\begin{equation}
    z^\star = \operatorname*{arg\,min}_{z \in \mathcal{Z}} f(z; \theta) \;\; \text{s.t.}\;\; g(z;\theta) \leq 0, \; h(z;\theta) = 0,
\end{equation}
where $f$ is an objective function, $g,h$ denote inequality and equality constraints, and $\mathcal{Z}$ is a feasible domain.
The forward pass computes the optimal solution $z^\star(\theta)$, while the backward pass relies on implicit differentiation of the optimality conditions to obtain $\tfrac{\partial z^\star}{\partial \theta}$; the layer is declarative, in that the model specifies what to solve rather than how to compute it.
Under this view, QP solvers~\cite{amos2017optnet}, convex optimization layers~\cite{agrawal2019differentiable}, differentiable geometry solvers~\cite{teed2021droid}, and differentiable logic program solvers~\cite{wang2019satnet,shindo2024learning} are all instantiations of the same layer, tailored to different domains of spatial or symbolic reasoning.


\subsubsection{Methodology}

\paragraph{QP and Convex Optimization Layers.}
\citet{amos2017optnet} introduced OptNet, which treats the solution of a quadratic program as a layer: the forward pass solves
\begin{equation*}
    z^\star = \operatorname*{arg\,min}_z \tfrac{1}{2} z^{\top} Q z + q^{\top} z
\end{equation*}
subject to affine constraints whose matrices can depend on previous network outputs, and the backward pass differentiates the KKT optimality conditions exactly, with a custom GPU interior-point solver that handles many small QPs in parallel and reuses its factorization for the backward pass.
OptNet learned to solve Sudoku from data where conventional architectures could not, which showed that implicit layers can handle discrete or constrained reasoning problems.
\citet{agrawal2019differentiable} generalized the idea to any disciplined convex program: CVXPY Layers~\cite{diamond2016cvxpy} canonicalize the problem into conic form, solve it, and differentiate through the cone program's optimality conditions, so a constrained least-squares, a constrained logistic regression, or a robot-motion-limit constraint can be inserted as a layer without deriving gradients by hand.

\paragraph{Implicit Differentiation.}
All of these layers share one mechanism.
Rather than unrolling an iterative solver, one differentiates the closed-form optimality (KKT) conditions, which for a convex program comprise primal feasibility, dual feasibility, complementarity, and stationarity~\cite{DeepImplicitLayers}; differentiating this system and solving one linear system yields $\frac{\partial z^\star}{\partial \theta}$ exactly, without back-propagating through solver iterations.
Where OptNet hand-codes this backward pass, Deep Declarative Networks~\cite{gould2021deep} apply automatic differentiation to the KKT system itself.
The same principle extends to root-finding, fixed-point, and equilibrium layers via the implicit function theorem; deep equilibrium models~\cite{bai2019deep} treat an entire weight-tied network as a single fixed-point layer differentiated this way, and differentiable physics engines and planners propagate gradients through the equilibrium or optimum conditions of a physical system.

\paragraph{Differentiable Geometry Solvers.}
A subclass is tailored to geometric problems that are themselves argmin problems, typically minimizing a reprojection error: camera pose from 2D--3D correspondences, bundle adjustment, inverse kinematics.
\citet{campbell2020solving} introduced a differentiable blind PnP solver that estimates 3D pose as a layer; \citet{teed2021droid} built a differentiable bundle-adjustment module into a SLAM system, solving for camera poses and 3D points by iterative optimization while letting gradients flow; EPro-PnP~\cite{chen2022epro} recasts PnP as a probabilistic layer with a learnable distribution over poses; and Theseus~\cite{pineda2022theseus} provides a reusable library for differentiable nonlinear least squares (Gauss--Newton, Levenberg--Marquardt) that underlies bundle adjustment, pose-graph optimization, and inverse kinematics.
These layers keep the domain-specific optimizer but make it differentiable, so the whole vision or robotics system trains end-to-end while the geometry stays correct and constrained, for instance so that estimated object poses agree across views; \citet{gould2021deep} refer to such constructs as declarative nodes and demonstrate them on visual Sudoku and robust geometry estimation.

\paragraph{Differentiable Logic Layers.}
Beyond continuous optimization, the same construction extends to symbolic reasoning.
SATNet~\cite{wang2019satnet} integrates a differentiable satisfiability solver as a layer: the forward pass solves a relaxed SAT problem (a quadratic unconstrained binary optimization) to produce assignments consistent with logical constraints, and the backward pass differentiates implicitly through the solver's fixed-point conditions; \citet{shindo2024learning} extend differentiable logic programs to abstract visual reasoning benchmarks such as Raven's Progressive Matrices.
These layers belong to the same class of constraint-based implicit modules as the QP and convex layers, specialized to combinatorial and relational dependencies.


\subsubsection{Benchmarks}
Differentiable optimization layers have been validated across symbolic, geometric, and control settings.
Sudoku~\cite{amos2017optnet,wang2019satnet} and its visual variant~\cite{gould2022exploiting} test whether a layer models constraint satisfaction and propagates gradients through discrete reasoning; Raven's Progressive Matrices (RPM / I-RAVEN)~\cite{shindo2024learning} test whether differentiable logic captures relational structure in visual analogy puzzles; cart-pole, inverted-pendulum, and reaching tasks in MuJoCo and OpenAI Gym~\cite{agrawal2019differentiable} test whether differentiable MPC~\cite{amos2018differentiable} or convex layers improve control while enforcing physical feasibility; and maze navigation and obstacle avoidance test QP-based collision-avoidance layers by whether the agent reaches its target without violating safety constraints.
Evaluation combines task metrics with criteria specific to the embedded solver: geometric accuracy (translation / rotation error in pose estimation, trajectory or map error for a bundle-adjustment layer), task success rate (puzzles solved, feasible plans produced) together with the count or magnitude of violated constraints and, for physics-informed layers, residual energy or momentum drift, and solver cost and gradient stability (forward / backward runtime, iterations to converge, and sensitivity of training to approximate versus exact implicit gradients).
Table~\ref{tab:p3_s7_2_4_perf} lists the representative layers by solver type.

\begin{table}[t]
    \centering
    \small
    \caption{Representative differentiable-optimization layers, organized by solver / layer type (\cf\ Figure~\ref{fig:p3_s7_2_4_differentiable_optimization_formulation}).}
    \label{tab:p3_s7_2_4_perf}
     \fontsize{9}{11}\selectfont
    \setlength{\tabcolsep}{1.0mm}
    \begin{tabular}{L{0.26\textwidth} L{0.10\textwidth} L{0.28\textwidth} L{0.28\textwidth}}
        \toprule
        \textbf{Method} & \textbf{Year} & \textbf{Solver / layer type} & \textbf{Task / benchmark} \\
        \midrule
        OptNet~\cite{amos2017optnet}                 & 2017 & QP layer, KKT implicit diff.            & Sudoku from data            \\
        Differentiable MPC~\cite{amos2018differentiable} & 2018 & convex MPC layer                    & planning / control          \\
        CVXPY Layers~\cite{agrawal2019differentiable} & 2019 & disciplined convex / cone layer        & control, geometry           \\
        SATNet~\cite{wang2019satnet}                 & 2019 & differentiable MAXSAT relaxation       & visual Sudoku               \\
        Deep Declarative Nets~\cite{gould2021deep}   & 2021 & implicit-function-theorem nodes        & robust geometry, Sudoku     \\
        DROID-SLAM~\cite{teed2021droid}              & 2021 & differentiable bundle adjustment       & monocular / stereo SLAM     \\
        Theseus~\cite{pineda2022theseus}             & 2022 & differentiable nonlinear least squares & pose-graph / BA / IK        \\
        EPro-PnP~\cite{chen2022epro}                 & 2022 & probabilistic PnP layer                & monocular 6-DoF pose        \\
        \bottomrule
    \end{tabular}
\end{table}




\subsubsection{Challenges and Trends}
Solving an optimization problem inside a layer is more expensive than a convolution or a matrix product, and the forward pass becomes the bottleneck as problem size grows, which is exactly the regime of SLAM, bundle adjustment, and motion planning.
OptNet's batched GPU solver shows how problem structure buys tractability for many small problems, but global route planning or large-scale 3D reconstruction still overwhelms existing solvers, forcing approximate solves, warm starts, or hierarchical decomposition to keep each optimization manageable.

Gradient stability is the second difficulty.
If the solution is non-unique or depends discontinuously on the input, as at boundary optima, at active-set changes in a linear program, or in the presence of discrete variables, the implicit gradient is unstable or undefined; strict convexity guarantees a well-defined Jacobian, and real problems frequently violate it.
Entropy smoothing, small regularizers, and perturbation methods restore usable gradients in many cases, and in others the backward pass has to be designed by hand.

The third is physical consistency.
An optimization layer enforces only the constraints it is given: if the formulation omits a friction constraint in a contact model, the network exploits the gap and converges to implausible solutions, and because the layer sits inside a trainable system, the surrounding network can learn to feed it inputs that satisfy the constraints superficially while violating their intended meaning.
The benefits of interpretability and correctness therefore depend on the optimization model reflecting the real dynamics, which in practice means careful constraint modeling, additional penalties, or iterative refinement of the solver design; current work refines solvers, numerical stability, and integration with deep-learning frameworks toward layers that are both theoretically grounded and deployable in spatial systems.



\subsection{Reinforcement Learning with Spatial Inductive Biases}
\label{subsec:p3_s7_2_5}

Where the previous four subsections constrain the learned representation, reinforcement learning shifts the locus of structure to the \emph{action-conditioned} setting, in which an agent must reason about spatial relationships as a consequence of its own behavior over a long horizon.
Many spatial tasks, such as navigation, exploration, manipulation, and embodied interaction, cannot be reduced to static perception or single-step prediction; they require reasoning about how actions transform the agent's position, orientation, and relations to surrounding objects over time, under geometric constraints, partial observability, and delayed reward.
Generic RL policies built on unstructured function approximators struggle to generalize across unseen layouts, so a growing body of work introduces \emph{spatial inductive biases}, explicit planning iterations, map-like latent representations, differentiable search operators, or relational state abstractions, that constrain the hypothesis space of policies or value functions to reflect the compositional and geometric structure of spatial decision-making.


\subsubsection{Task Definition}
At each step $t$ the agent receives an observation $o_t$ from the environment $\mathcal{E}$ (RGB / RGB-D frames, egocentric maps, point clouds, language goals) together with its action history $a_{<t}$, chooses an action $a_t$ from a continuous control space or a discrete set $\mathcal{A}=\{1,\ldots,K\}$, and receives a scalar reward $r_t$.
The policy or value function is parameterized with a spatial structural prior $S$ (a differentiable planning module, a learned map or belief representation, a relational or graph encoder, an equivariant policy head),
\begin{equation}
  \pi_\theta(a_t \mid o_t, a_{<t};\, S),\qquad
  \max_\theta\ \mathbb{E}\!\left[\sum_{t=0}^{\infty} \gamma^t r_t\right],
\end{equation}
so that $S$ shapes \emph{which} policies and value functions the agent considers in the first place, not only how it generalizes after the fact; the optimization itself remains a standard RL objective (Bellman, policy gradient, actor--critic, or preference-based).

\begin{figure}[!ht]
    \centering
    \includegraphics[width=0.7\linewidth]{figures/part3/section7/rl_overview.pdf}
    \caption{The general framework in RL: the agent acts on the environment, receives an observation and a reward, and updates its policy from this feedback.}
    \label{fig:rl_overview}
\end{figure}


\subsubsection{Methodology}

Figure~\ref{fig:rl_overview} shows the loop every method below shares: the agent observes, acts, receives a reward and a new state, and updates its policy to maximize long-term return.
Learning a good policy requires balancing \textbf{exploration} of untried actions against \textbf{exploitation} of the actions currently believed best, the exploration--exploitation trade-off that underlies all of RL.
The subsections that follow walk through the optimization backbones in the order the field developed them, value-based, policy-based, actor--critic, PPO, and the recent preference-based and group-relative objectives, and note at each step where spatial structure enters.


\begin{figure}[!ht]
    \centering
    \includegraphics[width=0.70\linewidth]{figures/part3/section7/value-based-rl.pdf}
    \caption{Workflow of value-based RL.}
    \label{fig:value_based_rl}
\end{figure}

\paragraph{Value-Based RL.}
Value-based methods learn a value function that estimates the expected cumulative return of acting under a policy $\pi$.
The optimal state-action value function is
\begin{equation}
Q^*(s,a) = \max_\pi \mathbb{E}_\pi \left[ \sum_{t=0}^{\infty} \gamma^t r_t \,\middle|\, s_0 = s, a_0 = a \right],
\end{equation}
where $\gamma \in (0,1)$ is the discount factor, and it satisfies the \textit{Bellman equation}: if the optimal values $Q^*(s',a')$ of the next state are known for every action $a'$, the best action now is the one that maximizes $r+\gamma Q^*(s',a')$,
\begin{equation}
    Q^*(s, a) = \mathbb{E}_{s' \sim \mathcal{E}} [r+\gamma \max_{a'} Q^{*}(s', a') | s, a].
\end{equation}
Iterating this identity as an update, $Q_{i+1}(s, a) = \mathbb{E}_{s' \sim \mathcal{E} }[r+\gamma\; \max_{a'} \; Q_{i}(s', a') | s, a]$, converges to $Q^*$, but a table of values indexed by every state is useless in high-dimensional observation spaces, so a function approximator $Q(s,a;\theta)\approx Q^*(s,a)$ is learned instead, a neural network in the deep-RL setting (a Q-network).
Training minimizes a sequence of regression losses whose target is the bootstrapped Bellman estimate,
\begin{equation}
    L_i(\theta_i) = \mathbb{E}_{s, a \sim \rho(\cdot)} [(y_i - Q(s, a; \theta_i))^{2}],
\end{equation}
with $y_i = \mathbb{E}_{s' \sim \mathcal{E}} [r + \gamma \; \max_{a'} Q(s', a';\theta_{i-1}) | s, a]$ computed from the previous parameters $\theta_{i-1}$, held fixed, and $\rho(s,a)$ the distribution of visited state--action pairs; the greedy policy $\pi(s;\theta) = \arg\max_{a} Q(s, a;\theta)$ then follows from the learned values.

Deep Q-networks (DQN)~\cite{mnih2015human} made this scale to high-dimensional observations with experience replay and a slowly updated target network; their convergence is analyzed theoretically by \citet{fan2020theoreticalQlearning}, and Figure~\ref{fig:value_based_rl} shows the resulting workflow.
Beyond control, value-based RL drives \emph{active visual perception}: \citet{caicedo2015active} cast object localization as a sequential decision process in which a Q-network iteratively transforms an attention window until it tightly bounds the target, turning detection into a learned spatial search policy, a formulation since applied to active detection, visual search, and region selection for fine-grained recognition, where the agent learns \emph{where to look} rather than scanning exhaustively.
In spatial domains the limitation of vanilla DQN is that any planning behavior must be learned implicitly from data, which generalizes poorly to novel layouts, so several works embed planning structure in the value estimate itself.
Value Iteration Networks (VIN)~\cite{tamar2016valueiterationnetworks} unroll Bellman updates as convolutions over a spatial grid, a differentiable approximation of value iteration; Gated Path Planning Networks (GPPN)~\cite{lee2018gatedpath} replace the recurrence with gated units to improve optimization stability and long-range dependencies; and under partial observability QMDP-Net~\cite{karkus2017qmdpnet} integrates belief updates and approximate POMDP planning (QMDP)~\cite{littman1995learningQMDP} into one differentiable network.
In each case the learned value function is constrained to follow spatially structured dynamic programming rather than arbitrary function approximation.


\begin{figure}[!ht]
    \centering
    \includegraphics[width=0.7\linewidth]{figures/part3/section7/policy-based-rl.pdf}
    \caption{Workflow of policy-based RL.}
    \label{fig:policy_based_rl}
\end{figure}

\paragraph{Policy-Based RL.}
Policy-based methods~\cite{williams1992simple,zhang2021sampleefficientreinforcement} parameterize a policy $\pi_\theta(a|s)$ directly and maximize expected return,
\begin{equation}
J(\theta) = \mathbb{E}_{\pi_\theta} \left[ \sum_{t=0}^{\infty} \gamma^t r_t \right],
\end{equation}
by gradient ascent on
\begin{equation}
\nabla_\theta J(\theta) = \mathbb{E}_{\pi_\theta} \left[ \nabla_\theta \log \pi_\theta(a_t|s_t) \, R_t \right],
\end{equation}
where $R_t = \sum_{k=t}^{\infty} \gamma^{\,k-t} r_k$ is the return from time-step $t$ (Figure~\ref{fig:policy_based_rl}).
Optimizing action distributions directly suits spatial decision problems in which exploration and long-horizon credit assignment are central~\cite{chaplot2018active}, and policy gradients train embodied agents to actively gather visual evidence: Embodied Question Answering~\cite{das2018embodied} optimizes a navigation policy so that the agent moves to vantage points from which it can answer questions about a 3D scene.
The spatial competence, however, comes from the architecture rather than the objective; policy-gradient agents generalize across layouts when the policy carries a learned map~\cite{gupta2017cognitivemapping}, a grounded spatial representation~\cite{janner2018representationlearning}, memory, or a relational encoder~\cite{zambaldi2019relationaldeep}.


\begin{figure}[!ht]
    \centering
    \includegraphics[width=0.8\linewidth]{figures/part3/section7/AC-based-rl.pdf}
    \caption{Workflow of Actor-Critic-based RL.}
    \label{fig:actor_critic_based_rl}
\end{figure}

\paragraph{Actor-Critic RL.}
Actor--critic methods~\cite{degris2012modelfreerl} combine the two by jointly learning a policy $\pi_{\theta}(a|s)$ (the actor, which controls behavior) and a value function $V_\phi(s)$ (the critic, which scores the actions taken), as in Figure~\ref{fig:actor_critic_based_rl}.
The critic provides a low-variance estimate of the advantage,
\begin{equation}
A_{\phi}(s_t,a_t) \approx r_t + \gamma V_{\phi}(s_{t+1}) - V_{\phi}(s_t),
\end{equation}
or, when a Q-function is learned as well, $A_{\phi}(s_t,a_t) = Q_{\phi}(s_t,a_t) - V_{\phi}(s_t)$, and the advantage replaces the raw return $R_t$ in the policy-gradient update.
The asynchronous advantage actor-critic (A3C)~\cite{mnih2016asynchronous} made this stable and scalable and became a standard backbone for embodied navigation and spatial control: target-driven visual navigation~\cite{zhu2017target} conditions an A3C agent on an image of the goal and reaches semantically specified targets in unseen indoor scenes, and Neural SLAM~\cite{zhang2017neuralslam} gives an A3C agent a differentiable map memory.
The separation between actor and critic matches the structure of spatial tasks, with the actor responsible for geometry-aware decisions and the critic supplying global, long-horizon value signals, which is why actor--critic methods dominate navigation benchmarks that require memory, planning, and partial observability~\cite{mirowski2016learningnavigate}.



\begin{figure}[!ht]
    \centering
    \includegraphics[width=0.8\linewidth]{figures/part3/section7/ppo-based-rl.pdf}
    \caption{Workflow of PPO-based RL.}
    \label{fig:ppo_based_rl}
\end{figure}

\paragraph{PPO-Style Policy Optimization.}
Proximal Policy Optimization (PPO)~\cite{schulman2017proximalpo} is an actor--critic algorithm that stabilizes training by constraining each policy update through a clipped surrogate objective,
\begin{equation}
L^{\text{PPO}}(\theta) =
\mathbb{E} \left[
\min \left(
r_t(\theta) A_t,
\text{clip}(r_t(\theta), 1-\epsilon, 1+\epsilon) A_t
\right)
\right],
\end{equation}
where $r_t(\theta)$ is the probability ratio between the new and old policies (Figure~\ref{fig:ppo_based_rl}).
Its robustness and sample efficiency made PPO the de facto backbone for training agents in large-scale photorealistic 3D simulators such as Gibson~\cite{Xia2018Gibson} and Matterport3D~\cite{chang2017matterport3d}, and in manipulation suites such as panda-gym\footnote{\url{https://github.com/qgallouedec/panda-gym}}.
Scaled to 2.5 billion frames with DD-PPO~\cite{wijmans2019ddppo}, it essentially solves PointGoal navigation on Habitat (99.7\% success and 0.97 SPL on Gibson validation), and the same backbone underpins vision-and-language navigation~\cite{anderson2018vision}, object-goal navigation~\cite{chaplot2020object}, audio-visual navigation~\cite{chen2020soundspaces}, and planning over uncertain topological maps~\cite{beeching2020learninguncertaintopological}.

\begin{figure}[!ht]
    \centering
    \includegraphics[width=0.8\linewidth]{figures/part3/section7/dpo-based-rl.pdf}
    \caption{Workflow of DPO-based RL.}
    \label{fig:dpo_based_rl}
\end{figure}

\paragraph{Preference-Based Policy Optimization (DPO).}
Preference-based learning replaces the reward signal with comparisons.
Direct Preference Optimization (DPO)~\cite{rafailov2023directpo} learns a policy from preference pairs without an explicit reward model or online rollouts (Figure~\ref{fig:dpo_based_rl}): for two completions $(y_1,y_2)$ of an input $x$, a labeler marks the preferred one, $y_w \succ y_l \mid x$, and by expressing the implied reward of a Bradley--Terry preference model~\cite{bradley1952rankbt} through the optimal policy itself, DPO collapses the reward-learning-then-RL pipeline into a single maximum-likelihood objective on $\pi_\theta$:
\begin{equation}
    \mathcal{L}_{\text{DPO}}(\pi_{\theta};\pi_{\text{ref}}) = - \mathbb{E}_{(x,y_w, y_l)\sim \mathcal{D}} \big[ \log \sigma \big(\beta\log\frac{\pi_{\theta}(y_w|x)}{\pi_{\text{ref}}(y_w|x)} - \beta\log\frac{\pi_{\theta}(y_l|x)}{\pi_{\text{ref}}(y_l|x)}\big)\big],
\end{equation}
where $\pi_{\text{ref}}$ is a fixed reference policy and $\beta$ controls the deviation from it.
Developed for language-model alignment, DPO applies to spatial reasoning whenever supervision takes the form of trajectory preferences, plan rankings, or human judgments of spatial behavior~\cite{liang2025structuredpreference,huang2024making,qiu2025spatialPreference}, which arise naturally in navigation, manipulation, and interactive tasks where dense rewards are hard to specify.



\begin{figure}[!ht]
    \centering
    \includegraphics[width=0.8\linewidth]{figures/part3/section7/grpo-based-rl.pdf}
    \caption{Workflow of GRPO-based RL.}
    \label{fig:grpo_based_rl}
\end{figure}
\paragraph{Group-Based Policy Optimization (GRPO).}
Group Relative Policy Optimization (GRPO)~\cite{guo2025deepseekr1} removes the learned critic and replaces it with a group-relative baseline computed from multiple sampled trajectories (Figure~\ref{fig:grpo_based_rl}).
For each question $q$, it samples a group of outputs $\{o_1, o_2, \cdots, o_G\}$ from the old policy $\pi_{\theta_{\text{old}}}$ and optimizes $\pi_\theta$ by maximizing
\begin{equation}
\begin{split}
     \max_{\pi_\theta} &\;
    \mathbb{E}_{q \sim D_{\text{rl}}, \{o_i\}_{i=1}^{G} \sim \pi_{\theta_{\text{old}}}(O|q)}  \\
    & \Bigg[ \frac{1}{G}
        \sum_{i=1}^G
        \text{min}(\frac{\pi_\theta(o_i)}{\pi_{\theta_{\text{old}}}(o_i)} A_i, \text{clip}(\frac{\pi_\theta(o_i)}{\pi_{\theta_{\text{old}}}(o_i)}, 1-\epsilon, 1+\epsilon) A_i )
        - \beta \, \mathbb{D}_{\mathrm{KL}}\big(\pi_\theta \| \pi_{\mathrm{ref}}\big)
    \Bigg],
\end{split}
\end{equation}
where $\epsilon$ and $\beta$ are hyperparameters and the advantage $A_i$ is computed from the rewards $\mathbf{r}=\{r_1,\dots,r_G\}$ that a reward model assigns to the group, normalized within the group only:
\begin{equation}
    A_i = \frac{r_i - \text{mean}(\mathbf{r})}{\text{std}(\mathbf{r})}.
\end{equation}
Eliminating the value network reduces memory and stabilizes training.
Although explored first for language models, the formulation transfers to spatial decision-making wherever several candidate plans or trajectories can be scored against one another, and it has become the standard recipe for reinforcement fine-tuning of spatial vision--language models on visual reasoning, navigation, multi-step assembly, and video spatial reasoning~\cite{tan2025reasonrft,gao2026octonav,tang2025legoPuzzles,ouyang2025spacer}.


\paragraph{Architectural Spatial Inductive Biases.}
The backbones above are domain-agnostic; what makes a policy \emph{spatial} is the prior $S$ in its architecture, and three families recur across the examples already met.
\textbf{Map-based representations} give the agent an explicit allocentric memory: cognitive maps~\cite{gupta2017cognitivemapping} and Neural Map~\cite{parisotto2018neural} write observations into a spatially indexed latent grid, while Neural SLAM~\cite{zhang2017neuralslam} and Active Neural SLAM~\cite{chaplot2020learning} couple differentiable mapping with a hierarchical exploration policy.
\textbf{Differentiable planning modules} embed dynamic programming in the network, as in the VIN / GPPN / QMDP-Net line discussed under value-based RL, or learn a value-equivalent model and plan by lookahead search, as in MuZero~\cite{schrittwieser2020mastering}.
\textbf{Relational abstractions} replace flat feature maps with object- or graph-structured state, letting attention or message passing capture relations that transfer across layouts~\cite{zambaldi2019relationaldeep}.
In each case the RL objective is unchanged; the spatial competence comes from the prior, which is why these architectures generalize across geometries where unstructured policies overfit.


\subsubsection{Benchmarks}

Two complementary benchmark families shape the empirical landscape.
Embodied navigation and exploration benchmarks stress long-horizon planning under partial observability in photorealistic 3D scenes: Habitat~\cite{savva2019habitat} is the dominant platform, paired with real-world scans from Matterport3D~\cite{chang2017matterport3d} and Gibson~\cite{Xia2018Gibson} for large, scene-diverse PointGoal and ObjectGoal navigation that exercises geometry-aware exploration and cross-scene generalization; SoundSpaces~\cite{chen2020soundspaces} adds audio-visual cues as a multimodal test of spatial state estimation; and SOON~\cite{zhu2021soon} extends vision-and-language navigation to scenario-oriented object goals in unseen environments.
Symbolic and relational benchmarks test whether structured priors yield combinatorial and compositional generalization: PDDLGym\footnote{\url{https://github.com/tomsilver/pddlgym}} builds Gym environments from PDDL domains with relational observations and actions, and lightweight gridworld suites such as MiniGrid\footnote{\url{https://github.com/Farama-Foundation/Minigrid}} support controlled evaluation of planning, memory, and compositional generalization; the differentiable planners (VIN, GPPN, QMDP-Net) are evaluated on procedurally generated gridworld mazes and POMDPs of increasing size, where the question is whether the architectural prior generalizes to unseen layouts rather than memorizing training mazes.
Across both families the choice of RL algorithm matters less than the choice of prior: PPO has converged as the standard optimizer for embodied navigation, and differences between agents mostly reflect the map, planning module, or relational encoder on the policy, while on gridworld benchmarks the methods that embed planning structure consistently dominate unstructured DQN, with the largest gains where the task is most planning-shaped.


\subsubsection{Challenges and Trends}

Sample inefficiency comes first.
Spatial tasks involve long horizons, sparse rewards, and combinatorial state spaces, so na\"ive exploration is prohibitively expensive; architectural priors such as learned maps, planning modules, and relational representations improve efficiency~\cite{tamar2016valueiterationnetworks,zhang2021sampleefficientreinforcement}, but current RL systems still need orders of magnitude more interaction than humans or classical planners to acquire robust spatial competence.

Generalization and transfer come second.
Many agents perform well in their training environments and degrade sharply in unseen layouts, novel geometries, or altered dynamics, which suggests that much of the learned behavior rests on superficial correlations rather than abstract spatial structure; structured state representations, world models, and modular policy design~\cite{chaplot2018active,gupta2017cognitivemapping} help, but how to enforce spatial invariances and compositional generalization in RL in a principled way remains open.

Reward specification is the third.
Spatial tasks often lack dense, well-defined rewards, especially when success depends on long-term plan quality, safety, or efficiency rather than immediate outcomes.
Preference-based methods such as DPO~\cite{rafailov2023directpo} and group-relative objectives such as GRPO~\cite{guo2025deepseekr1} learn from comparative feedback instead, at the price of reward-model bias, preference inconsistency, and the difficulty of aligning learned preferences with the true task objective in complex spatial domains.

The directions that follow from these three are hybrid methods that integrate learning with explicit planning or search~\cite{tamar2016valueiterationnetworks,karkus2017qmdpnet}, structure-aware representations (object-centric, relational, graph-based) that scale spatial reasoning beyond grid-based and navigation-centric benchmarks~\cite{zambaldi2019relationaldeep}, and evaluation protocols that weight out-of-distribution generalization, robustness to environmental variation, and reasoning quality over performance on fixed benchmarks.




\subsection{Multi-agent Spatial Reasoning}
\label{subsec:p3_s7_2_6}


Where the previous five subsections study a single learning agent, multi-agent spatial reasoning (MASR) moves to the \emph{joint} regime: multiple embodied or abstract agents perceive, communicate, and act within a shared environment to accomplish cooperative or competitive goals, and the structural prior now governs decentralized decisions over a shared spatial state.
MASR links low-level sensorimotor streams (proprioception, occupancy grids, trajectories) to relational decision processes (joint planning, role assignment, coalition formation); its core problem is \emph{credit assignment and coordination} under spatial constraints, that is, how agents with partial, decentralized observations reason over nearby entities, infer latent intent, and synthesize actions that produce coherent team-level behavior.
Deep learning has moved the field from rule-driven coordination to end-to-end policies that integrate perception, communication, and control under partial observability and non-stationarity~\cite{lowe2017maddpg,foerster2018coma,rashid2018qmix}, with autonomous driving and interactive motion forecasting over road geometry (WOMD~\cite{ettinger2021womd}, nuScenes~\cite{caesar2020nuscenes}), multi-robot exploration, formation control, and patrol under attrition~\cite{goeckner2024magec}, and real-time strategy and team-sports simulation~\cite{samvelyan2019smac,kurach2020grf} as its representative domains.
We stress at the outset that the canonical algorithms surveyed below (VDN, QMIX, COMA, MADDPG) are themselves spatially agnostic: they solve a generic cooperative Dec-POMDP and see geometry only through the environment and the reward.
Explicit spatial structure enters chiefly through the benchmarks and through the graph-based family, where the interaction graph $\mathcal{G}_t$ makes neighborhood relations a first-class part of the policy.

\begin{figure}[t]
    \centering
    \includegraphics[width=0.95\linewidth]{figures/part3/section7/multi-agent_spatial_reasoning.pdf}
    \caption{Multi-agent spatial reasoning tasks (top) and mechanisms (bottom). Top: (a) cooperative micromanagement under partial observation (SMAC~\cite{samvelyan2019smac}), (b) interactive motion forecasting (WOMD, nuScenes~\cite{ettinger2021womd,caesar2020nuscenes}), and (c) multi-robot patrol after agent loss~\cite{goeckner2024magec}. Bottom: (d) centralized critic (MADDPG, COMA~\cite{lowe2017maddpg,foerster2018coma}), (e) value decomposition (VDN, QMIX~\cite{sunehag2018vdn,rashid2018qmix}), (f) graph coordination, and (g) attention across agents and time. Graph and sequence architectures serve prediction or control and do not by themselves imply decentralized execution. Red marks the lost agent; dashed trajectories are candidate futures and dashed graph links are intermittent interactions. Examples are schematic.}
    \label{fig:p3_s7_2_6_masr_formulation}
\end{figure}


\subsubsection{Task Definition}
We cast cooperative multi-agent spatial reasoning as a graph-structured Dec-POMDP
\begin{equation}
\mathcal{M}=(\mathcal{V},\{\mathcal{A}_i\},\mathcal{S},P,R,\{\mathcal{O}_i\},O,\rho_0,\gamma,T),
\end{equation}
with a (possibly time-varying) interaction graph $\mathcal{G}_t=(\mathcal{V},\mathcal{E}_t)$ derived from state (\eg, proximity or visibility) or estimated locally.
At time $t$, agent $i$ receives a partial egocentric observation $o_t^i\!\sim\!O(\cdot\mid s_t,i)$ (RGB-D, LiDAR, occupancy, kinematics) and optional neighbor messages $m_t^i$; we treat $\tilde{o}_t^i=(o_t^i,m_t^i,\mathcal{G}_t^i)$ as the input available to agent $i$, with any bandwidth cost included in $R$.
Agents act with decentralized, history-dependent policies
\begin{equation}
a_t^i \sim \pi_\theta^i(\cdot \mid h_t^i),\quad
h_t^i=\mathrm{upd}(h_{t-1}^i,\tilde{o}_t^i),
\end{equation}
the joint policy factorizes as $\pi_\theta(\mathbf a_t\mid \mathbf h_t)=\prod_i \pi_\theta^i(a_t^i\mid h_t^i)$, and the objective is
\begin{equation}
\max_{\theta}\ \mathbb{E}_{s_1\sim\rho_0,\ \pi_\theta,\ P}
\Big[\sum_{t=1}^{T}\gamma^{t-1}R(s_t,\mathbf a_t)\Big],\qquad \mathbf a_t=(a_t^1,\dots,a_t^n).
\end{equation}
In predictive variants, auxiliary heads output multi-agent forecasts or spatial fields with confidences for planning or losses; these are \emph{not} environment actions.
Settings differ along two further axes: online methods act in streaming settings under tight latency, whereas offline methods learn from logged interaction data (\eg, driving datasets) with batch updates and conservative behavior constraints; and centralized training with decentralized execution (CTDE) exploits global state to train joint value functions or critics while enforcing decentralized policies at test time~\cite{foerster2018coma,rashid2018qmix}, whereas fully decentralized learning never accesses global state, trading sample efficiency for deployability.


\subsubsection{Methodology}
Figure~\ref{fig:p3_s7_2_6_masr_formulation} summarizes the tasks above and the method families discussed below.

\paragraph{Classical Paradigms.}
Before learning, coordination relied on potential fields, formation keeping via graph Laplacians, and optimal multi-agent pathfinding with conflict-based search~\cite{Sharon2015CBS}, treated in §\ref{subsec:p3_s7_1_3}; these methods offer strong guarantees but struggle with perception noise, combinatorial scaling, and model mismatch in open worlds.

\paragraph{CTDE and Value Decomposition.}
Two lines made CTDE work.
Centralized critics resolve credit assignment and non-stationarity while keeping execution decentralized: COMA uses a counterfactual baseline computed by a centralized critic for decentralized actors~\cite{foerster2018coma}, and MADDPG extends deterministic policy gradients with agent-specific centralized critics~\cite{lowe2017maddpg}.
Value factorization instead decomposes the joint value: VDN sums per-agent utilities to approximate $Q_{\text{tot}}$ for cooperative tasks~\cite{sunehag2018vdn}, and QMIX generalizes the sum to a monotonic mixing network that preserves individual-global argmax consistency with more representational capacity~\cite{rashid2018qmix}.
These approaches made MARL scalable on SMAC~\cite{samvelyan2019smac}, where the value-decomposition line dominated the early cooperative leaderboard while the centralized-critic methods anchored other regimes (MADDPG on the particle environments, COMA on StarCraft micromanagement); the monotonic constraint and homogeneous-role assumptions limit expressivity, as the challenges below detail.

\paragraph{Graph Neural Networks for Coordination.}
Graph-based MARL encodes agents as nodes and interactions as edges, and message passing or attention over the neighborhood learns role-specialized coordination~\cite{jiang2018gcrl}; this family improves sample efficiency and robustness under topology changes and communication limits, and MAGEC~\cite{goeckner2024magec} shows resilience to agent attrition and communication disturbance in multi-robot patrol by combining GNN encoders with PPO-style multi-agent optimization.
Graph MARL targets a different axis from the leaderboard methods, robustness to topology disturbances rather than headline return, which exposes a gap between coordination-quality benchmarks (SMAC, MPE) and deployment-quality ones (multi-robot patrol, driving forecasting on WOMD and nuScenes) that standardized suites still under-represent.

\paragraph{Sequence Models and Transformers.}
Casting MARL as sequence modeling applies attention over multi-entity, multi-horizon contexts: the Multi-Agent Transformer~\cite{wen2022mat} maps the sequence of agents' observations to a sequence of actions with an encoder--decoder, using the multi-agent advantage decomposition to turn joint policy search into sequential decisions with linear complexity in the number of agents.
Transformer controllers, often hybridized with graph layers, address long-range temporal and relational dependencies, unify perception and control, and support offline RL through trajectory modeling; learned memory alleviates partial observability, while quadratic attention and credit assignment over long horizons remain open.


\subsubsection{Benchmarks}
Benchmarks span two regimes: simulation suites that stress cooperation, partial observability, and communication learning under controlled physics and reward shaping, and real-world multi-agent trajectory datasets that test interactive prediction and motion forecasting at deployment scale.
Table~\ref{tab:MASR_datasets} summarizes the most widely used resources.

\begin{table}[t]
\centering
\caption{Common MASR benchmarks spanning simulated coordination to real-world multi-agent trajectory datasets.}
\label{tab:MASR_datasets}
\fontsize{9}{11}\selectfont
\setlength{\tabcolsep}{3mm}
\begin{tabular}{@{}lccc@{}}
\toprule
\textbf{Dataset / Suite} & \textbf{Year} & \textbf{Scale} & \textbf{Action/Task Space} \\
\midrule
SMAC \cite{samvelyan2019smac} & 2019 &
  22+ micro scenarios & Discrete, decentralized \\
Multi-Agent Particle Env.\ (MPE) \cite{lowe2017maddpg} & 2017 &
  10+ 2D templates & Continuous/discrete \\
GRF \cite{kurach2020grf} & 2020 &
  Dozens of tasks / full-game & Continuous/discrete  \\
nuScenes \cite{caesar2020nuscenes} & 2020 &
  1,000 scenes $\times$ \SI{20}{s} @ \SI{2}{Hz} &
  Trajectory prediction \\
WOMD \cite{ettinger2021womd} & 2021 &
  $10^{5}$ scenes $\times$ \SI{20}{s} @ \SI{10}{Hz} &
  Trajectory prediction \\
Neural MMO \cite{mmo20aamas, MMO23nips} & 2020 &
  Persistent worlds; $O(10^2)$ agents & Long-horizon survival \\
\bottomrule
\end{tabular}
\end{table}

Evaluation depends on whether the setting is control or prediction.
Control is scored by mean episodic return $\mathbb{E}_{\pi}\!\left[\sum_t r_t\right]$, success or win rate, time-to-goal, collision and rule-violation rates, fairness or load-balancing indices for cooperative logistics, and communication overhead; prediction and planning are scored by spatial accuracy, from displacement errors over the forecast horizon (minADE, minFDE, miss rate on WOMD and nuScenes) to intersection-over-union between predicted and ground-truth occupied regions or trajectory tubes.


\subsubsection{Challenges and Trends}

Multi-agent spatial reasoning faces structural difficulties that distinguish it from the single-agent settings of §\ref{subsec:p3_s7_2_1}--§\ref{subsec:p3_s7_2_5}, organized here along four axes within the classical multi-agent paradigm.

\paragraph{Credit Assignment under Partial Observability and Long Horizons.}
The decentralized-execution constraint of the Dec-POMDP forces each agent to act on its own history $h^i_t$ alone, so relating a local action $a^i_t$ to a delayed team-level reward $r_T$ across a horizon $T$ produces an irreducibly noisy policy gradient.
Counterfactual baselines such as COMA \cite{foerster2018coma} marginalize over alternative actions of agent $i$ while holding the rest of $\mathbf{a}_t$ fixed and reduce gradient variance substantially, but they still rely on a centralized critic that, at scale, becomes a moving target for itself.
Value decomposition (VDN, QMIX) \cite{sunehag2018vdn,rashid2018qmix} side-steps the centralized critic by factorizing $Q_\mathrm{tot}$ across agents, yet the monotonicity constraint $\partial Q_\mathrm{tot}/\partial Q_i \ge 0$ that guarantees decentralized \emph{argmax} consistency cannot represent strictly non-monotonic cooperation games, a representational ceiling made explicit by the later non-monotonic factorizations QTRAN, weighted QMIX, and QPLEX, and visible on cooperative matrix games and SMAC's hardest scenarios~\cite{son2019qtran,rashid2020wqmix,wang2021qplex,samvelyan2019smac}.
Under sparse reward shaping, policies trained against these signals tend to chase short-term spatial cues (``move toward the closest enemy'', ``follow the nearest teammate'') rather than team-level objectives.

\paragraph{Scaling Laws of Joint Reasoning and Communication.}
Sequence-model multi-agent Transformers and dense graph message passing both incur $\mathcal{O}(n^2)$ complexity in the number of agents, plus additional quadratic terms across temporal contexts, which collides with the real-time latency budgets of multi-robot \cite{goeckner2024magec} and autonomous-driving deployments \cite{ettinger2021womd,caesar2020nuscenes}.
Sparse, learned interaction graphs are the natural mitigation, but they introduce their own credit-assignment difficulty: the edges $\mathcal{E}_t$ are themselves part of the learned policy, and gradients flowing through hard edge selection are biased or high-variance.
Communication bandwidth is a related gap: in the formulation above any message $m^i_t$ should carry a cost that appears in $R$, but the standard simulators (SMAC \cite{samvelyan2019smac}, MPE \cite{lowe2017maddpg}, GRF \cite{kurach2020grf}) idealize communication as free and noiseless, leaving little benchmark surface against which to push bandwidth-aware coordination.

\paragraph{Non-stationarity, Heterogeneous Roles, and Partner Generalization.}
Even in fully cooperative settings, intra-team policies co-evolve, so from the perspective of any single agent the environment is non-stationary \cite{lowe2017maddpg}, the failure mode the centralized-critic family of MADDPG / COMA \cite{lowe2017maddpg,foerster2018coma} was introduced to address.
The problem sharpens when teammates are heterogeneous, may drop out (the patrol-with-attrition setting of MAGEC \cite{goeckner2024magec}), or are drawn from an unseen partner distribution at evaluation time, the zero-shot coordination problem.
Persistent open-world simulators such as Neural MMO \cite{mmo20aamas,MMO23nips} push along this axis with procedurally generated maps and long-horizon survival, but standardized protocols for partner generalization and role specialization remain immature compared with the single-agent supervised benchmarks of §\ref{subsec:p3_s7_2_1} and §\ref{subsec:p3_s7_2_3}.

\paragraph{Ambiguous Spatial Semantics and Sim-to-Real Transfer.}
The egocentric observations $o^i_t$ that drive decentralized policies are partial and noisy: occluded teammates, uncertain affordances, and depth-discontinuity-induced map errors distort both prediction and control.
Driving benchmarks such as WOMD \cite{ettinger2021womd} and nuScenes \cite{caesar2020nuscenes} expose the prediction side of this gap, but the closed-loop control side (whether the same forecasts survive when an agent acts on its own predictions) is comparatively under-evaluated.
The sim-to-real gap on multi-robot systems is wider still: perception noise and actuation latency interact non-linearly with the coordination signal in ways that simulators rarely model, and resilience under realistic communication disturbance and agent attrition has only recently received first-class benchmark coverage \cite{goeckner2024magec}.

Three routes follow from these axes.
Topology-aware sparse interaction (learned edges, hierarchical groupings, or attention masks parameterized by $\mathcal{G}_t$) preserves long-range coordination while scaling sub-quadratically, with the GNN+PPO patrol architecture of MAGEC \cite{goeckner2024magec} as an early instantiation and graph encoders combined with sequence-model attention over short temporal windows as the natural next step.
The monotonicity ceiling of VDN / QMIX \cite{sunehag2018vdn,rashid2018qmix} can be relaxed by attention-weighted mixers, weighted projection schemes, or dueling-style decompositions that keep individual-global \emph{argmax} consistency while admitting non-monotonic cooperation, so that COMA-style counterfactual credit assignment \cite{foerster2018coma} no longer depends on a joint-action critic that becomes the bottleneck.
And explicit objectives for agent attrition and communication disturbance \cite{goeckner2024magec}, paired with failure-centric curricula and hooks for formal verification, move resilience from an emergent property of training to a first-class objective, decisive in safety-critical multi-robot and driving deployments \cite{ettinger2021womd,caesar2020nuscenes} where worst-case performance dominates the deployment decision; closed-loop evaluation suites that jointly report team return, safety, calibration, and latency are the missing measurement substrate, and the benchmarks for partner generalization, bandwidth-aware coordination, and sim-to-real transfer that would play the role SMAC and GRF played for CTDE are still under construction.





% \subsection{Summary}
% \label{subsec:p3_s7_2_summary}

% Read from top to bottom, the six paradigms move the structural prior from the representation to behavior: implicit reasoning leaves it latent, explicit and geometry-aware reasoning lifts it into a program or a symmetry group, optimization layers place it in the forward pass, and the two reinforcement-learning subsections let it shape the policy of one agent and then of a team.
% Table~\ref{tab:p3_s7_2_paradigm_comparison} summarizes the six along the axes used for the symbolic paradigms in §\ref{sec:p3_s7_1}.

% \begin{table}[t]
%     \centering
%     \fontsize{8}{10}\selectfont
%     \setlength{\tabcolsep}{2pt}
%     \caption{Comparison of the six paradigms of learning-based spatial reasoning covered in §\ref{sec:p3_s7_2}, on the same axes as the comparison table of §\ref{sec:p3_s7_1}; the \emph{locus} column gives where each paradigm injects its structural prior.}
%     \label{tab:p3_s7_2_paradigm_comparison}
%     \begin{tabularx}{\textwidth}{@{}p{2.9cm} p{3.3cm} p{3.3cm} X X@{}}
%         \toprule
%         \textbf{Paradigm} & \textbf{Structural prior (\emph{locus})} & \textbf{Canonical inputs / outputs} & \textbf{Strengths} & \textbf{Limitations} \\
%         \midrule
%         §\ref{subsec:p3_s7_2_1} Implicit neural reasoning      & Latent prior: attention, relation networks, FiLM modulate the learned feature space                & Pixel / feature inputs $\to$ answer, caption, or relation label                                  & End-to-end fluency on perception-heavy benchmarks; minimal modeling burden                & Limited inspectability; compositional generalization gaps; shortcut learning under sparse supervision \\
%         §\ref{subsec:p3_s7_2_2} Explicit structured reasoning  & Program- or graph-level prior: parsed layouts, scene graphs, neural-module / LLM-generated programs & Image $+$ compositional question $\to$ executable program $\to$ stepwise inspectable answer       & Compositional generalization; auditability; modular reuse of perception tools              & Parser / program-generator quality is a hard ceiling; brittle on long-tail vocabulary; supervision-hungry \\
%         §\ref{subsec:p3_s7_2_3} Geometry-aware reasoning       & Symmetry-group prior: invariance / equivariance to a chosen group of input transformations         & Image / point cloud / molecule $\to$ classification, pose, or property prediction                & Sample efficiency; out-of-distribution rotation / translation robustness                    & Group choice is a silent error mode; classical equivariant filters scale with $\lvert G \rvert$; brittle when inputs violate the assumed symmetry \\
%         §\ref{subsec:p3_s7_2_4} Differentiable optim.\ layers  & Optimization prior: a constrained convex / SAT problem is solved inside the forward pass            & Image / pose / observation $\to$ optimization-layer call $\to$ optimal solution $+$ KKT gradients & Certified feasibility / structure on the prediction; end-to-end training preserves geometry & Solver wall-clock dominates the forward pass; biased / high-variance gradients near constraint boundaries \\
%         §\ref{subsec:p3_s7_2_5} RL with spatial priors          & Action-conditioned prior: differentiable planning, learned maps, relational policies               & Observation $+$ action history $\to$ action under MDP / POMDP; reward signal                      & Long-horizon credit assignment; generalization across spatial layouts; embodied tasks       & Sample inefficiency vs.\ supervised paradigms; reward-shaping fragility; sim-to-real gap \\
%         §\ref{subsec:p3_s7_2_6} Multi-agent spatial reasoning   & Joint action-conditioned prior: Dec-POMDP, CTDE, graph-MARL, sequence-model coordination           & Per-agent egocentric obs.\ $+$ messages $+$ interaction graph $\to$ joint decentralized actions   & Team-level coordination; topology-aware robustness; resilience under attrition              & Credit assignment under partial observability $+$ non-stationarity; $\mathcal{O}(n^2)$ scaling; partner generalization immature \\
%         \bottomrule
%     \end{tabularx}
% \end{table}

% Every paradigm in the table trades the certificates of §\ref{sec:p3_s7_1} (feasibility proofs, force-closure verdicts, $do(\cdot)$-effect estimates) for end-to-end trainability, and each keeps a slice of the classical structure as its prior; which slice it keeps determines both its strengths and its blind spots.
% As the locus of that prior migrates from representation to behavior, the evaluation regime migrates with it, from offline benchmark accuracy at the top of the table to closed-loop team return at the bottom.
% What no single row supplies is \emph{cross-locus composition}: a geometry-aware backbone (§\ref{subsec:p3_s7_2_3}), an inspectable program-level reasoner (§\ref{subsec:p3_s7_2_2}), a differentiable feasibility layer (§\ref{subsec:p3_s7_2_4}), and an action-conditioned policy (§\ref{subsec:p3_s7_2_5}--§\ref{subsec:p3_s7_2_6}) stacked into one pipeline without losing the guarantee each contributes.
% That composition problem is the bridge to §\ref{sec:p3_s7_3}, where the prior is repositioned inside a large pre-trained model, and to §\ref{sec:p3_s7_4}, where it is distributed across a hybrid of classical and learned components.

% \newpage


\section{Foundation Model-driven Spatial Reasoning}
\label{sec:p3_s7_3}

\FramedInitial{F}{oundation} model-driven spatial reasoning turns to large pre-trained models (LLMs, vision-language models, vision-language-action models, and large vision models) as the engine of spatial inference. The motivation is empirical: such models display a striking facility with spatial concepts and relations even without task-specific fine-tuning, and a large body of work now recruits them as the core of systems for spatial question answering, planning, and explanation.

What separates these approaches is the locus and granularity at which the foundation model intervenes. At one end the model is the reasoning core itself, answering spatial queries directly (LLM/VLM-centric reasoning) or additionally emitting low-level actions (VLA-based reasoning). Its reasoning can be drawn out step by step, as intermediate deductions in language (chain-of-thought) or as an executable program synthesized from the task (code and DSL generation). It can be coupled to external machinery such as planners, renderers, simulators, and perception modules (tool-augmented reasoning), or grounded in external spatial memory (retrieval- and memory-augmented reasoning). And it can be handed control of a closed interaction loop with the environment (agentic reasoning) or paired with a learned neural simulator of dynamics that serves as the spatial substrate (world-model-based reasoning).


\subsection{LLM/VLM-centric Spatial Reasoning}
\label{subsec:p3_s7_3_1}
    % - Focus: the llm/vlm model for spatial reasoning
% [202511] Deconstructing Spatial Intelligence in Vision-Language Models
% [ICML2025] Why Is Spatial Reasoning Hard for VLMs? An Attention Mechanism Perspective on Focus Areas

Large language and vision-language models have become a default substrate for spatial reasoning, not because they compute geometry explicitly, but because large-scale pretraining leaves them with a usable, if imperfect, prior over how objects are arranged, oriented, and related. This subsection examines that prior directly: what spatial competence text-only LLMs acquire from language alone, how visual grounding sharpens it in VLMs, and where both still fall short of the metric, viewpoint-consistent reasoning that physical tasks demand. The methods are organized by the intervention used to elicit or strengthen this competence, from prompting to architectural change.

\begin{figure}[t]
    \centering
    \includegraphics[width=\linewidth]{figures/part3/section7/LLM_VLM-centric_spatial_reasoning.pdf}
    \caption{LLM/VLM-centric spatial reasoning, illustrated by the query ``Which object is left of the lamp?'': (i) text-only or visual inputs with optional spatial evidence, (ii) a pretrained LLM/VLM with four complementary, not sequential, intervention strategies (prompt, training, data, model), and (iii) the answer, other task outputs, and evaluation categories. Shading marks the queried region, not a measured attention map.}
    \label{fig:p3_s7_3_1_llm_vlm_formulation}
\end{figure}

\subsubsection{Task Definition}
Spatial reasoning, which embodied AI~\cite{driess2023palm, zhu2024spa}, autonomous driving~\cite{bogdoll2025muvo,zhou2024hktsg}, and spatial content generation~\cite{yu2025wonderworld, liu2025agentic} all depend on, is cast here as inference over learned representations rather than explicit geometric computation: a pre-trained model $\mathcal{M}_\theta$ answers a spatial query $q$ from context $c$,
\begin{equation}
  y=\mathcal{M}_\theta(q, c),\qquad c\in\{\varnothing,\ I,\ I+\text{aux}\},
\end{equation}
where the context may be absent (a text-only LLM), an image or image set $I$ (a VLM), or an image together with auxiliary evidence such as a scene graph, region features, or depth, and the output $y$ is a relation, a metric estimate, an orientation, or a navigation plan.
For text-only LLMs the spatial prior is built from statistical patterns in massive corpora~\cite{yamada2023evaluating,li2024advancing,wang2024picture}: such models infer that ``sky'' is above ``ground'', but the reasoning is \emph{ungrounded}, with no mapping to physical reality or visual signal, and therefore prone to physically impossible inferences.
VLMs ground the same prior by aligning the LLM with a visual encoder, and spatial competence then extends beyond recognizing objects and their intrinsic attributes to reasoning about relations among objects and to \emph{spatial extrapolation}, inferring unseen states or running mental simulations from what is perceived~\cite{liu2025deconstructing,chen2024spatialvlm}.




\subsubsection{Methodology}
% Training-free Prompting
% Model-Centric Architecture
% Explicit Spatial Injection
% Data-Centric Approaches

Interventions against the spatial deficiencies of foundation models fall into four levels, in increasing order of invasiveness: prompt engineering leaves the model untouched, training strategies update its weights, data-centric augmentation changes what it learns from, and model-centric enhancement alters the architecture.

\paragraph{Prompt Engineering.}
Prompting elicits latent spatial ability at inference time without weight updates.
Textual prompting applies chain-of-thought reasoning~\cite{liao2024reasoning} so that the model decomposes a spatial task into intermediate steps (identifying reference objects before deducing relative positions), and inserting scene graphs or coordinates extracted by expert models into the prompt raises accuracy further~\cite{mitra2024compositional}.
Visual prompting guides attention through markers superimposed on the image itself (masks, coordinate axes, keypoints), and in zero-shot settings such explicit visual cues can outperform fine-tuned models~\cite{yang2023set,liu20253daxisprompt}.

\paragraph{Training Strategies.}
Contrastive pretraining objectives induce a ``bag-of-concepts'' behavior that under-supervises spatial structure, so stronger supervisory signals are introduced to enforce spatial alignment.
Auxiliary losses such as image reconstruction or coordinate prediction compel the model to attend to pixel-level spatial detail rather than global semantics alone~\cite{wang2024reconstructive}.
Reinforcement learning with spatial reward functions (PPO or GRPO~\cite{wu2025reinforcing,shen2025fine}; view-consistency rewards~\cite{wang2025svqa}) trains models to perform mental visual simulation or to produce more precise spatial descriptions, moving spatial training from imitation toward preference-based optimization.

\paragraph{Data-Centric Augmentation.}
Because high-quality spatially annotated data is scarce, augmentation drives much of the progress.
Existing 2D/3D datasets are re-annotated with explicit spatial relations, metric distances, or orientations~\cite{wang2024all}, and simulation engines synthesize data with perfect geometric ground truth, which is especially effective for extrapolation tasks such as mental rotation and path planning where real-world annotation is impractical~\cite{tang2024sparkle,jiang2022pseudo}.

\paragraph{Model-Centric Enhancement.}
Architectural changes target the fine-grained visual-spatial features that mainstream encoders miss.
Because CLIP-style encoders~\cite{radford2021learning} prioritize global semantics over local geometry, higher-resolution encoders or fused features from multiple visual experts (SAM~\cite{kirillov2023sam, ravi2024sam2, carion2025sam}, DINOv2~\cite{dinov2}) sharpen the perception of shape and position; region-of-interest features or scene graphs supplied as additional input tokens decouple objects from their relationships and tighten fine-grained alignment~\cite{guo2024regiongpt}; and for 3D reasoning, models ingest point clouds, voxel grids, or depth maps directly~\cite{wangspatial} or lift 2D features into implicit 3D representations from multi-view images~\cite{zhu2024llava}.
A recent alternative internalizes 3D awareness through auxiliary supervision instead of extra inputs: SpatialSV~\cite{tang2026spatialsv} forces hidden visual features to predict explicit 3D targets (depth, ray maps, point clouds) during training, injecting geometric priors at no inference cost.



\subsubsection{Benchmarks}
% [EgoOrientBench]: Is 'right' right? enhancing object orientation understanding in multimodal large language models through egocentric instruction tuning

% [GeoMeter]: Understanding depth and height perception in large visual-language models
% [CV-Bench]: Cambrian-1: A fully open, vision-centric exploration of multimodal llms
% [What's Up]: What's" up" with vision-language models? investigating their struggle with spatial reasoning
% [SEED-Bench]: Seed-bench: Benchmarking multimodal llms with generative comprehension
% [SRBench]: Mind the gap: Benchmarking spatial reasoning in vision-language models
% [MINDCUBE]: Spatial mental modeling from limited views
% [RealWorldQA]: Grok-1.5v vision preview
% [OmniSpatial]: Omnispatial: Towards comprehensive spatial reasoning benchmark for vision language models

% \cite{jung2025right,azad2025understanding,tong2024cambrian,KamathHC23a,li2023seed,stogiannidis2025mind,yin2025spatial,jia2025omnispatial}

Evaluation has moved from general VQA to benchmarks that probe spatial capability along four levels~\cite{liu2025deconstructing}.
Intrinsic-attribute benchmarks test perceptual precision on single objects, with EgoOrientBench~\cite{jung2025right} evaluating object orientation from an egocentric viewpoint and GeoMeter~\cite{azad2025understanding} metric distance and dimension, and they show that even state-of-the-art models are weak at reference-free perception.
Relational benchmarks, the most populated class, test the parsing of spatial relations through vision-centric questions such as relative position (CV-Bench~\cite{tong2024cambrian}, What's Up~\cite{KamathHC23a}), and this is where general-purpose VLMs peak.
Deep-spatial-logic benchmarks go beyond visible relations: SRBench~\cite{stogiannidis2025mind} evaluates mental rotation and navigation, and MindCube~\cite{yin2025spatial} the construction of an internal 3D mental model from multiple views.
Comprehensive suites such as RealWorldQA~\cite{xai2024grok15v} and OmniSpatial~\cite{jia2025omnispatial} cover all levels at once; a wider coverage map of these static probes appears in Table~\ref{tab:spatial_reasoning_benchmark} of §\ref{subsec:p3_s7_3_7}.



\subsubsection{Challenges and Trends}
Three obstacles limit what foundation models can be made to do spatially.
Mainstream vision encoders such as CLIP optimize for global semantic alignment and capture coarse semantics at the expense of fine-grained geometric detail~\cite{wang2023makes,tong2024eyes}, and the contrastive objective compounds this by inducing the ``bag-of-concepts'' behavior that under-supervises structural relations among objects~\cite{zhou2025logic,wang2025spatialclip}.
Training data with precise metric annotation (quantitative distance, 3D coordinates) is scarce, and existing datasets rarely support the extrapolation tasks that require reasoning beyond direct observation~\cite{liu2023visual,agrawal2023stupd}.
Spatial reference is ambiguous: ``to the left'' may mean the viewer's left (egocentric), the object's left (allocentric), or another agent's left, and VLMs seldom switch frame of reference from context, which misaligns interpretation in embodied use~\cite{liu2023visual,lee2025perspective}.

Two directions respond to these.
Reinforcement learning with process rewards such as GRPO moves training from static VQA toward interaction, letting the model acquire spatial logic by trial and error in simulation or the real world~\cite{wu2025reinforcing}.
World-model-informed pretraining attacks the 2D-to-3D gap directly, learning physical regularities (gravity, collision) alongside image-text correspondence from synthetic corpora with exact geometric ground truth~\cite{yang2026mindjourney,tang2024sparkle,chen2024spatialvlm}; the use of such world models at inference time is the subject of §\ref{subsec:p3_s7_3_8}.




\subsection{VLA-based Spatial Reasoning}
\label{subsec:p3_s7_3_2}
    % - Vision-language-action model
% [Arxiv202509] Pure Vision-Language-Action (VLA) Models A Comprehensive Survey
% [Arxiv202505] Vision-Language-Action Models Concepts, Progress, Applications and Challenges

Vision-language-action (VLA) models push spatial reasoning past description and into execution: the same latent competence that lets a model answer where an object sits must now drive an end effector toward it without violating depth, contact, or stability constraints. Spatial reasoning here is therefore inseparable from the agent's body and the physics of its workspace, and a plan that reads correctly on paper still fails if it ignores reachability or collision. The methods below are organized by their action-generation paradigm (autoregressive, diffusion-based, reinforcement-learning-driven, and hybrid).

\begin{figure}[t]
    \centering
    \includegraphics[width=\linewidth]{figures/part3/section7/VLA_spatial_reasoning.pdf}
    \caption{VLA-based spatial reasoning as a perception--action loop: (i) embodied inputs, (ii) a VLA policy that generates actions autoregressively, by diffusion (denoising $A_K$ to $A_0$), or with a hybrid, plus optional RL refinement, and (iii) an action representation executed by a controller, with environment feedback closing the loop. $\Delta p,\Delta R, g, q^{*}$ are end-effector translation and rotation increments, gripper command, and joint targets. The robot scenes illustrate the initial state and intended outcome, not experimental observations.}
    \label{fig:p3_s7_3_2_vla_formulation}
\end{figure}

\subsubsection{Task Definition}

A VLA model extends the reasoning core of §\ref{subsec:p3_s7_3_1} with an action head: given a visual observation $o_t$, a language instruction $\ell$, and the robot's proprioceptive state $s_t$, it emits an executable action chunk
\begin{equation}
  a_{t:t+H}\sim\pi_\theta(\cdot\mid o_t,\ \ell,\ s_t),
\end{equation}
where each action encodes end-effector increments, gripper commands, joint targets, or pose waypoints that a controller converts into motion~\cite{zhang2025pure,sapkota2025vision}.
Spatial reasoning under this paradigm is \emph{embodied grounding}: language tokens must be anchored not only to visual features but to the three-dimensional geometry, affordances, and constraints of the workspace, so that the generated actions respect depth, stability, and contact rather than merely naming where objects are.



% 模型在具身环境中，跨越视觉感知和自然语言指令，理解、定位并操纵物理实体的能力 。

\subsubsection{Methodology}
VLA methods for spatial reasoning fall into four classes by their action-generation and reasoning paradigm.

% 自回归
\paragraph{Autoregression-Based Methods.}
Autoregressive methods cast spatial manipulation as next-token prediction over a unified token stream of discrete action tokens.
Gato~\cite{reed2022generalist} showed that heterogeneous modalities can be jointly tokenized and trained in one autoregressive model; PointVLA~\cite{li2025pointvla} injects point-cloud representations into a pretrained VLA for lightweight 3D reasoning without redesigning the architecture; and TraceVLA~\cite{zheng2024tracevla} adds visual trace prompting, explicit spatiotemporal cues that improve long-horizon spatial awareness and task consistency during generation.

% 扩散
\paragraph{Diffusion-Based Methods.}
Diffusion methods model action generation as conditional denoising, which suits trajectories that must respect physical constraints and uncertainty.
SE(3)-DiffusionFields~\cite{urain2022se} runs the diffusion process in $\mathrm{SE}(3)$ pose space and jointly optimizes grasping and motion planning; 3D Diffuser Actor~\cite{ke20243d} conditions the denoiser on structured 3D scene representations so that manipulation trajectories stay geometrically consistent; and ForceVLA~\cite{yu2025forcevla} brings force and tactile feedback in through a mixture-of-experts design for contact-rich manipulation.

% 强化学习
\paragraph{Reinforcement Learning-Driven Methods.}
Reinforcement learning refines VLA policies through interaction with the environment, particularly on long-horizon tasks (the broader treatment of RL with spatial inductive biases is §\ref{subsec:p3_s7_2_5}).
VIP~\cite{ma2022vip} learns a goal-conditioned, action-independent value function that implicitly encodes spatial distance to assess task progress; NaVILA~\cite{cheng2024navila} couples VLA representations with RL policies so that quadruped robots navigate and avoid dynamic obstacles on difficult terrain; and SafeVLA~\cite{zhang2025safevla} adds safety critics that constrain high-risk behavior during manipulation.

% 混合(自回归+扩散)

\paragraph{Hybrid Architectures.}
Hybrids combine generative mechanisms and add explicit spatial representations to strengthen 3D reasoning.
3D-VLA~\cite{zhen20243d} joins autoregressive language modeling with diffusion-based action prediction inside a generative 3D world model; VoxPoser~\cite{huang2023voxposer} has an LLM compose 3D voxel value maps, separating goal interpretation from geometric planning; ReKep~\cite{huang2024rekep} models spatiotemporal dependencies as relational keypoint constraints for precise manipulation; and RoboPoint~\cite{yuan2024robopoint} predicts affordance maps that encode where actions are feasible, as a prior for downstream planning and control.


\subsubsection{Benchmarks}
Evaluation rests on two complementary resources: generalization-oriented benchmarks in simulation and large-scale real-world embodied datasets.
VIMA-Bench~\cite{jiang2023vima} measures spatial generalization through a four-level protocol that runs from random object placement to fully novel task compositions; ALFRED~\cite{shridhar2020alfred} poses long-horizon household tasks that combine high-level goals with low-level instructions; and CALVIN~\cite{mees2022calvin} and LOTA-Bench~\cite{choi2024lota} test language-conditioned manipulation across environments, scoring generalization to unseen objects and configurations.
Among datasets, Open X-Embodiment~\cite{o2024open} is the most comprehensive, integrating 22 robot datasets from 21 institutions and 527 skills to enable cross-platform generalization; BridgeData~\cite{ebert2021bridge} spans 71 tasks in 10 environments and shows that joint training on limited out-of-domain data raises success rates substantially; and RH20T~\cite{fang2024rh20t} adds vision, force, and audio across 147 tasks, suited to one-shot learning and experience-based policy transfer.
Across these suites performance falls sharply along every generalization axis (L1 to L4 in VIMA-Bench, seen to unseen splits in ALFRED, single to multi-step chains in CALVIN), while data scale and embodiment diversity are the most reliable lever for cross-platform transfer and multimodal sensing (force, audio) becomes decisive where vision alone underdetermines the contact configuration.



\subsubsection{Challenges and Trends}
The generalization drop above reflects a \textbf{semantic--physical misalignment}: VLA systems inherit their competence from LLMs and VLMs that excel at semantic abstraction and cross-modal alignment but have had little exposure to physical interaction, so instructions are often interpreted correctly and executed incorrectly.
The same gap appears as \textbf{pseudo-interaction}: actions are generated from statistical co-occurrence in the training data rather than from an understanding of environment dynamics or causality, so when the environment is unfamiliar or changes quickly the model cannot revise its strategy from sensory feedback.

Two further limits are practical.
\textbf{Long-horizon consistency} is unsolved for autoregressive and diffusion policies alike: over extended action sequences the spatial representation drifts or hallucinates, and existing embodied datasets, including Open X-Embodiment, are dominated by tabletop manipulation and grasping, so environmental diversity and large unstructured scenes are under-represented.
\textbf{Inference latency} compounds across autoregressive decoding of large Transformers, and for high-frequency tasks such as dynamic grasping or mobile navigation the compute and memory demands remain prohibitive on embedded platforms.

The directions that follow are a unified token stream over vision, language, and action that lets VLA systems grow into proto world models (the world-model view is developed in §\ref{subsec:p3_s7_3_8}); active perception and causal reasoning, in which the agent probes its environment and updates its policy from feedback, as the answer to pseudo-interaction; and virtual-real data synthesis (high-fidelity simulation, synthetic generation, multi-robot data sharing) as the route to embodied data at the scale that text corpora provided for language models.


\subsection{Chain-of-Thought for Spatial Reasoning}
\label{subsec:p3_s7_3_3}
\label{ssec:cot-reasoning}

Spatial reasoning poses unique challenges for large language models (LLMs), as it requires the manipulation of geometric relations, perspective-taking, and dynamic transformations: skills that are not purely symbolic but grounded in embodied cognition.
Chain-of-Thought (CoT) prompting~\cite{wei2022chain}, originally developed as a mechanism to elicit intermediate reasoning steps in arithmetic and commonsense tasks, has therefore been imported into spatial reasoning not as a generic recipe for longer answers, but as a way to externalize the intermediate spatial commitments that a model must keep consistent: which object is left of which, which frame of reference is active, which action changes the scene, and which hypothesized 3D configuration explains the observation.
The evidence is mixed.
In some settings, explicit rationales help decompose a spatial query into auditable subgoals; in others, they induce a fluent but ungrounded verbal trace that makes spatial performance worse~\cite{kancheti2026cotdegrades}.

\begin{figure}[t]
    \centering
    \includegraphics[width=\linewidth]{figures/part3/section7/chain-of-thought_spatial_reasoning.pdf}
    \caption{Spatial reasoning tasks (top) and intermediate-trace structures (bottom). Top: (a) textual deduction in a north-up frame, (b) visual counting, where the boxed red cube is the only one left of the green sphere, and (c) embodied planning revised by new observations. Bottom: CoT, ToT, and GoT differ in trace topology; PoT and CoA ground the trace in program execution and environment interaction. $\tau=(t_1,\ldots,t_L)$ is the intermediate trace and $\hat y=g(c,q,\tau)$ the final answer. Scenes and trace diagrams are illustrative, not recorded model computation.}
    \label{fig:p3_s7_3_3_cot_formulation}
\end{figure}


\subsubsection{Task Definition}
The task is to answer a spatial query $q$ about a scene or environment, given context $c$ that may be purely textual, visual, 3D, temporal, or embodied.
A CoT-style reasoner augments the answer with an intermediate trace
\begin{equation}
    \tau = (t_1,\ldots,t_L), \qquad \hat{y}=g(c,q,\tau),
\end{equation}
where each $t_\ell$ records a partial spatial commitment: a relation, coordinate estimate, subgoal, viewpoint change, object binding, or predicted action effect.
The trace is useful only insofar as it improves the grounded answer $\hat{y}$ or makes the failure mode inspectable; many apparent ``reasoning steps'' are not grounded in metric or perceptual evidence and merely rationalize a final guess, so the central question is not whether a model can produce more text, but whether the trace helps preserve spatial invariants across transformations, viewpoints, time steps, and actions.

\subsubsection{Methodology}
CoT, Tree-of-Thought (ToT)~\cite{yao2023tot}, and Graph-of-Thought (GoT)~\cite{besta2024got} differ mainly in the topology of the intermediate trace: CoT is a single path through partial spatial commitments; ToT branches over alternative hypotheses or subplans and uses a value function to prune; GoT allows shared sub-derivations, such as reusing a computed room layout or object graph across several queried relations.
All three search over partial reasoning states with a proposal mechanism, an evaluator $V$, and a halting rule $H$,
\begin{equation}
    \tau^\star=\arg\max_{\tau} V(\tau)\quad\text{s.t.}\quad t_{\ell+1}\in\mathrm{Propose}(t_{\le \ell}),\ \ H(\tau)=1,
\end{equation}
and differ in whether the search structure is a path, a tree of alternatives, or a graph of reusable sub-results, and in whether the trace stays verbal or becomes executable through programs (Program-of-Thought, PoT~\cite{chen2022pot}) or environment-changing actions (Chain-of-Action, CoA~\cite{pan2024coa}).
The use cases below organize how these traces are deployed.

\paragraph{Text-Only Spatial Deduction.}
In linguistic spatial puzzles, route descriptions, and relative-position QA, the model must maintain a latent arrangement without visual feedback.
CoT can expose intermediate bindings (\eg, ``A is north of B; C is west of A'') and reduce one-shot relational confusion~\cite{hu2023chain}, while ToT-style branching can preserve competing layouts until enough constraints have been applied.
The persistent failure mode is coordinate drift: even with explicit steps, LLMs lose track of the origin, orientation, or accumulated displacement in textual mazes and grid-navigation tasks~\cite{zhang2025mitigating}, so text-only CoT is most reliable when the trace is tied to a compact symbolic artifact (a table, coordinate sketch, or adjacency list) rather than to unconstrained prose.

\paragraph{Multimodal Scene Reasoning.}
In image- or video-grounded QA, a CoT trace must bind language to perceptual evidence: which object instance is referred to, what relation is visible, which is occluded, and which viewpoint assumption is made.
Free-form rationales can verbalize a relation that the image does not support, so recent work moves toward grounded intermediate artifacts (structured layouts, object bindings, diagrams, or executable module calls): SCoT-style planners use layout-coordinate traces to guide diffusion models on the generation side~\cite{chen2026spatialcotgeneration}, and program-mediated systems call perception modules rather than describe them~\cite{hsu2025makes}, the paradigm of §\ref{subsec:p3_s7_3_4}.
Here ToT and GoT serve less as generic search algorithms than as hypothesis-management devices that let a model compare object correspondences, camera frames, or depth-order explanations before committing to an answer.

\paragraph{Embodied Planning and VLA Reasoning.}
In navigation, manipulation, and vision-language-action tasks, spatial reasoning unfolds over time, and the trace must connect perceptual state, subgoal choice, action feasibility, and feedback from the environment.
CoT-style subgoals support long-horizon planning in navigation~\cite{lin2025navcot,liu2025navr1}, embodied instruction following~\cite{lu2023thinkbot,mu2023embodiedgpt}, and VLA models that expose intermediate visual or action commitments: CoT-VLA pairs reasoning with predicted visual states~\cite{zhao2025cotvla}, ACoT-VLA formulates the trace in action space through coarse action intents and reference trajectories~\cite{zhong2026acotvla}, and scene graphs stabilize the trace by giving the model a persistent substrate for object identities and relations~\cite{zhang2025embodiedvsr}.

\subsubsection{Benchmarks}
Evaluation mirrors the three use cases.
Spatial puzzles, relative-position QA, and textual maze navigation are evaluated with prompt-based CoT on the suites of~\cite{hu2023chain,zhang2025mitigating}; visual question answering and scene-grounded compositional queries are evaluated under CoT with program execution (§\ref{subsec:p3_s7_3_4}), where the rationale's correctness is mediated by program output rather than by the rationale string; and long-horizon planning with sensory feedback is evaluated through CoT-augmented navigation (NavCoT~\cite{lin2025navcot}, NavR1~\cite{liu2025navr1}), CoT-augmented instruction following (ThinkBot~\cite{lu2023thinkbot}, EmbodiedGPT~\cite{mu2023embodiedgpt}), visual- and action-trace supervision for VLA models~\cite{zhao2025cotvla,zhong2026acotvla}, and scene-graph-augmented embodied reasoning~\cite{zhang2025embodiedvsr}.


\subsubsection{Challenges and Trends}
The transfer of CoT to spatial intelligence is uneven.
\textbf{Coordinate hallucination} persists under simple textual CoT: even with ToT-style beam expansion, prompt-based planners mistrack absolute positions across multi-step chains in maze and grid-navigation problems~\cite{zhang2025mitigating}.
\textbf{Ungrounded supervision} can make things worse: textual CoT applied without a structured or rationale-grounded training signal degrades spatial VLM performance, because the model is forced to verbalize a reasoning trace it does not possess~\cite{ji2025enhancing,kancheti2026cotdegrades}, especially when the prompt rewards plausible explanation length over geometric correctness.
\textbf{Search cost} scales super-linearly in solution depth for ToT and GoT (and in the branching factor for ToT), which limits real-time or on-device deployment.
\textbf{Trace--action mismatch} arises in embodied settings when a VLA model produces a coherent verbal rationale while its action head follows a different latent policy; action-aware CoT aligns rationales with action tokens and visual predictions~\cite{zhong2026acotvla}, but reliable attribution between the visible trace and the executed behavior remains unresolved.

\emph{Grounded rationale supervision} pairs CoT with explicit spatial signals (scene graphs~\cite{zhang2025embodiedvsr}, depth maps, 3D coordinate frames, action traces, or layout-coordinate traces for spatial generation~\cite{chen2026spatialcotgeneration}) so the rationale is verifiable rather than free-form.
A complementary line makes the trace \emph{execution-grounded}: PoT-style traces compile and run against a geometric engine or vision--language tool, with execution outputs serving as a truth signal, the natural bridge from CoT to code / DSL generation in §\ref{subsec:p3_s7_3_4}; for embodied agents the parallel bridge runs from CoT to CoA and the agentic orchestration loops of §\ref{subsec:p3_s7_3_7}.


\subsection{Code/DSL Generation for Spatial Reasoning}
\label{subsec:p3_s7_3_4}
Code/DSL generation treats spatial reasoning as program synthesis: an LLM or MLLM translates a natural-language task into an executable program, and a more deterministic system (a robot control API, a physics simulator, a library of vision functions) executes it.
Relative to the explicit structured reasoners of §\ref{subsec:p3_s7_2_2}, the specialized parser and fixed module set are replaced by the world knowledge and reasoning ability of the LLM, so the operation set is no longer confined to a predefined vocabulary; relative to the chain-of-thought traces of §\ref{subsec:p3_s7_3_3}, the trace is a program whose execution against an external environment supplies a truth signal that a free-form rationale cannot.
Figure~\ref{fig:p3_s7_3_4_code_formulation} shows the two operating modes.

\begin{figure}[t]
    \centering
    \includegraphics[width=\linewidth]{figures/part3/section7/code_DSL_spatial_reasoning.pdf}
    \caption{Code/DSL generation for spatial reasoning in two modes: understanding (top), where a counting query becomes a program over perception calls whose call trace is inspectable, and planning (bottom), where an instruction becomes robot API calls whose execution feedback drives program revision. Function names are illustrative pseudocode, not a specific system's API, and the scenes do not show experimental execution.}
    \label{fig:p3_s7_3_4_code_formulation}
\end{figure}


\subsubsection{Task Definition}
Given a high-level, often ambiguous, natural-language query $q$ and access to an API $\mathcal{A}$ (a robot's control interface, a programming-language interpreter, or a set of vision-language functions), the LLM acts as a high-level planner that emits a program and offloads execution to the external system $\mathcal{E}$,
\begin{equation}
  P=\mathrm{LLM}(q,\mathcal{A}),\qquad y=\mathrm{Exec}(P;\mathcal{E}),
\end{equation}
where $P$ is written in a general-purpose language such as Python or in a domain-specific language and composed of calls that manipulate objects, query spatial relations, or control an agent, and the output $y$ is a change in the environment's state (a robot moving an object), an answer to the query, or both.

\subsubsection{Methodology}

\paragraph{Code/DSL for Understanding.}
For static scenes the LLM generates code that dissects a visual query by composing calls to perception modules, so the program is itself an explicit, verifiable reasoning chain.
VisProg~\cite{gupta2023visual} prompts an LLM to produce a program in a high-level DSL that sequentially invokes modules for image understanding, image manipulation, and knowledge retrieval, implemented either heuristically in Python or by neural models such as OWL-ViT~\cite{minderer2022simple} and Stable Diffusion~\cite{rombach2022ldm}.
ViperGPT~\cite{suris2023vipergpt} instead exposes two classes, \texttt{ImagePatch} and \texttt{VideoSegment}, whose methods (\texttt{find}, \texttt{exists}, \texttt{compute\_depth}, \texttt{verify\_property}) each call a pretrained model, and generates a Python script that chains these calls to answer a complex query about an image or video, composing existing vision tools zero-shot.
Both depend on a static, human-defined DSL; VADAR~\cite{marsili2025visual} removes that limit by synthesizing the API implementations dynamically, in an API-generation stage followed by program synthesis, and LEFT~\cite{hsu2023s} makes the approach learnable by having the LLM generate differentiable logic programs that chain trainable perception networks, so concept grounding can be learned where no pretrained module exists.
The same recipe lifts to 3D: LLM-Grounder~\cite{yang2024llmgrounder} decomposes a grounding query into Python code that calls 3D grounding tools such as OpenScene~\cite{peng2023openscene} or LERF~\cite{kerr2023lerf} to locate objects in a scene.

\paragraph{Code/DSL for Planning.}
Beyond static understanding, the generated code can serve as a policy that acts in an environment, with the LLM as high-level planner and the interpreter handling deterministic low-level execution.
Code as Policies~\cite{liang2023code} has the LLM write Python that calls a robot's control API, so that ``put the apple in the bowl'' becomes \textit{robot.pick('apple')} followed by \textit{robot.place('bowl')}, with implicit steps and preconditions supplied by the model's common sense.
In simulation, Voyager~\cite{wang2023voyager} explores Minecraft by writing, executing, and refining its own code: it prompts the LLM for code to achieve a goal, receives environment feedback (success, failure, new inventory items), debugs accordingly, and adds successful routines to a growing library of reusable skills.


\subsubsection{Benchmarks}
Code/DSL generation for spatial reasoning is typically evaluated under three benchmark families that mirror the two-mode split of the Methodology above. For 2D scene understanding, methods are evaluated on compositional visual question answering (GQA~\cite{hudson2019gqa}), referring expression comprehension (RefCOCO / RefCOCO+ / RefCOCOg~\cite{yu2016modeling,mao2016generation,kazemzadeh2014referitgame}), and image-editing benchmarks. For 3D scene understanding, evaluation has converged on 3D referring expression and 3D visual grounding benchmarks (\eg,\ ScanRefer~\cite{chen2020scanrefer}, Multi3DRefer~\cite{zhang2023multi3drefer}, ScanQA~\cite{azuma2022scanqa}), which exercise both the perception-API composition and the long-tail of object-attribute reasoning. For embodied planning, the canonical evaluation suites are ALFWorld~\cite{shridhar2020alfworld} and VirtualHome~\cite{puig2018virtualhome} (long-horizon household tasks) and Minecraft (open-world goal acquisition, as exemplified by Voyager~\cite{wang2023voyager}); the headline metric in this family is task-success rate over multi-step goals rather than single-step accuracy. Across all three families, the recurring methodological observation is that code-generation methods inherit the accuracy of their underlying perception or control modules, so benchmark performance tracks the module suite more tightly than it tracks the LLM backbone.

\subsubsection{Challenges and Trends}
The interpretable trace comes with three recurring limits.
A \textbf{fixed, human-authored DSL} (ViperGPT's \texttt{find} / \texttt{compute\_depth} / \texttt{verify\_property}) restricts the model to a predefined operation set, so spatial concepts outside the API are systematically out of reach; VADAR's dynamic API synthesis~\cite{marsili2025visual} is a response but does not solve the coverage problem.
\textbf{Brittle execution with silent failure} follows from the dependence on perceptual modules (OWL-ViT~\cite{minderer2022simple}, OpenScene~\cite{peng2023openscene}, LERF~\cite{kerr2023lerf}): an upstream error propagates through the program and yields a plausible but wrong answer with no diagnostic signal.
\textbf{Sample-inefficient debug loops} limit self-refinement agents such as Voyager~\cite{wang2023voyager} and Code as Policies~\cite{liang2023code}, whose LLM calls per task scale poorly in high-dimensional or stochastic environments.

Two directions follow.
PoT-style execution-grounded rationales are being folded back into chain-of-thought training, with the generated program serving as both rationale and verification signal, and \emph{learnable DSLs} such as LEFT~\cite{hsu2023s} replace fixed perception with differentiable logic programs whose components are trainable, pointing toward a DSL that adapts under task supervision rather than being hand-engineered.


\subsection{Tool-augmented Reasoning}
\label{subsec:p3_s7_3_5}
\label{ssec:tool-augmented}
Tool-augmented reasoning in spatial intelligence is the paradigm in which an LLM or VLM acts not as an isolated solver but as an orchestrator over a catalogue of external operators.
These operators may be planners, renderers, CAD systems, physics simulators, retrieval engines, perception modules, or metric estimators.
The section boundary is deliberately operational: tools are treated as black boxes with input--output contracts, and the central research problem is the orchestration policy that decides when to call them, how to parameterize the call, and how much to trust the returned result.
When the executable program itself becomes the reasoning trace, the detailed treatment belongs to code / DSL generation (§\ref{subsec:p3_s7_3_4}); when persistent external memory becomes the learned object, it belongs to retrieval- and memory-augmented reasoning (§\ref{subsec:p3_s7_3_6}); when the loop becomes closed-loop control, it becomes agentic spatial reasoning (§\ref{subsec:p3_s7_3_7}); and when the tool is a learned predictive simulator, the specialized discussion is world-model-based reasoning (§\ref{subsec:p3_s7_3_8}).

\begin{figure}[t]
    \centering
    \includegraphics[width=\linewidth]{figures/part3/section7/tool_augmented_spatial_reasoning.pdf}
    \caption{Tool-augmented spatial reasoning as iterative orchestration over black-box tools: (i) scene, instruction, and tool catalogue $\mathcal{T}$, (ii) an LLM/VLM that selects tool $i_t$ with arguments $q_t$, receives result $r_t$, and appends to history $h_t$ until STOP, and (iii) the answer or plan with its tool-call trace. Execution, representation, validation, and memory are tool roles, not sequential stages. Solid and dashed boxes mark the selected cube and target compartment; the scene is a generated illustration with annotated, not detected, overlays.}
    \label{fig:p3_s7_3_5_tool_formulation}
\end{figure}

\subsubsection{Task Definition}
A tool-augmented reasoning system is formally defined as a tuple $(\mathcal{M}_\theta, \mathcal{T}, \pi_\phi)$, where $\mathcal{M}_\theta$ denotes the reasoning core, $\mathcal{T} = \{T_1, \dots, T_k\}$ is a catalogue of external tools, and each tool $T_i : \mathcal{I}_i \rightarrow \mathcal{O}_i$ is exposed only through a structured input--output interface.
The orchestration policy $\pi_\phi$ governs tool selection, argument construction, response parsing, retry / fallback behavior, and feedback integration.
Given an input $x \in \mathcal{X}$, the system maintains a time-ordered reasoning history $h_t$ (initially $h_0 = \varnothing$) and proceeds by the iterative update
\begin{equation}
\begin{aligned}
s_t &= \mathcal{M}_\theta(x, h_{t-1}), &
(i_t, q_t) &= \pi_\phi(s_t), \\
r_t &= T_{i_t}(q_t), &
h_t &= h_{t-1} \mathbin{\Vert} (s_t, i_t, q_t, r_t),
\end{aligned}
\end{equation}
where $s_t$ is the intermediate reasoning state, $(i_t, q_t)$ is the selected tool-and-query pair, $r_t$ is the returned result, and $\Vert$ denotes appending to the interaction history.
The loop halts when the policy or model emits a STOP decision, after which an output decoder $g$ produces the final answer $y = g(h_T)$.
Both the reasoning core and the tool internals are black-box; what is exposed, optimized, and evaluated is the interface history and the policy that controls it.

\subsubsection{Methodology}
Two design choices interact: the \emph{tool catalogue} determines which black-box capabilities are available, and the \emph{orchestration policy} determines how they are sequenced, checked, and folded back into the model's spatial answer.
The catalogue roles below are complementary, not sequential.

\paragraph{Execution-Augmented Reasoning.}
Execution tools bridge the semantic--geometric gap between high-level language and low-level feasible action~\cite{zhang2025scaling}: they expose callable primitives such as motion planners, robot skills, affordance estimators, navigation controllers, or simulator actions, and the orchestrator translates a spatial intention into a sequence of calls, monitors whether the result satisfies the goal, and recovers when a tool fails.
Representative systems include affordance-grounded planning, where language-model outputs are modulated by value functions for skill feasibility~\cite{ahn2022icanisay}; API- or code-mediated manipulation, where perception and control primitives are composed through generated calls~\cite{liang2023code,huang2023voxposer}; and action-token models that fold robot actions into the model vocabulary~\cite{brohan2022rt,zitkovich2023rt2,driess2023palm}.
SpatialClaw~\cite{cho2026spatialclaw} studies the action interface itself: a stateful Python kernel pre-loaded with perception and geometry primitives lets the agent execute one analysis step, observe intermediate outputs, and adapt the next call rather than committing to a full static program.
Closed-loop feedback and hierarchical decomposition~\cite{huang2022innermonologue,singh2023progprompt,yang2024vlmtamp,chen2024autotamp,bhat2024brainbodyllm} push this family toward agentic orchestration.

\paragraph{Representation-Augmented Reasoning.}
Representation tools externalize spatial working memory through sketches, 3D renders, CAD systems, scene graphs, game engines, or visualization buffers, motivated by the observation that transformers struggle to manipulate fine-grained geometry internally~\cite{hu2024visualsketchpad}.
Externalization reduces hallucination by moving spatial computation out of parametric memory~\cite{hu2024visualsketchpad} and enables iterative self-correction through visual feedback loops~\cite{wu2025reinforcing}; the claim is that spatial reasoning improves when models can \textit{see what they think}, as in Vision-Augmented Prompting~\cite{xiao2024enhancing}, which interleaves image synthesis with chain-of-thought so that textual and visual reasoning can be reconciled.
Recent spatial-tool agents make the evidence acquisition active: PERIA~\cite{li2026peria} pairs perception tools that expose symbolic or spatial evidence with interaction tools that manipulate the visual context, learning multi-tool behavior from synthesized trajectories and reinforcement learning, and S-Agent~\cite{dai2026sagent} casts the VLM as a semantic planner over a hierarchy of spatial tools that lift 2D object evidence into 3D geometric support while accumulating scene and agent memory across multi-view images or video.

\paragraph{Validation-Augmented Reasoning.}
Validation tools check whether intermediate outputs are logically, causally, geometrically, or physically consistent: the orchestrator translates a candidate answer or plan into a verifiable representation and queries a planner, solver, judge, or simulator.
% symbolic planners
Generated spatial plans can be encoded as constraint-satisfaction problems or PDDL specifications and validated by symbolic planners or SMT solvers~\cite{liu2023llmp,hao2023reasoning,yang2023coupling}, which checks preconditions, geometric relations, and causal effects;
% physics simulators as verifiers
physics simulators such as MuJoCo~\cite{todorov2012mujoco} and Isaac Sim\footnote{\url{https://developer.nvidia.com/isaac/sim}} act as physical oracles for the feasibility of proposed movements, collisions, and trajectories~\cite{liu2022mindseye,huang2023voxposer,chen2024spatialvlm};
% Neuro-symbolic methods
neuro-symbolic systems integrate Answer Set Programming solvers that translate natural language into formally verifiable logical rules for relational and topological coherence~\cite{yang2020neurasp,premsri2024neurosymbolic};
% plan-critic
and plan-critic mechanisms~\cite{lei2025stma,wang2025multiagent} or an LLM-as-a-Judge that evaluates and refines the plans of a primary planner close a self-correcting loop~\cite{hariharan2025plan}.


\paragraph{Memory-Augmented Reasoning.}
Memory tools expose external storage and retrieval through a black-box query interface: the orchestrator asks for relevant prior observations, spatial subgraphs, object histories, or task episodes and decides how to incorporate what is returned, which addresses the \textit{``crisis of context''} in long-horizon robotic tasks whose sensory streams exceed the finite context window of an LLM.
At this level the memory may be a hierarchical cognitive architecture such as RoboMemory or an open-vocabulary 3D scene graph such as ConceptGraphs (both treated in §\ref{subsec:p3_s7_3_6}), a classical 3D or dynamic scene graph~\cite{armeni20193dsg,rosinol20203ddsg}, or a system-level memory manager: MemOS~\cite{li2025memos} introduces paging and scheduling that swap spatial contexts in and out of the reasoning core, treating the LLM as a CPU and the tool as RAM, and spatio-temporal retrieval engines such as ReMEmbR~\cite{anwar2025remembr} and MemoTime~\cite{tan2025memotime} index video history in vector databases so that an agent can answer ``where did I leave the keys five minutes ago?''.

\paragraph{Predictive Tools and World-Model Feedback.}
Some catalogues include tools that simulate counterfactual or future spatial states rather than observe or verify the present one: generative video models, learned dynamics models, and predictive latent models are called as expensive but powerful instruments for ``what would happen if'' or ``what would this look like from another viewpoint''~\cite{liu2022mindseye,yang2025cambrians}.
Here they are one class of callable tool; their feedback forms, imagination cost, and fidelity are the subject of §\ref{subsec:p3_s7_3_8}.


\subsubsection{Benchmarks}
Evaluation asks whether the orchestrator selects useful tools and converts their outputs into better spatial decisions.
For execution-augmented tool use the canonical suites are SayCan~\cite{ahn2022icanisay} (LLM-grounded household manipulation), ALFWorld~\cite{shridhar2020alfworld} and VirtualHome~\cite{puig2018virtualhome} (long-horizon text-and-embodied environments), and Voyager-on-Minecraft~\cite{wang2023voyager} (open-world skill accretion through code-as-policy).
For spatial-reasoning-specific tool use, SpaceTools~\cite{chen2025spacetools} targets the metric gap between VLM perception and physical manipulation, and Cambrian-S~\cite{yang2025cambrians} introduces the spatial-supersensing framework that exercises tool augmentation from perception (Stage~I) to predictive world modeling (Stage~IV).
Validation- and memory-augmented reasoning are evaluated on the memory and agentic benchmarks of §\ref{subsec:p3_s7_3_6} and §\ref{subsec:p3_s7_3_7}, and predictive tools on those of §\ref{subsec:p3_s7_3_8}, where the central question is whether imagined evidence improves the final answer enough to justify its cost.


\subsubsection{Challenges and Trends}
The recurring limitations originate in the LLM--tool interface rather than in any individual tool.
The \textbf{exploration bottleneck} is combinatorial: each tool adds an action dimension and tool composition multiplies them, so naive reinforcement learning over tool use is inefficient; Double Interactive Reinforcement Learning (DIRL)~\cite{chen2025spacetools} warm-starts from expert demonstrations before open exploration, but the scaling of the action space itself is unresolved.
The \textbf{metric gap} moves rather than disappears: VLMs excel at semantic perception but fail at the precise metric estimation manipulation requires, and offloading geometry to depth estimators and pose solvers lets the LLM reason over measurements instead of hallucinations~\cite{chen2025spacetools,zhang2025scaling}, yet query formulation, output parsing, and error attribution at the interface become the new precision bottleneck.
\textbf{Interface calibration} compounds this: tool outputs arrive with heterogeneous units, confidence scores, coordinate frames, and failure modes, and an orchestrator that cannot judge whether a returned depth map, solver result, retrieved memory, or imagined rollout is trustworthy can perform worse than a direct model response.

\emph{Learned orchestration} is one direction: policies are moving from hand-written call rules toward imitation- or reinforcement-learned controllers that decide which tool is worth its latency and uncertainty budget~\cite{chen2025spacetools}.
A parallel shift is the \emph{System-1 / System-2 architectural split}, in which the neural backbone handles fast, intuitive perception (System~1) while the toolset handles slow, deliberative, verifiable reasoning (System~2)~\cite{yang2025neurosymbo}, achieving both flexibility and rigor.


\subsection{Retrieval- \& Memory-augmented Spatial Reasoning}
\label{subsec:p3_s7_3_6}
\label{ssec:memory_spatial_reasoning}

Retrieval- and memory-augmented spatial reasoning is best understood through a \emph{Planner-Memory-Retriever} decomposition. The Planner $\mathcal{P}$ decides what spatial information is needed for the current instruction; the Memory Bank $\mathcal{M}$ stores an external, updateable state of the world; and the Retriever $\mathcal{R}$ selects the subset of that memory that should condition the next reasoning or action step. Unlike one-shot retrieval-augmented generation, spatial memory must also support \emph{write} and \emph{update}: an agent may discover that a door is blocked, that an object has moved, or that a previous topological belief was wrong, and the memory state must change accordingly.

The memory bank is therefore not merely a text corpus. In spatial domains, $\mathcal{M}$ may contain vector embeddings of views and instructions, metric maps and voxel belief states, dynamic 3D scene graphs, frontier snapshots, past trajectories, user preferences, and brain-inspired episodic archives. The central problem is to make these heterogeneous memories queryable: semantic retrieval must locate relevant objects and events, while geometric retrieval must preserve relations such as containment, support, reachability, and path connectivity.

\begin{figure}[t]
    \centering
    \includegraphics[width=\linewidth]{figures/part3/section7/retrieval_memory_augmented_spatial_reasoning.pdf}
    \caption{Retrieval- and memory-augmented spatial reasoning: (i) a query and a memory bank of vector records, spatial scene graphs, and episodic archives, (ii) retrieval and update via spatial RAG, dynamic 3D scene graphs, or layered memory, and (iii) alternative output forms rather than joint outputs of one query. In the example, the keys were on the kitchen counter five minutes ago and are in the hallway now. Scenes and retrieval diagrams are illustrative, not measured similarity scores.}
    \label{fig:p3_s7_3_6_memory_formulation}
\end{figure}

\subsubsection{Task Definition}
Where §\ref{subsec:p3_s7_3_5} treats memory as one black-box tool among others, here the \emph{external memory itself} is the central learning target. Formally, a retrieval- \& memory-augmented system is the triple $(\mathcal{P}, \mathcal{M}, \mathcal{R})$, where $\mathcal{P}$ is a reasoning core (an LLM, VLM, or embodied planner), $\mathcal{M}$ is an external memory store, and $\mathcal{R}: \mathcal{Q} \times \mathcal{M} \to \mathcal{C}$ maps a query $q_t$ and the current memory state to a relevant context fragment $c_t$. Given an instruction $x$, an observation stream $o_t$, and a reasoning history $h_t$ (initially $h_0 = \varnothing$), the system proceeds by the iterative update
\begin{equation}
\begin{aligned}
q_t &= \mathrm{Query}_{\mathcal{P}}(x, o_t, h_{t-1}), &
c_t &= \mathcal{R}(q_t, \mathcal{M}_{t-1}), &
s_t &= \mathcal{P}(x, o_t, c_t, h_{t-1}), \\
\mathcal{M}_t &= \mathrm{Write}(\mathcal{M}_{t-1}, o_t, s_t, c_t), &
h_t &= h_{t-1} \cup \{(o_t, q_t, c_t, s_t)\},
\end{aligned}
\end{equation}
where $\mathrm{Query}_{\mathcal{P}}$ extracts the planner's information need, $\mathcal{R}$ returns the retrieved context, and $\mathrm{Write}$ inserts, revises, or invalidates memory entries using the newest observation and decision state. The decomposition brings three concrete advantages: \emph{scalability}, by letting agents operate over large environments while fetching only task-relevant subgraphs or vectors; \emph{personalization}, through targeted retrieval of user profiles, historical preferences, and episodic interactions; and \emph{explainability}, since decisions can be traced to the specific memories, rules, or map fragments that were retrieved.

\subsubsection{Methodology}
Three families of memory architecture instantiate the $(\mathcal{P}, \mathcal{M}, \mathcal{R})$ triple at different levels of abstraction and scale.

\paragraph{Retrieval-Augmented Generation for Navigation and Instruction.}
Vision-and-Language Navigation (VLN) demands a tight coupling of visual perception, language understanding, and action generation, and retrieval-augmented generation (RAG) bridges high-level planning and local observation by externalizing spatial context; the design question is what the memory bank should store and how the retriever should index it.

\emph{Textual abstraction vs.\ visual fidelity.} NavRAG~\cite{wang2025navrag} argues for hierarchical abstraction: a ``Hierarchical Scene Description Tree'' converts the 3D environment into structured text (Region $\rightarrow$ Object $\rightarrow$ Attribute), so the retriever can access information at the granularity the query requires, at the cost of visual nuance. VisRAG~\cite{yu2024visrag} takes the opposite stance with a multimodal pipeline that bypasses text parsing entirely: a VLM retrieves images directly from the query, preserving texture, lighting, and style that text-only retrieval discards.

\emph{Privileged maps vs.\ egocentric mental simulation.} A second line replaces privileged global maps with human-like episodic memory. GoViG~\cite{wu2025govig} generates instructions solely from egocentric observations of start and goal states, alternating between predicting the next visual view and generating the corresponding instruction, so that the navigational narrative is built on a hallucinated path; retrieval shifts from static facts to potential futures.

\emph{Hybridizing spatial and semantic logic.} Retrieval in physical spaces requires logic that pure semantic search cannot provide (\eg,\ ``near'', ``inside'', ``along''). Spatial-RAG~\cite{yu2025spatial} combines sparse spatial filtering (SQL-based constraints) with dense semantic matching, formulating retrieval as multi-objective optimization, which is what grounds CoT reasoning~\cite{liu2023retrieval} so that multi-step plans respect geometric feasibility.

\paragraph{Dynamic 3D Scene Graphs and Subgraph Retrieval.}
As environments grow, flat memory representations such as point clouds become intractable, and hierarchical abstraction through 3D scene graphs (3DSGs) has taken over: Hydra~\cite{hughes2022hydra} builds metric-semantic scene graphs in real time, and open-vocabulary successors embed vision--language features directly into graph nodes~\cite{gu2024conceptgraphs}, turning scene graphs from rigid structural maps into queryable vector stores. Because global graphs overwhelm planner context windows, processing has moved from holistic graph analysis to sparse subgraph retrieval, through semantic pruning~\cite{rana2023sayplan} or RAG-based mechanisms~\cite{xie2024embodiedrag} that isolate task-relevant topological neighborhoods. The convergence of graph and vector retrieval matters for spatial memory in particular: graph edges preserve physical relations, while vector neighborhoods preserve open-vocabulary semantic similarity.

\paragraph{Unified and Brain-Inspired Memory Architectures.}
RAG provides access to information; the organization of that information is equally critical, and recent architectures borrow the brain's division of labor. RoboMemory~\cite{lei2025robomemory} integrates working, short-term, and long-term memory in one cognitive framework so that the agent iteratively refines its model of the environment, and 3D-Mem~\cite{yang20253dmem} takes a visual-centric route, using ``frontier snapshots'' to manage exploration and hallucinate potential future states. 3DLLM-Mem~\cite{hu20253dllmmem} grounds the integration of retrieved memory into the inference stream with a memory fusion layer that cross-attends from the current instruction into the stored 3D memory bank:
\begin{equation}
f^Q_{\text{fuse}} = \mathrm{Softmax}\!\left( \frac{f^Q_t (f^K)^\top}{\sqrt{C}} \right) f^V,
\end{equation}
where $f^Q_t$ are the query features of the current instruction and $f^K$, $f^V$ the keys and values derived from the memory bank. Addressing specific failure modes, VLM$^2$~\cite{liu2025vision} introduces a dual-memory module for the semantic--geometric misalignment of video-based spatial reasoning, and Vision-to-Geometry~\cite{cai2025vision} incorporates geometric cues such as field-of-view coverage to prevent redundant re-exploration.



\subsubsection{Benchmarks}
Progress in memory-augmented spatial reasoning requires separating \emph{memory-specific} tests from broad spatial probes.

\paragraph{Memory-specific evaluation.}
3DMem-Bench~\cite{hu20253dllmmem} serves as a standard for assessing long-term persistence: comprising over 26{,}000 trajectories, it challenges agents with Spatial--Temporal QA (\eg,\ ``What color was the bed in the room you visited five minutes ago?''), requiring the retrieval of past observations that are no longer in the active context. Moving beyond static environments, EmbodiedBench~\cite{yang2025embodiedbench} tests long-horizon planning where objects are moved or manipulated during the task, benchmarking the \emph{write/update} capability of memory systems, their ability to overwrite obsolete data without catastrophic forgetting. OpenEQA~\cite{majumdar2024openeqa} offers a complementary episodic-memory setting by asking questions about an environment the agent has previously explored, making it useful for measuring whether retrieved experience rather than immediate perception supports the answer.

\paragraph{Broad spatial probes as diagnostic complements.}
Benchmarks such as SpaCE-10~\cite{gong2025space10}, SITE~\cite{wang2025site}, and the static probes collected in Table~\ref{tab:spatial_reasoning_benchmark} (OmniSpatial~\cite{jia2025omnispatial}, Spatial-DISE~\cite{huang2025spatialdise}, Open3DVQA~\cite{zhang2025open3dvqa}) evaluate spatial orientation, perspective taking, metric estimation, open-space 3D relations, and intrinsic/extrinsic static or dynamic reasoning. They are valuable for diagnosing what information a memory system should retain (viewpoint-dependent relations, object-centric attributes, multi-step multi-view dependencies), but they are not memory benchmarks unless the protocol hides necessary information in a prior episode and requires retrieval from $\mathcal{M}$: a model can fail a broad spatial probe because it lacks geometric perception, because it retrieved the wrong memory, or because the memory was never written in the first place.


\subsubsection{Challenges and Trends}
The recurring limitations originate in the Planner-Memory-Retriever interface rather than in any individual component.
\textbf{Query formation under unbounded context} is the central bottleneck, not memory size alone: the act of deciding what to retrieve requires compressing the current task, viewpoint, uncertainty, and history into a finite representation, and memory paging and scheduling alleviate but do not remove the need for a learned retrieval policy.
\textbf{Write/update consistency} is where append-only memories become dangerous: agents need mechanisms for reconciling conflicting observations, invalidating stale object locations, and representing uncertainty when the world may have changed since the last visit; EmbodiedBench-style protocols make the issue visible, but robust update semantics remain open.
\textbf{Integration mechanics} are unresolved between in-context learning, which concatenates retrieved text to the prompt (simple but token-inefficient), and cross-attention fusion, which injects retrieved 3D features into attention layers (3DLLM-Mem~\cite{hu20253dllmmem}; expressive but architecture-coupled), with no clear winner on the fidelity--modularity trade-off.
\textbf{Modality mismatch} has no general theory: RAG-over-text (NavRAG~\cite{wang2025navrag}) discards visual nuance, RAG-over-images (VisRAG~\cite{yu2024visrag}) loses structural priors, and when to retrieve text, images, or 3D representations remains task-dependent.

Memory is shifting \emph{from static maps to dynamic semantic belief states}: systems store probabilities, semantic labels, and interaction histories rather than coordinates, enabling agents to reason under uncertainty about where an object \emph{likely} resides given commonsense and past observation, as RoboMemory~\cite{lei2025robomemory} typifies.
Graph- and vector-based retrieval are converging into hybrid architectures that use vector search to identify graph entry nodes and topological traversal for neighborhood reasoning.
Most abstractly, \emph{meta-cognitive memory} such as Meta-Memory~\cite{mao2025meta} generates task-specific ``cognitive maps'' and retrieves \emph{strategies} rather than raw data, pointing toward learnable retrieval policies and bridging into the agentic orchestration of §\ref{subsec:p3_s7_3_7}.


\subsection{Agentic Spatial Reasoning}
\label{subsec:p3_s7_3_7}

Agentic spatial reasoning asks how a model organizes \emph{interaction} with space. The core object is not a single chain-of-thought prompt, tool call, memory store, VLA policy, or world model, but an orchestration policy that repeatedly decides what to perceive, what to infer, what to do, and how to revise its state after feedback. In this closed-loop view, a wall is not only a geometric plane but also a constraint on motion; a cup is not only an object category but also a potential affordance; and a failed action is not only an error but also evidence that should reshape the next plan.

``Agentic'' is therefore a composition principle: the rationales, programs, tools, memory, VLA policies, and world models of the preceding subsections become resources inside a perceive--reason--act--reflect loop, and the research question is how the loop is scheduled and corrected: when should the agent gather more evidence, retrieve memory, call a geometric tool, simulate a future, ask another agent, or execute an action?

\begin{figure}[t]
    \centering
    \includegraphics[width=\linewidth]{figures/part3/section7/agentic_spatial_reasoning.pdf}
    \caption{Agentic spatial reasoning as closed-loop orchestration. Top: (i) instruction, observation, robot state, tools, and memory, (ii) the perceive--reason--act--reflect loop under planning core $\mathcal{P}$ and policy $\pi$, and (iii) complementary, not simultaneous, output forms; the door-state sequence is a prediction and the failed counter search is feedback from a different step. Bottom: (a)--(d) mechanisms that support the loop. The scene is a generated illustration, not an experimental observation.}
    \label{fig:p3_s7_3_7_agentic_formulation}
\end{figure}

\subsubsection{Task Definition}
Whereas §\ref{subsec:p3_s7_3_5} treats tools as black-box operators and §\ref{subsec:p3_s7_3_6} treats external memory as the central learning target, agentic spatial reasoning composes both into a closed-loop cycle in which the orchestration policy itself becomes the learning target. Formally, an agentic spatial reasoning system is the tuple $(\mathcal{P}, \mathcal{S}, \mathcal{A}, \mathcal{F}, \pi)$, where $\mathcal{P}$ is a planning core (an LLM or VLM), $\mathcal{S}$ is the spatial state space (perception, internal map, history), $\mathcal{A}$ is the action space (sensorimotor commands or sub-goals), $\mathcal{F}$ is a reflection / critic function that scores proposed actions or rationales, and $\pi$ is the orchestration policy. Given an instruction $x$, the system maintains a history $h_t$ (initially $h_0 = \varnothing$) and proceeds by the iterative perceive--reason--act--reflect cycle
\begin{equation}
\begin{aligned}
o_t &= \mathrm{Perceive}(\text{env}), &
s_t &= \mathcal{P}(o_t, h_{t-1}), \\
a_t &\sim \pi(\cdot \mid s_t), &
r_t &= \mathcal{F}(s_t, a_t, o_{t+1}), \\
h_t &= h_{t-1} \cup \{(o_t, s_t, a_t, r_t)\},
\end{aligned}
\end{equation}
where $\mathrm{Perceive}$ extracts the current spatial observation, $\pi$ samples a proposed action, $\mathcal{F}$ supplies a reflection signal (verbal critique, planner-critic verdict, or value estimate), and the history accumulates all four. This formulation makes explicit that ``agentic'' is not a new capability but rather a structural promotion: the orchestration policy $\pi$ becomes the central learning target rather than the planner $\mathcal{P}$ or any individual tool.

\subsubsection{Methodology}
Four orchestration problems organize the methods: grounding observations into executable state, closing the feedback loop, coordinating across agents, and using action and world-model modules without reducing agency to any one of them.

\paragraph{Grounding Observations into Executable State.}
The first responsibility of an agentic loop is to convert an instruction and observation into a state that can support action, preserving referential consistency between the agent's internal model and the physical environment.
Probabilistic affordance alignment, as in SayCan~\cite{ahn2022icanisay}, PaLM-E~\cite{driess2023palm}, and RT-2~\cite{zitkovich2023rt2}, weights the likelihood of an instruction (``pick up the apple'') by the geometric feasibility of the action in the current scene; spatial reasoning stays implicit, and the orchestration layer must decide whether that is sufficient or whether to request an explicit representation.
Explicit intermediate representations supply the precision implicit alignment lacks: SpatialCoT~\cite{liu2025spatialcot} compels the VLM to output metric coordinates inside its chain of thought, a ``think-coordinate-act'' cycle that creates a verifiable spatial bottleneck; Embodied-R1~\cite{yuan2025embodiedr1} uses reinforcement learning to refine a pointing representation so that the model outputs 3D reference points (spatial anchors) that decouple high-level planning from motor control; and procedural querying treats the environment as a queryable database, through dynamically synthesized APIs (VADAR, §\ref{subsec:p3_s7_3_4}) or a stateful code kernel that executes one cell at a time and revises the next from the observed outputs (SpatialClaw, §\ref{subsec:p3_s7_3_5}).
Beyond the immediate percept, agents ground reasoning in persistent structured abstractions: object-centric 3D scene graphs such as ConceptGraphs (§\ref{subsec:p3_s7_3_6}) let the agent reason over topological relations (``behind'', ``supported by'') rather than raw geometry, and MapGPT~\cite{chen2024mapgpt} converts topological maps into textual prompts that give the LLM a global view of connectivity for long-horizon planning.

\paragraph{Closing the Feedback Loop.}
Agentic spatial reasoning is a temporal \emph{process} rather than a single inference, and the methodological shift is from producing one rationale to maintaining a loop that observes the consequences of its own choices.
The chain-of-thought paradigm (§\ref{subsec:p3_s7_3_3}) improved reasoning depth but remained open-loop; ReAct~\cite{yao2022react} closed the loop by interleaving reasoning traces with executable actions so that spatial outcomes are observed before the next step is planned. Observation alone is often insufficient for correction, so Reflexion~\cite{shinn2023reflexion} introduced ``verbal reinforcement learning'', in which the agent writes explicit critiques of its failures into short-term context, and DEPS~\cite{wang2023describe} forces the agent to \emph{explain} failures in geometric terms before choosing new sub-goals, turning retrying into debugging.
For coherence across long horizons, agents need memory beyond the context window: Voyager (§\ref{subsec:p3_s7_3_4}) stores successful spatial behaviors as executable code in a vector database and retrieves them when a similar context recurs, consolidating transient experience into procedural knowledge.
Reflection can also precede execution. Reasoning via Planning (RAP)~\cite{hao2023reasoning} and Tree of Thoughts~\cite{yao2023tot} explore branching spatial possibilities in a latent space; STMA~\cite{lei2025stma} ties such simulation to dynamic knowledge graphs with planner--critic loops that validate whether a proposed path is topologically valid; and vision-language world models such as VLWM~\cite{chen2025vlwm} roll out semantically described future states so that the orchestration policy can compare hypothetical futures before committing.

\paragraph{Coordinating Multiple Spatial Agents.}
In shared environments spatial reasoning becomes distributed, and movement has to be negotiated across perspectives and intents.
Multi-agent reinforcement learning induces spatial cooperation without explicit communication, as in the SMAC formations learned by decentralized actor--critic and value-decomposition methods~\cite{samvelyan2019smac} (§\ref{subsec:p3_s7_2_6}), but there each agent treats the others as dynamic obstacles and has no symbolic model of its partners.
A \emph{spatial theory of mind} adds the ability to model another agent's geometric viewpoint, transforming egocentric perception into an allocentric frame (``my left'' is ``your right''): RoCo~\cite{mandi2024roco} uses foundation models for embodied dialog in which agents disambiguate spatial references by simulating the partner's visual field, and ProAgent~\cite{zhang2024proagent} shows agents communicating to update a shared map and synchronize movement, converging on a common ``language of space''.

\paragraph{Orchestrating Action and World-Model Modules.}
Generalist VLA policies and world models are components of agentic systems, but they do not by themselves define agency; they become agentic when an orchestration policy uses them inside the loop, with a VLA module executing a sub-goal, a world model predicting the consequence of a candidate action, and a reflection module deciding whether to retry, replan, or gather more evidence.
What makes an action representation useful to a loop is that it exposes intermediate commitments the loop can inspect, revise, or hand to a controller: beyond the unified action interfaces of RT-2~\cite{zitkovich2023rt2}, Gato~\cite{reed2022generalist}, and Octo~\cite{team2024octo}, world-language-action models~\cite{yang2026wla} jointly predict textual subtasks, subgoal images, and robot actions, and ACoT-VLA~\cite{zhong2026acotvla} reasons through coarse action intents before fine-grained control.
Cross-embodiment transfer (§\ref{subsec:p3_s7_3_2}; RoboCat's self-improvement cycles adapt to new bodies~\cite{bousmalis2023robocat}) becomes an orchestration problem too, since the same high-level plan may require different perception checks, affordance tests, and recovery routines on different bodies.
World models (§\ref{subsec:p3_s7_3_8}), from DreamerV3~\cite{hafner2023mastering} and Genie~\cite{genie3} to WorldPilot-style prediction inside planning loops~\cite{lin2026worldpilot}, serve the loop as cognitive simulators: the agentic use is not video generation for its own sake but closed-loop decision support, in which the agent imagines, critiques, and then acts.





\subsubsection{Benchmarks}
Evaluating agentic spatial reasoning requires benchmarks that go beyond static visual question answering: a useful protocol tests whether the agent can gather missing observations, choose tools or actions, use memory across long horizons, and revise its plan after feedback, and static spatial probes should be read as component diagnostics rather than direct evidence of closed-loop agency.
Two protocols exercise the loop directly: iVISPAR~\cite{mayer2025ivispar} evaluates VLMs as agents in an interactive sliding-tile setting with visual 2D, visual 3D, and text inputs, where the answer depends on multi-step state transitions rather than one-shot recognition, and TIR-Bench~\cite{li2025tirbench} requires agents to use tools (zoom, crop, rotate) to manipulate images while solving spatial puzzles.
Memory-centered settings (OpenEQA~\cite{majumdar2024openeqa}, §\ref{subsec:p3_s7_3_6}) and open-3D probes such as Open3DVQA~\cite{zhang2025open3dvqa} test what the loop must retain, and the broad visual-spatial probes (SpatialEval~\cite{wang2024picture}, VSI-Bench~\cite{yang2025vsibench}, OmniSpatial~\cite{jia2025omnispatial}, Spatial-DISE~\cite{huang2025spatialdise}) identify weak spatial primitives along the axes of static vs.\ dynamic content, single vs.\ multi-view, ego vs.\ exo perspective, and seen vs.\ unseen queried information; Table~\ref{tab:spatial_reasoning_benchmark} maps the coverage.

\begin{table*}[!ht]
  \caption{Representative spatial-reasoning benchmarks discussed in §\ref{subsec:p3_s7_3_7}. Perspective marks whether ego and exo viewpoints are explicitly represented; Interactive marks protocols that exercise action or tool selection rather than only answer generation; Query Target is whether the queried information is visible (Seen), hidden (Unseen), or an interactive state (State). \textit{Undeclared} marks attributes the source paper does not state.}
  \label{tab:spatial_reasoning_benchmark}
  \centering
  \resizebox{\textwidth}{!}{
  \begin{tabular}{@{}lccccccccc}
    \toprule
      \multirow{2}{*}{\textbf{Benchmark}}
      & \multirow{2}{*}{\textbf{QA / Tasks}}
      & \multirow{2}{*}{\textbf{Modality}}
      & \multirow{2}{*}{\textbf{Scene}}
      & \multicolumn{2}{c}{\textbf{View Type}}
      & \multicolumn{2}{c}{\textbf{Perspective}}
      & \multirow{2}{*}{\textbf{Interactive}}
      & \multirow{2}{*}{\textbf{Query Target}} \\
      \cmidrule(lr){5-6}
        \cmidrule(lr){7-8}
        & & &
        & \textbf{Single} & \textbf{Multi}
        & \textbf{Ego} & \textbf{Exo}
        & & \\
    \midrule
    3DSRBench~\cite{ma20253dsrbench} & 2{,}772 & Image & Static & \checkmark  & \crossmark & \crossmark  &  \checkmark & \crossmark & Seen\\
    InternSpatial-Bench~\cite{deng2025internspatial} &  6{,}008 & Image & Static & \checkmark  & \crossmark & \crossmark  &  \checkmark & \crossmark & Seen \\
    OmniSpatial~\cite{jia2025omnispatial} & $\sim$8{,}400 & Image & Static & \checkmark & \crossmark & \checkmark & \checkmark  & \crossmark & Seen\\
    Spatial-DISE~\cite{huang2025spatialdise} & 559 & Image & Mixed & \checkmark & \checkmark & \multicolumn{2}{c}{\cellcolor{lightgray}\textit{Taxonomy}} & \crossmark & Seen\\
    Open3DVQA~\cite{zhang2025open3dvqa} & 9{,}000 & 3D / Sim & Static & \checkmark & \crossmark & \checkmark & \checkmark & \crossmark & Seen\\
    iVISPAR~\cite{mayer2025ivispar} & Puzzles & 2D/3D/Text & Dynamic & \checkmark & \crossmark & \multicolumn{2}{c}{\cellcolor{lightgray}\textit{N/A}} & \checkmark & State\\
    Spatial457~\cite{wang2025spatial457} & 23{,}752 & Image & Static & \checkmark & \crossmark & \multicolumn{2}{c}{\cellcolor{lightgray}\textit{Undeclared}}  & \crossmark & Seen\\
    All-Angles Bench~\cite{yeh2025seeing} & $\sim$2{,}100 & Image & Static  & \crossmark & \checkmark & \multicolumn{2}{c}{\cellcolor{lightgray}\textit{Undeclared}} & \crossmark  & Seen \\
    EmbSpatial-Bench~\cite{du2024embspatial} & 3{,}640 & Image & Static & \checkmark  & \crossmark & \checkmark & \crossmark & \crossmark &  Seen \\
    MMSI-Bench~\cite{yang2025mmsi} & 1{,}000 & Image & Static & \crossmark & \checkmark  & \multicolumn{2}{c}{\cellcolor{lightgray}\textit{Undeclared}} & \crossmark & Seen\\
    MindCube~\cite{yin2025spatial} & 21{,}154 & Image & Static & \crossmark & \checkmark & \multicolumn{2}{c}{\cellcolor{lightgray}\textit{Undeclared}}  & \crossmark & Unseen \\
    VSI-Bench~\cite{yang2025vsibench} & 5{,}130 & Video & Static & \checkmark & \crossmark & \checkmark & \crossmark & \crossmark & Seen\\
    VLM4D~\cite{zhou2025vlm4d} & $\sim$1{,}800 & Video & Dynamic & \checkmark & \crossmark &  \checkmark & \checkmark & \crossmark & Seen \\
    STI-Bench~\cite{li2025sti} & 2{,}064 & Video & Dynamic & \checkmark & \crossmark & \checkmark & \checkmark & \crossmark &  Seen\\
    DSI-Bench~\cite{zhang2025dsi} & $\sim$1{,}700 & Video & Dynamic & \checkmark & \crossmark & \multicolumn{2}{c}{\cellcolor{lightgray}\textit{Undeclared}} & \crossmark &  Seen\\
    \bottomrule
  \end{tabular}
  }
\end{table*}

Table~\ref{tab:spatial_reasoning_benchmark} should be read as a coverage map, not as a leaderboard for agentic systems. Most entries evaluate static answer generation or localized visual-spatial primitives; only interactive protocols such as iVISPAR and tool-use settings directly exercise action selection and feedback. Claims about agentic spatial reasoning should therefore report loop-level metrics (observation gathering, action choice, recovery after failure, memory update, and final task success) rather than single-turn accuracy alone.

\subsubsection{Challenges and Trends}
Boundaries are collapsing between language and action (VLAs), between robot bodies (cross-embodiment), and between perception and imagination (world models), and two difficulties gate progress toward generalist spatial agents.
The \textbf{sim-to-real fidelity gap}: high-fidelity simulators (predictive world models, generative video) accelerate training but introduce systematic biases that physical execution exposes, and the cost--precision frontier for transfer remains open.
\textbf{Long-horizon spatiotemporal coherence}: across multi-hour episodes the agent's internal map drifts, and episodic memory and reflection reduce but do not eliminate the drift, because the bottleneck is the finite context window of the planning core.

\emph{Explicit orchestration over observation--state--action--reflection tuples} is one direction: VLA and WLA systems are beginning to expose intermediate subgoals, action intents, and predicted states that a loop can inspect rather than hiding all control inside a single forward pass. A second is \emph{mental simulation via generative world models}: DreamerV3-style, Genie-style, VLWM-style, and WorldPilot-style models let agents verify spatial logic in imagination before acting, shifting the paradigm from ``trial-and-error'' toward ``trial-in-imagination'' while keeping the agentic focus on when to simulate, when to act, and when to revise.


\subsection{World-Model-based Spatial Reasoning}
\label{subsec:p3_s7_3_8}


% Test-Time Scaling with World Models for Spatial Reasoning
% When and How Much to Imagine: Adaptive Test-Time Scaling with World Models for Visual Spatial Reasoning
% From 2D to 3D Cognition: A Brief Survey of General World Models
% Vagen: Reinforcing world model reasoning for multi-turn vlm agents
% Scaling and Beyond: Advancing Spatial Reasoning in MLLMs Requires New Recipes
% SpatialDreamer: Incentivizing Spatial Reasoning via Active Mental Imagery
% Think with 3D: Geometric Imagination Grounded Spatial Reasoning from Limited Views



% 在这里想要重新梳理一下应该描述的内容：
% 在这个章节中更想强调的是如何使用world model as a tool来实现spatial reasoning task，而不是去讨论world model本身的训练和发展。也就是说，world model是一个工具，spatial reasoning task是一个目标。我们要定义的是如何用这个工具去实现这个目标。
% 并且需要在文中也要提到的是，world model的具体构建和训练是在后续的章节8.6 World Modeling中

% world 作为一个工具其核心的能力就是模拟了世界，因此能够在spatial reasoning过程中提供一些反馈，这个反馈可以是多种形式的，比如生成未来的图片，或者是生成未来的状态，或者是生成未来的动作序列。然后这些反馈可以被用来提升spatial reasoning。
% 另外，world model 也可以得到 latent representation，比如利用JEPA得到的 latent representation，这些latent representation也可以被用来提升spatial reasoning。


A world model is a neural simulator of how world states transition, conditioned on an agent's actions and/or external multi-modal context.
Trained on large-scale episode trajectories, it acquires a generalizable, latent model of environment dynamics.
This subsection treats that model as a \emph{tool}: the object of study is the spatial-reasoning task, and the world model is the instrument used to solve it.
Its defining capability is that it \emph{simulates} the world and can therefore return a predictive feedback signal, a rollout of imagined futures or a dynamics-aware latent, that a reasoner consumes to answer a spatial query or to plan.
We accordingly organize the discussion by \emph{how the world model is queried and what feedback it returns}, rather than by how it is built.
The complementary question of how world models are themselves constructed, trained, and scaled is the subject of §\ref{sec:worldmodeling} (World Modeling); here we take a pretrained world model as given.
Figure~\ref{fig:p3_s7_3_8_world_model_formulation} schematizes this tool-use formulation; its two usage modes, understanding-oriented and planning-oriented, cut across the feedback forms surveyed below.

\begin{figure}[!t]
    \centering
    \includegraphics[width=0.99\linewidth]{figures/part3/section7/world_model_based_spatial_reasoning.pdf}
    \caption{World-model-based spatial reasoning read as tool use: (i) a spatial query $q$ with context $c$ and optional actions $a_{0:T-1}$; (ii) a pretrained world model (constructed and trained as in §\ref{sec:worldmodeling}) whose predictive feedback $f$ takes four possibly overlapping forms, namely (a) future observations, (b) future states, (c) action-conditioned trajectories, and (d) predictive latents; and (iii) a readout $y=g(c,f)$ that gives an inferred relation (understanding) or a selected plan (planning). The robot scenes are illustrative, not experimental observations.}
    \label{fig:p3_s7_3_8_world_model_formulation}
\end{figure}

\subsubsection{Task Definition}
% 这里应该是定义world model如何帮助spatial reasoning task，而不是什么是world model。也就是说，world model是一个工具，spatial reasoning task是一个目标。我们要定义的是如何用这个工具去实现这个目标。

The task addressed in this subsection is spatial reasoning itself:
given a query $q$ about a scene together with the available context $c$ (a textual description, one or more observations, or a scene specification), a reasoner must produce a spatial answer $y$, either an inferred spatial property $o$ (\eg, a relation, a structure, or a physical attribute) or a plan / policy $\pi$ that acts on the scene.
What distinguishes the methods surveyed here is \emph{how} they arrive at $y$: rather than mapping $c$ to $y$ in a single feed-forward pass, the reasoner consults a \emph{world model} as an auxiliary instrument that can imagine how the scene would look or evolve, and reasons over that imagined evidence.
Concretely, the world model is a pretrained simulator exposing a distribution
$\mathbb{P}(s_0, s_1, \dots, s_T \mid a_0, a_1, \dots, a_{T-1}, c)$
over world states $s_t$, or over their partial observations $\operatorname{proj}(s_t)$ when full states are unavailable, optionally conditioned on an action sequence $\left\{a_t\right\}_{t=0}^{T-1}$ (null when the model is action-free).
A world-model-based reasoner treats $\mathbb{P}$ as a queryable tool: given the query, it (i) conditions $\mathbb{P}$ on the context $c$ and, when planning, on a candidate action sequence; (ii) draws predictive feedback $f$ from $\mathbb{P}$, a rollout of future observations, future states, or action-conditioned trajectories, or a dynamics-aware latent; and (iii) applies a readout $g$ that maps the feedback back to the spatial answer, $y = g(c, f)$.
The world model supplies the imagined evidence $f$, but producing the spatial answer $y$ remains the objective, and the fidelity of $f$ bounds how much the world model can help.





\subsubsection{Methodology}

We treat the world model as a pretrained instrument and organize methods by \emph{the form of predictive feedback it returns} to the reasoner.
Two families emerge: \emph{generative feedback}, in which the model is rolled out to synthesize explicit futures that the reasoner inspects, and \emph{latent-representation feedback}, in which the model's dynamics-aware embedding is consumed directly without decoding; the understanding-versus-planning distinction cuts across both.

\paragraph{Generative Feedback via Future Observations.}
The most direct use of a world model is to synthesize future or counterfactual \emph{pixels} and read spatial evidence off them, an analysis-by-synthesis strategy in which the reasoning outcome $o$ (a spatial relationship, structure, or physical property) is whatever best reproduces the observed or imagined imagery under $\mathbb{P}$.
Predating deep generative models, Galileo~\cite{wu2015galileo} instantiates this idea with a physics engine in place of a neural simulator: it tracks object velocities across raw videos and recovers 3D shape, mass, and friction as the parameters that best reproduce the observed motion, so that the projected state $\operatorname{proj}(s_t)$ is the object velocity $v_t$ and the videos themselves define the transition distribution $\mathbb{P}(v_0, \dots, v_T \mid c)$.
Modern instances replace the engine with an Internet-scale image or video diffusion prior: DragAPart~\cite{li2024dragapart}, Puppet-Master~\cite{li2024puppetmaster}, ATOP~\cite{vora2025articulate}, and DreamArt~\cite{lu2025dreamart} fine-tune such priors into articulation simulators $\mathbb{P}(I_0, \dots, I_T \mid c)$ whose frames $I_t = \operatorname{proj}(s_t)$ render an articulated object under a hypothesized configuration, and recover the 3D articulation (\eg, its motion axes) by aligning the rendered hypothesis with the synthesized video.
A second, inference-time use synthesizes \emph{unobserved viewpoints} to resolve spatial questions that a single view underdetermines: MindJourney~\cite{yang2026mindjourney} lets a vision-language model ``imagine'' the scene from a new camera pose with a view-conditioned world model before answering, turning the world model into a test-time reasoning aid rather than a training-time simulator.
Because such imagination is expensive and not always helpful, AVIC~\cite{yu2026avic} learns \emph{when} and \emph{how much} to imagine, gating world-model invocation on the sufficiency of the current visual evidence and matching or exceeding fixed-imagination strategies at a fraction of the world-model calls.
Pushing imagination from a fixed sweep to an \emph{active} policy, SpatialDreamer~\cite{cao2025spatialdreamer} closes the loop: a reinforcement-learned agent chooses which viewpoints to imagine through the world model and then grounds its answer in the imagined evidence, under a geometric policy-optimization objective that rewards geometrically consistent rollouts.

\paragraph{Generative Feedback via Future States.}
Rather than decode pixels, a second line rolls the world model out in a \emph{structured or higher-dimensional state space}, trading photorealism for geometric consistency.
TesserAct~\cite{zhen2025tesseract} learns a 4D (3D-plus-time) world model whose rollouts stay spatially coherent enough to support robotic manipulation, and persistent-3D world models~\cite{zhou2025learning} maintain an explicit 3D memory so that revisited regions remain consistent across long horizons.
State-space feedback is attractive when the downstream query is itself geometric, since it sidesteps the appearance-modeling burden that pixel rollouts carry.

\paragraph{Generative Feedback via Action-Conditioned Trajectories.}
When the goal is to \emph{act}, the world model is queried as a simulator: candidate action sequences are rolled out and scored by the favorability of the resulting states, yielding a policy $\pi:(s_{\leq t}, a_{<t}) \mapsto a_t$ that steers the episode toward a preferred terminal state $s_T \in S_{\text{good}}$ (the model-based-RL view; the broader treatment of RL with spatial inductive biases is §\ref{subsec:p3_s7_2_5}).
The paradigm originates with Ha and Schmidhuber~\cite{ha2018world}, who introduced the term ``world model'' in its current sense, training a generative recurrent model of future observations and using it as an internal simulator for policy search, and with Dreamer~\cite{hafner2019dreamer}, which learned the simulator in a compact latent space and optimized the policy purely on imagined rollouts.
Scaling this recipe, UniSim~\cite{yang2023learning} builds a video-diffusion simulator from which real-world manipulation policies can be extracted, and DreamGen~\cite{jang2025dreamgen} uses video world models as trajectory generators, pseudo-labeling the synthesized rollouts with an inverse-dynamics model to train policies that generalize to new behaviors and environments.
Most recently, VAGEN~\cite{wang2025vagen} turns this feedback back onto the reasoner itself, using reinforcement learning to \emph{supervise} a vision-language agent's internal world-model reasoning: it decomposes that reasoning into state estimation (``what is the current state?'') and transition modeling (``what comes next?'') and rewards accurate latent prediction.

\paragraph{Latent-Representation Feedback.}
A final family dispenses with generation entirely and consumes the world model's \emph{predictive latent} as the spatial signal.
Joint-embedding predictive architectures learn to predict the representation of a future or masked observation rather than its pixels: I-JEPA~\cite{assran2023jepa} establishes the recipe for images, and V-JEPA~2~\cite{assran2025vjepa} extends it to video, yielding dynamics-aware embeddings that support understanding, prediction, and planning without ever rendering a frame, while VL-JEPA~\cite{chen2025vl} ports the idea to vision-language inputs.
A related line keeps the latent explicitly \emph{geometric}: 3DThinker~\cite{chen20253dthinker} performs 3D mental imagery during reasoning by generating a 3D latent and aligning it with that of a 3D foundation model, letting a vision-language model reason about spatial relations from only a few views without any 3D input or explicit rendering.
Because the feedback is an embedding rather than an image, latent world models avoid the appearance-synthesis cost and the pixel-space memory drift that limit generative rollouts, at the price of a signal that is harder to inspect directly; they are an increasingly prominent, query-efficient alternative for spatial reasoning.

\subsubsection{Benchmarks}
Evaluation follows the usage mode.
Planning-oriented use is scored on the embodied suites of §\ref{subsec:p3_s7_3_2} and on long-horizon household environments such as ALFWorld~\cite{shridhar2020alfworld} and VirtualHome~\cite{puig2018virtualhome}, where the question is whether policies extracted from imagined rollouts transfer to the real environment.
Understanding-oriented use is scored on the static spatial probes that single views underdetermine, VSI-Bench~\cite{yang2025vsibench} and MindCube~\cite{yin2025spatial} among them, under test-time-scaling protocols that vary how many imagined views the reasoner is allowed and whether a verifier selects among them~\cite{jha2025probing,yu2026avic}.
The world model itself is assessed on video-generation quality (FVD, FID) together with task-specific physics-consistency probes, and surveys of general world models track the field's shift from 2D toward 3D-consistent simulation~\cite{xie2025worldsurvey}, while recent analyses argue that advancing MLLM spatial reasoning demands purpose-built recipes and data rather than scale alone~\cite{zhang2025scaling}.
No single benchmark currently exercises a world model along perception fidelity, action conditioning, and long-horizon spatial memory at once, a gap the challenges below make explicit.


\subsubsection{Challenges and Trends}
Reading the world model as a tool clarifies where the difficulties lie: some are inherited from the instrument itself, and are properly the concern of §\ref{sec:worldmodeling}, while others are specific to how a reasoner queries it. The most active questions of 2025--2026 are no longer whether a world model \emph{can} imagine, but whether its imagined evidence can be \emph{trusted} and whether consulting it is worth the cost.

Imagination is not a free lunch. Systematic probes report that world-model imagination improves spatial reasoning in only a minority of cases, and that adding imagined views yields non-monotonic, sometimes negative accuracy gains at steep token and latency cost~\cite{yu2026avic,jha2025probing}; the sheer amount of imagination is an unreliable proxy for accuracy~\cite{zhang2025scaling}, which recasts ``when and how much to imagine'' as a first-class test-time-scaling problem, with adaptive gating~\cite{yu2026avic} an early answer and a predictive account of \emph{which} queries need the world model still missing.
Faithfulness is the companion problem: a world model can render fluent futures that are physically or relationally wrong, so likelihood under $\mathbb{P}$ is necessary but not sufficient, and when the imagined evidence lacks fidelity verification buys little, since heuristic verifiers give limited calibration gains and an information bottleneck sets in, as a reasoner cannot recover a signal the imagination never contained~\cite{jha2025probing}. Grounding imagined evidence to the actual scene and calibrating trust in it remain open; physics-consistency objectives and geometry-grounded latents~\cite{chen20253dthinker} are partial responses.

The reasoner can be no more reliable than the simulator it queries. Even state-of-the-art models such as Genie 3~\cite{genie3}, despite markedly longer-horizon consistency, still drift in spatial memory, and consistency mechanisms~\cite{zhou2025learning,li2025vmem} narrow but do not close this gap, with hallucination and long-horizon error compounding over rollout length; raising this ceiling is the province of world-model construction (§\ref{sec:worldmodeling}), and the reasoning methods surveyed here can only mitigate its symptoms.

The feedback forms also trade off along different axes, with no principled account of which a given query should use: pixel rollouts are inspectable but expensive and prone to appearance drift; state or 4D rollouts are geometrically consistent but require structured simulators; latent feedback is cheap but opaque; action rollouts commit early to a policy; and most systems commit to one form a priori. Extending world-model-based reasoning beyond pixels (\eg, to forces~\cite{gillman2025force}, tactile, or audio modalities) is gated by the absence of internet-scale pre-training data in those modalities.

Two shifts define the current frontier. \emph{From external tool to trained-in reasoning}: rather than invoking a frozen simulator at inference time, recent work \emph{trains} the reasoner's use of the world model, whether by reinforcement learning that supervises an internal state-estimation and transition-modeling loop~\cite{wang2025vagen} or that rewards geometrically consistent active imagination~\cite{cao2025spatialdreamer}, or by distilling the world model's imagined dynamics directly into a vision-language model so that no simulator need be queried at test time~\cite{zhang2026world2vlm}. \emph{From pixels to predictive latents}: the imagined branch is increasingly consumed as a latent process whose hidden states guide reasoning, rather than as a faithfully rendered video~\cite{assran2023jepa,assran2025vjepa,chen20253dthinker}. Both shifts move the field away from ever-larger, faithfully-rendered pixel simulators and toward cheaper, more trustworthy, query-efficient feedback that a reasoner can afford to consult, and can be trained to consult well.









% \subsection{Summary}
% \label{subsec:p3_s7_3_summary}

% The eight subsections vary \emph{where} and \emph{how} the foundation model intervenes in the reasoning loop: as the reasoning core itself (§\ref{subsec:p3_s7_3_1}--§\ref{subsec:p3_s7_3_2}), as a generator of intermediate rationales or programs (§\ref{subsec:p3_s7_3_3}--§\ref{subsec:p3_s7_3_4}), as an orchestrator over tools, memory, and a closed interaction loop (§\ref{subsec:p3_s7_3_5}--§\ref{subsec:p3_s7_3_7}), or beside a learned simulator that supplies the spatial substrate (§\ref{subsec:p3_s7_3_8}).
% Table~\ref{tab:p3_s7_3_paradigm_comparison} summarizes the eight along the axes used for the symbolic paradigms in §\ref{sec:p3_s7_1} and the learning-based paradigms in §\ref{sec:p3_s7_2}.

% \begin{table}[t]
%     \centering
%     \fontsize{8}{10}\selectfont
%     \setlength{\tabcolsep}{2pt}
%     \caption{Comparison of the eight paradigms of foundation-model-driven spatial reasoning covered in §\ref{sec:p3_s7_3}, on the same axes as the comparison tables of §\ref{sec:p3_s7_1} and §\ref{sec:p3_s7_2}.}
%     \label{tab:p3_s7_3_paradigm_comparison}
%     \begin{tabularx}{\textwidth}{@{}p{2.9cm} p{3.3cm} p{3.3cm} X X@{}}
%         \toprule
%         \textbf{Paradigm} & \textbf{Locus of foundation-model intervention} & \textbf{Canonical inputs / outputs} & \textbf{Strengths} & \textbf{Limitations} \\
%         \midrule
%         §\ref{subsec:p3_s7_3_1} LLM/VLM-centric reasoning   & Pre-trained model \emph{is} the reasoning core; intervention via prompt / train / data / model-centric enhancement & Text-only or image$+$query $\to$ qualitative relation, distance, orientation, navigation plan & Zero-shot deployment; rapid prompt iteration; benefits from any backbone improvement                  & Ungrounded under text-only; metric-precision gap under VLM; reference-ambiguity in embodied use \\
%         §\ref{subsec:p3_s7_3_2} VLA-based reasoning         & Pre-trained model also emits low-level actions (action tokens / deltas / waypoints)                                & Image $+$ proprioception $+$ instruction $\to$ executable action sequence                            & End-to-end semantic-to-motor pipeline; cross-embodiment positive transfer via large data            & Semantic--physical misalignment; pseudo-interaction; inference-latency under autoregressive decoding \\
%         §\ref{subsec:p3_s7_3_3} Chain-of-thought (CoT \& variants) & Pre-trained model, with intermediate textual rationale as the prior; five topologies (CoT/ToT/GoT/PoT/CoA) unified as guided search & Query $\to$ rationale trace $\to$ grounded answer (qualitative relation, coordinate, plan, action) & Multi-step decomposition; verifier-friendly intermediate states; search-quality tunable               & Spatial-coordinate hallucination under simple textual CoT; rationale not grounded in truth signal; ToT/GoT cost scales super-linearly in depth \\
%         §\ref{subsec:p3_s7_3_4} Code / DSL generation         & Pre-trained model generates an executable program; environment supplies the truth signal                          & Query$+$API $\to$ program $P$ $\to$ executed output (answer with trace, or environment change)   & Interpretable rationale trace; execution-grounded; composes existing perception or control tools     & Reliance on a fixed human-authored DSL; brittle code execution under upstream-module error; sample-inefficient debug loops \\
%         §\ref{subsec:p3_s7_3_5} Tool-augmented reasoning      & Pre-trained model orchestrates a tool catalogue $\mathcal{T}$ via policy $\pi_\phi$, along execution / representation / validation / memory axes & Query$+$observation$+$tool set $\to$ history-conditioned answer $+$ tool-call trace             & Modular extension of perceptual / geometric / verification capacity; orchestration policy as the learning target & Combinatorial exploration in tool-action space; metric gap in the LLM--tool interface; conceptual overlap with §\ref{subsec:p3_s7_3_4} / §\ref{subsec:p3_s7_3_6} \\
%         §\ref{subsec:p3_s7_3_6} Retrieval- \& memory-augmented & External memory bank itself is the central learning target; Planner--Memory--Retriever triple $(\mathcal{P}, \mathcal{M}, \mathcal{R})$ & Query$+$scene observation$+$memory $\to$ action plan, retrieved subgraph, retrieved past observation & Kilometer-scale environment coverage; personalization; explainable trace of retrieved memories       & Stateless inference under unbounded context; ICL vs.\ cross-attention fusion unresolved; modality mismatch in retrieval (text vs.\ image vs.\ 3D) \\
%         §\ref{subsec:p3_s7_3_7} Agentic reasoning             & Orchestration policy $\pi$ over the perceive--reason--act--reflect cycle becomes the central learning target; tools (§\ref{subsec:p3_s7_3_5}) and memory (§\ref{subsec:p3_s7_3_6}) composed into closed-loop control & Instruction$+$observation $+$ tools $+$ memory $\to$ sensorimotor action sequence, accumulated skill library, updated world model & Closed-loop error correction; long-horizon procedural-knowledge accumulation; explicit Spatial Theory-of-Mind in multi-agent settings & Sim-to-real fidelity gap; long-horizon spatiotemporal coherence; convergence with §\ref{subsec:p3_s7_3_3}--§\ref{subsec:p3_s7_3_6} blurs the agentic boundary \\
%         §\ref{subsec:p3_s7_3_8} World-model--based reasoning  & Learned generative simulator $\mathbb{P}(s_0,\dots,s_T \mid a_0,\dots,a_{T-1}, c)$ replaces / augments the foundation model as the spatial substrate & Context $c$ ($+$ actions if conditioned) $\to$ inferred spatial properties (understanding) or policies $\pi$ (planning) & Predictive-physics reasoning; inference-time mental simulation; 4D / object-centric structural priors & World-model upper bound (Genie-style long-term memory loss); modality coverage gated by pre-training data; likelihood-vs.-physics gap (likelihood necessary but not sufficient) \\
%         \bottomrule
%     \end{tabularx}
% \end{table}

% Down the table the locus of intervention moves from \textbf{model-as-core} (§\ref{subsec:p3_s7_3_1}--§\ref{subsec:p3_s7_3_2}) through \textbf{model-as-rationale-generator} (§\ref{subsec:p3_s7_3_3}--§\ref{subsec:p3_s7_3_4}) and \textbf{model-as-orchestrator} over external components (§\ref{subsec:p3_s7_3_5}--§\ref{subsec:p3_s7_3_7}) to a \textbf{learned simulator as the reasoning substrate} (§\ref{subsec:p3_s7_3_8}), and the principal learning target migrates with it, from model weights to prompt structure to the orchestration policy $\pi$ to the world model $\mathbb{P}$ itself.
% The paradigms are cumulative rather than parallel: from §\ref{subsec:p3_s7_3_3} onward each adds one dimension (programs, tools, memory, reflection, world models) on top of the previous ones, so the ``agentic'' boundary of §\ref{subsec:p3_s7_3_7} lies where the orchestration policy itself becomes a learning target rather than where any new capability appears, which is an open research question rather than a settled definition.
% The open problem that cuts across all eight is the \emph{precision--grounding asymmetry}: the foundation models that supply semantic competence (§\ref{subsec:p3_s7_3_1}--§\ref{subsec:p3_s7_3_4}) lack the metric precision physical manipulation requires, while the external components that supply precision (§\ref{subsec:p3_s7_3_5}--§\ref{subsec:p3_s7_3_8}) inherit the LLM--tool interface as a new bottleneck.
% That asymmetry is the bridge to §\ref{sec:p3_s7_4}, where foundation-model-driven and symbolic or learning-based components are composed to close the gap from both sides at once.

% \newpage
\section{Hybrid Spatial Reasoning}
\label{sec:p3_s7_4}

\FramedInitial{C}{onventional} approaches to spatial reasoning exhibit fundamental limitations.
Systems that rely solely on linguistic or symbolic reasoning typically lack perceptual grounding in real-world visual and physical structures (\eg, geometry, physics, spatial continuity), making them inadequate for handling cross-view transformations, temporal evolution, or physically coherent scenarios.
Conversely, purely perceptual or vision-based systems, such as those combining visual recognition with rule-based geometric reasoning, often fall short in performing high-level abstract reasoning, multi-step planning, long-term memory integration, or interactive decision-making.
These limitations have motivated a growing body of work advocating for a hybrid paradigm of spatial reasoning: integrating heterogeneous modules with distinct capabilities (perception, world modeling, memory, and action) so as to jointly leverage low-level spatial and physical cues together with high-level cognitive reasoning capabilities.
A perceptual or physics-informed module supplies geometry, relations, affordances, and dynamics; a symbolic or neural module plans, infers, and abstracts over them; memory and history extend the pipeline to multi-step, multi-view, and temporally extended tasks; and because no single modality, task, or dataset carries the whole system, the composition transfers across vision-only, language-only, embodied, and multimodal settings.

% 在空间推理 (spatial reasoning) 的研究中，传统单一范式存在显著局限性。以仅依赖语言/符号推理 (symbolic reasoning) 的系统为例，它们往往缺乏对真实世界视觉／物理结构 (geometry / physics / continuity) 的感知与空间记忆，难以处理跨视角 (cross-view)、跨时 (temporal) 和物理连贯 (physical consistency) 的复杂情境；另一方面，仅依赖感知或视觉模块 (例如视觉识别 + 规则几何推理) 的系统，则可能缺乏高阶抽象推理 (abstract reasoning)、规划 (planning)、长程记忆 (long-term memory) 以及交互式决策 (interactive decision) 的能力。
% 因此，很多工作主张一种 “混合 (hybrid)”空间推理 (spatial reasoning) 的研究范式 — 即通过整合异构 (heterogeneous) 模块 (perception, world-model, reasoning, memory, action) —— 以同时兼顾空间结构感知、几何 / 物理一致性、跨视角／跨模态理解、长期规划与推理。
% 这种混合范式的优点包括但不限于：
% 感知 + 抽象推理：底层视觉 / 物理世界模型 (world-model) 提供空间结构与物体关系 (spatial relationships / affordances / dynamics)，而高层符号 /语言 /逻辑推理 (symbolic or neural reasoning) 提供规划、决策、归纳／演绎等能力；
% 长期记忆与交互能力：通过将 memory / history / context 纳入 world-model + reasoning pipeline，使系统能够处理长期任务 (multi-step, multi-view, temporally-extended)；
% 泛化与灵活性 (generalization & flexibility)：混合模块使系统不依赖单一任务或单一数据集 (\eg pure 3D supervision, or pure NL reasoning)，从而更具拓展性 (scalability) 与适应性 (adaptability)。


\subsection{Task Definition}
\label{subsec:p3-sec7-4-task-definition}
A hybrid spatial reasoning system is a pair $(\{M_i\}_{i=1}^{k}, \pi)$: a set of $k$ heterogeneous modules $M_i \in \{M_\text{percept}, M_\text{world}, M_\text{symb}, M_\text{mem}, M_\text{plan}, M_\text{act}\}$, each contributing a distinct inductive bias (perception, world model, symbolic constraints, memory, planning, or action), and an orchestration function $\pi : \mathcal{S} \to \{1, \dots, k\}$ that selects which module acts on the current system state $s_t \in \mathcal{S}$ (perception, internal map, history). Starting from the input $s_0 = x$, the system evolves as
\begin{equation}
  s_{t+1} = M_{\pi(s_t)}(s_t), \qquad y = \mathrm{Read}(s_T),
  \label{eq:p3_s7_4_hybrid}
\end{equation}
and the three composition styles found in practice are special cases of Eq.~\eqref{eq:p3_s7_4_hybrid}: a \emph{pipeline} fixes $\pi$ to a constant module order, an \emph{iterative} system lets $\pi$ revisit modules until a verifier accepts, and an \emph{end-to-end} system makes every $M_i$ differentiable so that the whole trajectory is trained jointly. Hybrid spatial reasoning is the paradigm in which $\pi$ and the module set $\{M_i\}$ are designed together to minimize the precision-grounding asymmetry: the foundation models of §\ref{sec:p3_s7_3} supply semantic competence, and the symbolic and learning-based components of §\ref{sec:p3_s7_1}--§\ref{sec:p3_s7_2} supply metric precision and verifiable structure. Existing systems therefore differ along three axes: \emph{which} modules they compose, \emph{how} they compose them (pipeline, iterative, or end-to-end), and how the orchestration function is realized (hand-coded rules or a learned policy).

\subsection{Methodology}
\label{subsec:p3-sec7-4-methodology}

Five hybridization patterns recur in the literature, distinguished by which module the foundation model is paired with and by what supplies the truth signal.

\begin{figure}
    \centering
    \includegraphics[width=\linewidth]{figures/part3/section7/foundation_rl_hybrid.pdf}
    \caption{The foundation-model + reinforcement-learning hybrid pattern: an LLM/VLM hypothesis generator coupled with an RL controller and with a physics engine or symbolic solver for verification. The RL controller's policy (written $\pi$ in the figure) is $\pi_{\mathrm{RL}}$ in the text; the bottom bar is the orchestration function $\pi$ of Eq.~\eqref{eq:p3_s7_4_hybrid}.}
    \label{fig:p3_s7_4_foundation_rl_hybrid}
\end{figure}

\paragraph{Foundation Models with Reinforcement Learning.}
A prominent hybrid couples a foundation model's semantic competence with reinforcement learning's long-horizon decision-making, often closed by a verifier that supplies executable structure (Figure~\ref{fig:p3_s7_4_foundation_rl_hybrid}).
The foundation model proposes hypotheses, sub-goals, or candidate action sequences; an RL policy $\pi_{\mathrm{RL}}$ commits to and executes them over extended horizons; and the resulting reward, optionally augmented by a physics engine or symbolic checker, updates both the policy and the orchestration.
Embodied-R~\cite{zhao2025embodied} instantiates this division of labor: a large vision--language model handles perception, while a small language model trained by RL, with a reward that scores the logical consistency between its reasoning trace and its answer, performs the slow, deliberative spatial reasoning, evaluated on UrbanVideo-Bench and VSI-Bench.
The group-relative fine-tuning recipes of §\ref{subsec:p3_s7_2_5} (Reason-RFT, OctoNav, SpaceR~\cite{tan2025reasonrft,gao2026octonav,ouyang2025spacer}) are the same hybrid with the verifier folded into a verifiable reward and are not repeated here.
The foundation model contributes broad priors and RL supplies credit assignment and exploration; the cost is the sample inefficiency and reward-specification difficulty inherited from RL (§\ref{subsec:p3_s7_2_5}).

\paragraph{Foundation Models with Explicit Physical Simulation.}
A second pattern pairs the foundation model with a physics engine, so that the truth signal is a simulated trajectory rather than a learned score; the engine side is the physical reasoning of §\ref{subsec:p3_s7_1_2}.
The LLM proposes hypotheses, structural descriptions, or programs, the simulator executes them under conservation laws, kinematic constraints, contact dynamics, and reachability limits, and the discrepancy between simulation and observation is fed back for another round.
LLMPhy~\cite{cherian2024llmphy} is the clearest instance: the LLM writes Python programs that embed estimates of latent physical parameters (friction, stiffness, mass) and of the scene layout, a physics engine executes them, and the reconstruction error against the observed video is returned as feedback, so that building the digital twin becomes a black-box optimization driven by the LLM.
DiffGen~\cite{jin2024diffgen} makes the loop differentiable, generating robot demonstrations by descending gradients through a differentiable simulator, a differentiable renderer, and a vision--language model until the rendered outcome matches the language instruction in embedding space.
PhysVLM~\cite{zhou2025physvlm} moves the physics into the input instead: a robot's reachable workspace is abstracted into a robot-agnostic space-physical reachability map that the VLM consumes alongside the image, trained on Phys100K and evaluated on the EQA-phys benchmark.
At the far end, ChronoLLM~\cite{wang2025chronollm} fine-tunes LLMs to write PyChrono multibody-simulation scripts, from a single pendulum to a full vehicle on deformable terrain; here the model authors the simulation rather than reasoning inside it, which places it at the boundary with the code-generation paradigm of §\ref{subsec:p3_s7_3_4}.
What the simulator contributes is an inductive bias that stays valid for unseen objects and configurations, so the hybrid extrapolates where a purely statistical model would not.

\paragraph{Foundation Models with Symbolic Reasoning.}
A third pattern inserts an explicit symbolic layer (logical constraints, relational predicates, or formal planning specifications) between natural-language reasoning and executable spatial inference, replacing the hand-crafted parsers of §\ref{sec:p3_s7_1} with foundation-model-generated symbolic structure.
Two sub-patterns differ in when the symbolic machinery is used.
In the first it is called at inference time: NS3D~\cite{hsu2023ns3d} has an LLM translate a 3D referring expression into a first-order logic program that a neural executor grounds against the scene's objects and relations on ReferIt3D, and LEFT~\cite{hsu2023s} (§\ref{subsec:p3_s7_3_4}) makes the same programs differentiable so that concept grounding is learned rather than hand-specified.
In the second, symbolic knowledge enters only at training time as a constraint on a neural model: \citet{premsri2024neurosymbolic} train language models for reasoning over spatial language with spatial logic rules as auxiliary losses, and RelateSeg~\cite{lin2025neuro} extracts relations such as $\langle\text{cat}, \text{to-right-of}, \text{person}\rangle$ from open-vocabulary segmentation outputs, writes them as first-order logic, and relaxes them with fuzzy logic into a single auxiliary loss that improves mIoU on four benchmarks without adding parameters.
The Scene Language~\cite{zhang2025scenelanguage} (§\ref{subsec:p3_s8_2_5}) is the representation such hybrids increasingly target: programs for structure, words for semantics, embeddings for appearance.
The symbolic layer buys consistency and compositional guarantees; its price is a closed predicate vocabulary that the neural side constantly exceeds, a point taken up under Challenges and Trends below.

\paragraph{Foundation Models with Spatial Memory and World Models.}
A fourth pattern gives the foundation model a persistent spatial state that it does not have to carry in its context window: an explicit memory (episodic trajectories, topological maps, 3D layouts) or a learned world model that predicts how the scene changes under action.
The hybrid question is the interface between the two.
MSNav~\cite{liu2025msnav} splits zero-shot vision-and-language navigation into a memory module, a spatial module, and a decision module, with a small fine-tuned LLM extracting the landmarks and relations the spatial module needs, and reaches state-of-the-art zero-shot results on R2R and REVERIE.
MindExplore~\cite{li2025towardsvla} stacks a reasoning layer (task-specific chains of thought that emit meta-actions) over an acting layer (a mixture of policy experts with a multimodal diffusion policy) and uses a memory mechanism as the feedback channel between the two, running at 30~FPS in dynamic sand environments.
The 3D spatial memories of §\ref{subsec:p3_s7_3_6} (3DSPMR~\cite{cai2025vision}, Embodied-RAG~\cite{xie2024embodiedrag}, Meta-Memory~\cite{mao2025meta}) and the world-model coupling of §\ref{subsec:p3_s7_3_8} are the memory- and simulator-side components these hybrids draw on; MindJourney's test-time gain on SAT~\cite{yang2026mindjourney} (about eight points of top-1 accuracy from imagined viewpoints alone, with the VLM untouched) is the cleanest quantitative evidence that a world model used as a tool closes part of the gap.
In each case the foundation model keeps high-level inference, instruction following, and semantic grounding, and offloads persistence, geometry, and dynamics to the dedicated module.

\begin{figure}[!ht]
    \centering
    \includegraphics[width=0.8\linewidth]{figures/part3/section7/agentic_hybrid_orchestration.pdf}
    \caption{Illustration of the agentic hybrid orchestration pattern: a foundation-model controller dynamically selects perception, world-modeling, symbolic-reasoning, memory, and action tools within a closed-loop perceive--reason--act--reflect cycle.}
    \label{fig:p3_s7_4_agentic_hybrid}
\end{figure}

\paragraph{Agentic Orchestration of Hybrid Components.}
The fifth pattern makes the orchestration function itself the object of design: a foundation model acts as a deliberative controller that decides, at each step, which of the perception, world-model, symbolic, memory, and action modules to invoke, within a closed perceive--reason--act--reflect loop (Figure~\ref{fig:p3_s7_4_agentic_hybrid}).
It differs from the agentic reasoning of §\ref{subsec:p3_s7_3_7} in what is being scheduled: there, tools and memory; here, additionally symbolic engines and physics simulators, with the scheduler $\pi$ of Eq.~\eqref{eq:p3_s7_4_hybrid} itself a candidate for learning.
VADAR~\cite{marsili2025visual} (§\ref{subsec:p3_s7_3_4}) is the program-synthesis version of such a controller, generating the API it then calls.
``See it. Say it. Sorted''~\cite{zhang2025seeitsayit} is a training-free version for sketch-to-diagram generation: a critic VLM proposes a few qualitative relational edits, several LLMs synthesize candidate SVG programs, and a judge VLM selects among them, which keeps global constraints such as alignment and connectivity intact where numerical estimation would drift.
\citet{manh2025mindmeetspace} arrive at the same module inventory from neuroscience, arguing that such a controller should sit over multimodal sensing, multisensory integration, egocentric--allocentric conversion, a cognitive map, spatial memory, and spatial reasoning.
The gain is compositionality and long-horizon consistency that no static architecture provides; the cost is the latency and attribution problems discussed below.

% 之后
% 具体地，今年来，研究者越来越倾向于将大型基础模型 (foundation models) 与强化学习 (reinforcement learning, RL) 结合，以构建具备 spatial reasoning 能力的混合系统。
% 例如，最新工作 Embodied-R: Collaborative Framework for Activating Embodied Spatial Reasoning in Foundation Models via Reinforcement Learning (简称 Embodied-R) 就是一个典型代表。该工作将大型视觉—语言模型 (VLM) 用于感知 (visual observation) —— 提取视觉特征与空间线索；同时用较小的语言模型 (LM) 作为推理 /决策模块 (reasoner / planner)，并通过强化学习训练 (设计 reward 考量 think-answer 逻辑一致性) 实现 “slow-thinking” (慢思考) 的空间推理。
% 这些工作表明：将基础模型 (VLM / LLM) 与 world-model + RL 结合，是当前实现混合空间推理的最有前景路径之一 — 尤其适合应对复杂、动态、物理 /几何约束强的现实场景 (例如 3D 布局、具身导航、物体交互、场景理解等)。
% 此外，从更理论／宏观层面来看，有分析指出当前基础模型 (foundation models) 在决策 (decision-making) 与空间 /世界智能 (world-model) 方面存在结构性空缺。因此，通过 RL、world-model、模块化 /混合架构 (hybrid architecture) 来弥补这一空缺，有助于未来构建真正具有空间智能 (spatial intelligence) 的通用系统 (general-purpose agents / embodied AI / spatially-aware AI)。

%   通过将混合架构纳入智能体 (agent) 或规划 + 行动 (planning + action) 回路 (loop) 中，我们可以构造更为通用、灵活、可扩展 (scalable) 的空间智能系统。这种 “Agentic Hybrid Orchestration” 的设计核心包括以下几个关键维度：
% LLM / VLM 作为 Orchestrator (协调者／指挥者)：高层语言 /符号 /规划模块 (LLM) 负责决定 “何时感知 (perceive)”、 “如何构建或更新 world-model (scene graph / layout / spatial representation)”、 “何时进行推理 /规划 (plan)” 以及 “何时发起行动 /布局 /生成 / manipulation”。
% 工具选择 (tool selection) 与顺序 (sequencing)：系统可以预定义或动态维护一组 “工具 (tools)” —— 比如视觉感知模块、world-model builder、符号 /语言推理模块、布局 /渲染 模块 —— LLM 负责根据当前任务与上下文选择合适工具并安排调用顺序。
% 反馈驱动 (feedback-driven) 的迭代 (iterative) 推理 (reasoning)：通过 RL (或其他优化方式) 将最终目标 (task success, spatial consistency, physical plausibility, layout realism / utility) 的反馈 (reward / loss) 传播回 orchestrator 和各工具调用决策。
% 这种 Agentic Hybrid Orchestration 的好处是，它将混合架构系统转化为一个动态、交互式、多轮 (multi-step) 的代理 (agent) 回路 (loop)，从而具备长期规划、空间结构一致性、物理可行性、泛化 与适应性 — 正是混合空间推理 (Hybrid Spatial Reasoning) 的目标。

\subsection{Benchmarks}
\label{subsec:p3-sec7-4-benchmarks}
Hybrid systems are evaluated end to end on the benchmark of their target task rather than module by module, and the benchmark is often one the paper introduces.
Among the systems above, Embodied-R reports on UrbanVideo-Bench and VSI-Bench, MindJourney on SAT, NS3D on ReferIt3D, MSNav on R2R and REVERIE, and PhysVLM on RoboVQA and OpenEQA alongside its own EQA-phys; LLMPhy ships TraySim, CLEVRER-LLMPhy, and Real-TraySim, MindExplore ships SandGo-1k and SandThink-21k, and DrawingBench~\cite{kim2025drawingbench} is built specifically for agentic hybrids, scoring an LLM's sequence of low-level mouse actions against eight rule-based criteria for the spatial configuration it was asked to draw.
Physics-side components are still checked against PHYRE and ShapeStacks (§\ref{subsec:p3_s7_1_2}) and symbolic components against the constraint- and rule-based suites of §\ref{subsec:p3_s7_1_4}--§\ref{subsec:p3_s7_1_5}, but no current benchmark attributes an end-to-end score to the modules that produced it; that gap is taken up under Challenges and Trends below and in the cross-paradigm survey of §\ref{sec:p3_s7_5} (Table~\ref{tab:benchmark_comparison}, whose 15-row agentic subset is Table~\ref{tab:spatial_reasoning_benchmark} of §\ref{subsec:p3_s7_3_7}).

\subsection{Challenges and Trends}
\label{subsec:p3-sec7-4-challenges-and-trends}
Hybrid spatial reasoning has produced architectures that combine complementary inductive biases in principled ways, but several recurring limitations originate in the module-composition framework itself rather than in any individual module.

\textbf{Frame, Scale, and Geometric Consistency.}
Unlike generic modular systems, spatial hybrids must reconcile heterogeneous \emph{spatial} representations across modules: perception emits camera or egocentric frames, map and memory modules operate in world or allocentric coordinates, symbolic reasoners abstract to topological relations, and physics engines expect metric units. Misaligned reference frames, scales, or conventions silently corrupt downstream reasoning, and because each module reports confidence in its own terms, small upstream pose, depth, or scale errors propagate and compound across the pipeline without a coherent notion of uncertainty. Maintaining frame, scale, and geometric consistency is thus a first-order obstacle with no analogue in non-spatial hybrid systems.

\textbf{Interfaces, Latency, and Attribution.}
The strength of a hybrid system is bounded by its worst module interface. Symbolic-to-neural translation is particularly brittle when the symbolic side uses closed-vocabulary predicates while the neural side proposes open-vocabulary ones, and the interface design becomes a hidden first-class hyperparameter that current literature seldom isolates. Each module invocation also adds wall-clock latency that compounds across a perceive--reason--act--reflect loop, so high-frequency manipulation and navigation may not be reachable through orchestration alone. And when a hybrid fails, attributing the failure to a specific module is non-trivial; this complicates iteration and is often cited as a practical reason why end-to-end models can outperform hybrid stacks on aggregate benchmarks despite offering weaker per-module guarantees.

The orchestration function $\pi$ is migrating from hand-coded rules to \emph{learnable orchestration policies} trained with reinforcement or imitation learning, which blurs the boundary with §\ref{subsec:p3_s7_3_7}.
\emph{Differentiable hybrid stacks}, built from differentiable physics simulators, differentiable optimization layers (§\ref{subsec:p3_s7_2_4}), and differentiable retrieval, make end-to-end training of a hybrid viable and remove the historical interface penalty, offering a path through the precision--grounding asymmetry that motivates this section.
Evaluation is changing too, toward the \emph{holistic and diagnostic}: in response to the attribution problem, emerging practice pairs aggregate task metrics with per-module diagnostic probes and counterfactual module-swap ablations that isolate where a hybrid system succeeds or fails, a direction consolidated in the cross-paradigm benchmark survey of §\ref{sec:p3_s7_5}.

% \newpage
z\section{Spatial Reasoning Benchmarks}
\label{sec:p3_s7_5}

\FramedInitial{E}{arly} spatial-reasoning benchmarks positioned the model as a detached observer. CLEVR~\cite{johnson2017clevr} and VSR~\cite{liu2023visual} ask for relative geometric relations from a static third-person viewpoint, and their 2D inputs leave depth ambiguous, a weakness that 3DSRBench~\cite{ma20253dsrbench}, MMSI-Bench~\cite{yang2025mmsi}, DSI-Bench~\cite{zhang2025dsi}, and STI-Bench~\cite{li2025sti} now probe directly.
SQA3D~\cite{ma2022sqa3d}, OpenEQA~\cite{majumdar2024openeqa}, and EmbSpatial-Bench~\cite{du2024embspatial} instead place the model inside the scene, so that it must self-localize before it can answer, and ScanQA~\cite{azuma2022scanqa}, ScanRefer~\cite{chen2020scanrefer}, and Multi3DRefer~\cite{zhang2023multi3drefer} ground queries in point clouds and RGB-D data, where the answer is a coordinate rather than a relation.
The reasoning demanded has deepened in parallel: IntPhys~\cite{riochet2018intphys} and Physion~\cite{Bear2021Physion} test whether a model expects gravity, collision, and friction to hold over time, and OpenEQA and RealWorldQA~\cite{xai2024grok15v} ask functional questions (``Can I navigate through this gap?'', ``Where would I find a cloth to clean a spill?'') in which geometry matters only through what it affords.

Data provenance splits the field. Simulators, the sources behind SpatialVLM~\cite{chen2024spatialvlm}, AI2-THOR~\cite{kolve2017ai2thor}, and Habitat~\cite{szot2021habitat2}, yield unlimited, perfectly annotated data; real-world collections such as ScanNet~\cite{dai2017scannet} and OpenEQA supply the clutter, occlusion, and sensor noise that simulation does not reproduce.
The resulting practice is sim-to-real: pre-train on large synthetic corpora for general spatial logic, then fine-tune on the scarcer real data.
Annotated 3D data remain scarce and fine-grained metric estimation remains hard for large models; both limits recur throughout the table below.

% \wsq{总结现有的有关spatial reasoning 相关的benchmarks}


\subsection{Comparative Analysis}
\label{subsec:p3-sec7-5-comparative-analysis}

Table~\ref{tab:benchmark_comparison} compiles 61 spatial-reasoning benchmarks across four categories (physical-focused, rooted in §\ref{subsec:p3_s7_1_2}; vision-based, overlapping with §\ref{subsec:p3_s7_3_1}; 3D-focused, rooted in §\ref{subsec:p3_s7_3_6} and §\ref{subsec:p3_s7_3_8}; and robot-environment, rooted in §\ref{subsec:p3_s7_3_7} and §\ref{sec:p3_s7_4}).
A narrower 15-row subset focused specifically on agentic spatial reasoning is provided in Table~\ref{tab:spatial_reasoning_benchmark} of §\ref{subsec:p3_s7_3_7}; this subset filters Table~\ref{tab:benchmark_comparison} to image-and-video benchmarks that exercise agentic deployments.

\begin{table*}[!t]
  \caption{Spatial-reasoning benchmarks surveyed in this chapter, grouped by what they test. \texttt{Number} is the evaluation-set size where the source paper documents one (train / test when both are given, scene counts for simulators, $\sim$ for approximate figures). \texttt{Perspective} marks explicitly designed ego (first-person) and exo (third-person) settings; \textit{Undeclared} marks attributes the source paper does not state. \texttt{Data Source} and \texttt{Task Type} follow the source papers.}
  \label{tab:benchmark_comparison}
  \centering
  \fontsize{9}{10}\selectfont
  \setlength{\tabcolsep}{1.4mm}
  \resizebox{\linewidth}{!}{
  \begin{tabular}{@{}lcccccccc}
    \toprule
      \multirow{2}{*}{\textbf{Benchmark}}
      & \multirow{2}{*}{\textbf{Number}}
      & \multirow{2}{*}{\textbf{Modality}}
      & \multicolumn{2}{c}{\textbf{View Type}}
      & \multicolumn{2}{c}{\textbf{Perspective}}
      & \multirow{2}{*}{\textbf{Data Source}}
      & \multirow{2}{*}{\textbf{Task Type}} \\
      \cmidrule(lr){4-5}
        \cmidrule(lr){6-7}
        & &
        & \textbf{Single} & \textbf{Multi}
        & \textbf{Ego} & \textbf{Exo}
        & \\
    \midrule
    \rowcolor{lightgrey}\multicolumn{9}{c}{$\bullet$ \textit{Physical Focused Spatial Reasoning}} \\
    {IntPhys}~\cite{riochet2018intphys} & 15k / 4.3k & 3D   & \checkmark & \crossmark    & \crossmark & \checkmark   & Synthetic    & VoE plausibility\\
    {IntPhys 2}~\cite{bordes2025intphys2}& 1,416    & 3D   & \checkmark & \crossmark    & \crossmark & \checkmark       & Synthetic    & VoE plausibility\\
    {Physion}~\cite{Bear2021Physion}  & $\sim$1.2k  & 3D    & \checkmark & \crossmark  & \crossmark & \checkmark   & Synthetic & Physical prediction / choice \\
    {Physion$++$}~\cite{tung2023physion} & $\sim$10k  & 3D  & \checkmark & \crossmark   & \crossmark & \checkmark  & Synthetic  & Physical prediction / choice\\
    {Physical Concepts Dataset}~\cite{piloto2020learn}  & 300k    & 3D   & \checkmark & \crossmark   & \crossmark & \checkmark  & Synthetic   & classification\\
    {Physics-IQ}~\cite{motamed2025generative} & 396  & Video     & \checkmark & \crossmark   & \crossmark & \checkmark  & Real   & Generative physical prediction \\
    {SPhyR}~\cite{siedler2025sphyr} & 1,296 & Image & \checkmark & \crossmark & \multicolumn{2}{c}{\cellcolor{lightgray}\textit{Undeclared}} & Synthetic & Physical prediction / infill \\
    \hline
    \rowcolor{lightgrey}\multicolumn{9}{c}{$\bullet$ \textit{Vision-Based Spatial Reasoning}} \\
    CLEVR \citep{johnson2017clevr} & 853k & Image & \checkmark & \crossmark & \crossmark & \checkmark & Simulation & VQA \\
    VSR \citep{liu2023visual}& 10,972 & Image & \checkmark & \crossmark & \crossmark & \checkmark & Real & VQA \\
    What's Up \citep{KamathHC23a}  & $\sim$8,000 & Image & \checkmark & \crossmark & \crossmark & \checkmark & Real & VQA\\
    RocketScience~\citep{hoehing2025RocketScience} & 482 & Image & \checkmark & \crossmark  & \crossmark & \checkmark & Real & VQA \\
    SpaRE \citep{abs-2504-20648} & 3.4M & Image &  \checkmark & \crossmark & \crossmark & \checkmark & Real/Sim & VQA \\
    RealWorldQA~\cite{xai2024grok15v} & 765 & Image & \checkmark & \crossmark & \checkmark & \crossmark & Real & VQA \\
    3DSRBench~\cite{ma20253dsrbench} & 2,772 & Image & \checkmark  & \crossmark & \crossmark  &  \checkmark & Real/Sim & VQA\\
    InternSpatial-Bench~\cite{deng2025internspatial} &  6,008 & Image  & \checkmark  & \crossmark & \crossmark  &  \checkmark  & Real & VQA \\
    OmniSpatial~\cite{jia2025omnispatial} & $\sim$8,400 & Image & \checkmark & \crossmark & \checkmark & \checkmark   &  Real & VQA\\
    Spatial457~\cite{wang2025spatial457} & 23,752 & Image & \checkmark & \crossmark & \multicolumn{2}{c}{\cellcolor{lightgray}\textit{Undeclared}}  & Real & VQA\\
    All-Angles Bench~\cite{yeh2025seeing} & $\sim$2,100 & Image  & \crossmark & \checkmark & \multicolumn{2}{c}{\cellcolor{lightgray}\textit{Undeclared}} & Real  & VQA \\
    EmbSpatial-Bench~\cite{du2024embspatial} & 3,640 & Image & \checkmark  & \crossmark & \checkmark & \crossmark &  Real & VQA \\
    MMSI-Bench~\cite{yang2025mmsi} & 1,000 & Image & \crossmark & \checkmark  & \multicolumn{2}{c}{\cellcolor{lightgray}\textit{Undeclared}} &  Real & VQA\\
    MindCube~\cite{yin2025spatial} & 21,154 & Image & \crossmark & \checkmark & \multicolumn{2}{c}{\cellcolor{lightgray}\textit{Undeclared}}  & Simulation  & VQA \\
    BLINK~\cite{fu2024blink} & 3,807 & Image & \checkmark & \crossmark & \crossmark & \checkmark & Real & VQA \\
    NuScenes-SpatialQA~\cite{tian2025nuscenesspatialqa} & 3.3M+ & Video & \checkmark  & \checkmark  & \checkmark  & \checkmark  & Real & VQA \\
    VSI-Bench~\cite{yang2025vsibench} & 5,130 & Video & \checkmark & \crossmark & \checkmark & \crossmark & Real  & VQA\\
    VLM4D~\cite{zhou2025vlm4d} & $\sim$1,800 & Video & \checkmark & \crossmark &  \checkmark & \checkmark &  Real & VQA \\
    STI-Bench~\cite{li2025sti} & 2,064 & Video  & \checkmark & \crossmark & \checkmark & \checkmark & Real & VQA\\
    DSI-Bench~\cite{zhang2025dsi} & $\sim$1,700 & Video  & \checkmark & \crossmark & \multicolumn{2}{c}{\cellcolor{lightgray}\textit{Undeclared}} & Real & VQA\\
    SITE~\cite{wang2025site} & 8,068  & Video & \checkmark & \crossmark & \multicolumn{2}{c}{\cellcolor{lightgray}\textit{Undeclared}} & Real/Sim & VQA \\
    RoboSpatial \citep{SongBTTSB25} & 3M & Image/3D & \checkmark & \crossmark & \checkmark & \crossmark & Real/Sim & VQA \\
    OSD \citep{cheng2024spatialrgpt} & \textit{Undeclared} & Image & \checkmark & \crossmark & \crossmark & \checkmark & Real & VQA \\
    SpatialVLM \citep{chen2024spatialvlm} & 2B & Image & \checkmark & \crossmark & \crossmark & \checkmark & Real/Sim & VQA \\
    ProVision \citep{abs-2412-07012} & 65k & Image & \checkmark & \crossmark & \crossmark & \checkmark & Simulation & VQA \\
    WhatIsWhere \citep{LiangZZZ16} & 6,000 & Image & \checkmark & \crossmark & \crossmark & \checkmark & Real & VQA \\
    SURDS~\cite{guo2024surds} & 9,250 & Image & \checkmark & \crossmark & \checkmark & \crossmark & Real & VQA \\
    EarthSpatialBench~\cite{xu2026earthspatial} & 325k+ & Image & \checkmark & \crossmark & \crossmark & \checkmark & Real & VQA \\
    PureSpace~\cite{li2026purespace} & 1,200 & Image & \checkmark & \crossmark & \multicolumn{2}{c}{\cellcolor{lightgray}\textit{Undeclared}} & Simulation & VQA \\
    SAT~\cite{ray2024sat} & 175k / 150 & Image & \crossmark & \checkmark & \checkmark & \checkmark & Simulation & VQA \\
    SpatiaLQA~\cite{xie2026spatialqa} & 9,605 & Image & \checkmark & \crossmark & \crossmark & \checkmark & Real & VQA \\
    Q-Spatial++~\cite{wang2025svqa} & 101 & Image & \checkmark & \crossmark & \crossmark & \checkmark & Real & VQA \\
    \hline
    \rowcolor{lightgrey}\multicolumn{9}{c}{$\bullet$ \textit{3D Focused Spatial Reasoning}} \\
    ScanQA~\cite{azuma2022scanqa} & 41,363 & 3D & \crossmark & \checkmark & \crossmark & \checkmark & Real (ScanNet) & 3D VQA \\
    SQA3D~\cite{ma2022sqa3d} & 33,460 & 3D & \checkmark & \crossmark & \checkmark & \crossmark & Real (ScanNet) & 3D VQA \\
    MMScan QA~\cite{lyu2024mmscan} &  295K & 3D+Image & \checkmark & \checkmark & \crossmark & \checkmark & Real (EmbodiedScan) & 3D VQA \\
    OpenEQA~\cite{majumdar2024openeqa} & 1,632 & Video/3D & \checkmark & \crossmark & \checkmark & \crossmark & Real (HM3D, ScanNet) & Embodied QA (EM / A splits) \\
    Space3D-Bench~\cite{szymanska2024space3dbench} & 1,000 & 3D & \checkmark & \crossmark & \checkmark & \crossmark  & Real(Replica) & 3D QA \\
    Scan2Cap~\cite{chen2021scan2cap} & 39,775 & 3D & \crossmark & \checkmark & \crossmark & \checkmark & Real (ScanNet) & 3D Captioning \\
    MMScan Captioning~\cite{lyu2024mmscan} & 172k & 3D+Image & \checkmark & \checkmark & \crossmark & \checkmark & Real (EmbodiedScan) & 3D Captioning \\
    ScanRefer~\cite{chen2020scanrefer} & 51,583 & 3D & \crossmark & \checkmark & \crossmark & \checkmark & Real (ScanNet) & 3D Grounding \\
    Multi3DRefer~\cite{zhang2023multi3drefer} & 42,735 & 3D & \crossmark & \checkmark & \crossmark & \checkmark & Real (ScanNet) & 3D Grounding \\
    SR3D~\cite{achlioptas2020referit3d} & 83,572 & 3D & \crossmark & \checkmark & \crossmark & \checkmark & Real (ScanNet) & 3D Grounding \\
    NR3D~\cite{achlioptas2020referit3d} & 41,503 & 3D & \crossmark & \checkmark & \crossmark & \checkmark & Real (ScanNet) & 3D Grounding \\
    SpaceEra++~\cite{guan2026spaceera} & 925k & Video/3D & \crossmark & \checkmark & \checkmark & \crossmark & Synthetic & 3D VQA \\
    Ego3D-Bench~\cite{gholami2025ego3d} & 8,600+ & Image & \crossmark & \checkmark & \checkmark & \crossmark & Real & VQA \\
    SpatialScore~\cite{wu2025spatialscore} & 28,093 (Hard: 1,400) & Image/3D & \crossmark & \checkmark & \multicolumn{2}{c}{\cellcolor{lightgray}\textit{Undeclared}} & Real/Sim & VQA \\
    iVISPAR~\cite{mayer2025ivispar} & 300 & 3D Env & \checkmark & \crossmark & \crossmark & \checkmark & Simulation & Interactive reasoning \\
    \hline
    \rowcolor{lightgrey}\multicolumn{9}{c}{$\bullet$ \textit{Robot}} \\
    AI2-THOR~\cite{kolve2017ai2thor} & 120 (Scenes) & 3D Env & \crossmark & \checkmark & \checkmark & \crossmark & Simulation & Nav/Manip \\
    TDW~\cite{gan2020threedworld} & \textit{Infinite} & 3D Env & \crossmark & \checkmark & \checkmark & \checkmark & Simulation & Physical Sim \\
    Habitat 2.0~\cite{szot2021habitat2} & 100 (Scenes) & 3D Env & \crossmark & \checkmark & \checkmark & \crossmark & Simulation & Navigation \\
    iGibson 2.0~\cite{Li2021iGibson2} & 15 (Scenes) & 3D Env & \crossmark & \checkmark & \checkmark & \crossmark & Simulation & Interactive Nav \\
    Habitat 3.0~\cite{puig2023habitat3} & 211 (Scenes) & 3D Env & \crossmark & \checkmark & \checkmark & \crossmark & Simulation & Social Nav \\
    R2R~\cite{anderson2018vision} & 21,567 & Image/Path & \crossmark & \checkmark & \checkmark & \crossmark & Real & Navigation \\
    RxR~\cite{ku2020roomacrossroom} & 126,069 & Image/Path & \crossmark & \checkmark & \checkmark & \crossmark & Real & Navigation \\
    \bottomrule
  \end{tabular}
  }
\end{table*}

\emph{A modality migration} runs from the top of the table to the bottom: the physical-focused benchmarks rely on synthetic 3D renderings, the vision-based block is image- and video-heavy with growing 3D extensions, the 3D-focused block is built on point clouds and RGB-D scans, and the robot block has moved largely to 3D environments.
\emph{The Perspective axis is imbalanced}: a large fraction of vision-based benchmarks carry \textit{Undeclared} ego/exo annotations, the physical-focused benchmarks lean exo (a third-person observer judging physical plausibility), and the robot block is entirely ego; the field is migrating from spectator-style observation to situated cognition faster than benchmark annotations record it.
\emph{The Data-Source spectrum is bifurcated} between predominantly synthetic or simulated (top block and robot block) and predominantly real (middle blocks), with the Real/Sim benchmarks (SpaRE, 3DSRBench, RoboSpatial, SpatialVLM, SITE) as the bridge that the sim-to-real training practice of the opening relies on.

\subsection{Future Directions}
\label{subsec:p3-sec7-5-future-directions}

\paragraph{Unified Spatial-Cognition Benchmark.}
No single benchmark in Table~\ref{tab:benchmark_comparison} exercises a model along all of perception fidelity, symbolic precision, memory persistence, and long-horizon action conditioning simultaneously. 
A unified benchmark covering all four §\ref{sec:reasoning} paradigms (§\ref{sec:p3_s7_1} rule-based, §\ref{sec:p3_s7_2} learning-based, §\ref{sec:p3_s7_3} foundation-model-driven, §\ref{sec:p3_s7_4} hybrid) would close the evaluation-attribution challenge flagged under Challenges and Trends in §\ref{sec:p3_s7_4}.

\paragraph{Cross-Modal Evaluation.}
The current axes (Modality, View, Perspective, Data Source, Task Type) under-specify the cross-modal demand of contemporary spatial reasoning. 
Future benchmarks should explicitly probe modality-to-modality transfer (\eg, 2D-trained models evaluated on 3D queries; text-conditioned policies evaluated on video-grounded tasks), so that the modality-mismatch challenge flagged in §\ref{subsec:p3_s7_3_6} (text-vs.-image-vs.-3D retrieval) becomes a first-class evaluation axis rather than a per-paper case study.

\paragraph{Embodied Long-Horizon Coverage.}
The Robot block in Table~\ref{tab:benchmark_comparison} covers navigation thoroughly but under-represents long-horizon manipulation and multi-room interactive tasks. 
The §\ref{subsec:p3_s7_3_7} 15-row subset (Table~\ref{tab:spatial_reasoning_benchmark}) further restricts attention to image-and-video benchmarks; the embodied long-horizon suites ALFWorld and VirtualHome (§\ref{subsec:p3_s7_3_4}) and cross-embodiment evaluation on Open X-Embodiment (§\ref{subsec:p3_s7_3_2}) are the nearest existing coverage, and standardizing on them would close the long-horizon-spatiotemporal-coherence challenge flagged under Challenges and Trends in §\ref{subsec:p3_s7_3_7}.

\paragraph{From Descriptive to Interactive Spatial Reasoning.}
The Task Type column of Table~\ref{tab:benchmark_comparison} is dominated by visual question answering: the overwhelming majority of benchmarks ask a model to \emph{describe} static spatial relations (``where is the object''), and only the small Robot block evaluates reasoning that is conditioned on the agent's own actions.
Yet the trajectory traced in this chapter, from passive observation to situated interaction, calls for benchmarks that test \emph{interactive} spatial reasoning: predicting how an action changes the spatial state, judging affordances and reachability (``can I fit, reach, or grasp this''), and reasoning counterfactually about manipulation outcomes.
Such evaluation couples the action-conditioned reasoning of §\ref{subsec:p3_s7_2_5} with the embodied hybrid pipelines of §\ref{sec:p3_s7_4}, extending the descriptive benchmarks surveyed here toward the manipulation- and affordance-grounded settings that embodied deployment ultimately demands.



\section{Summary and Outlook}
\label{sec:p3_s7_6}

This chapter followed spatial reasoning along the spectrum on which it developed: inference over explicit geometric, physical, topological, logical, and causal models (§\ref{sec:p3_s7_1}); reasoning acquired from data (§\ref{sec:p3_s7_2}); reasoning elicited from large pre-trained models (§\ref{sec:p3_s7_3}); hybrids that combine these components (§\ref{sec:p3_s7_4}); and the benchmarks that measure them (§\ref{sec:p3_s7_5}).

\textbf{Precision and grounding sit at opposite ends of the spectrum.} Rule-based methods are inspectable and can certify feasibility or safety, but they assume exact, fully specified inputs and a closed vocabulary of predicates, and they break down where perception is noisy or a task falls outside the authored rules. Learning-based methods acquire spatial competence from data, yet often reach high benchmark accuracy through dataset-specific shortcuts that do not survive novel compositions (§\ref{subsec:p3_s7_2_1}), and they need far more interaction than classical planners. Foundation models bring broad semantic competence, but their visual encoders retain only coarse geometry, textual reasoning mistracks coordinates over multi-step chains (§\ref{subsec:p3_s7_3_3}), and vision-language-action models often interpret an instruction correctly and execute it incorrectly (§\ref{subsec:p3_s7_3_2}). Hybrid systems aim to pair the semantic reach of the last family with the metric precision and verifiability of the first two.

\textbf{The failures concentrate at interfaces.} Across tool-augmented, memory-augmented, and hybrid reasoning, the recurring limitations arise less within any single component than where components meet: in forming a query that a memory or tool can answer, in translating between closed and open vocabularies, and in keeping coordinate frames and scales consistent across modules (§\ref{sec:p3_s7_4}).

\textbf{Evaluation lags the shift toward situated reasoning.} The field is moving from questions posed to a detached observer toward agents that must localize themselves, gather evidence, and act, yet no current benchmark exercises perception fidelity, symbolic precision, memory persistence, and long-horizon action together, and few attribute an end-to-end score to the components that produced it (§\ref{sec:p3_s7_5}).

Reasoning in this chapter takes the evidence of Chapter~\ref{chapter:3:sec:6} as given. Chapter~\ref{sec:generation} turns to the inverse problem of producing spatial content, and its world models (§\ref{sec:p3_s8_6}) supply the simulators that world-model-based reasoning (§\ref{subsec:p3_s7_3_8}) consults.
